% Lesson 6 — Vocabulary: Words and expressions related to taking part in communal work — Anglais 1ère A — Cours signé OnBuch+
% Programme officiel MINESEC. Compiler avec : tectonic ENA06-vocabulary-words-and-expressions-related-to-taking-part.tex
\newcommand{\DOCMATIERE}{Anglais}
\newcommand{\DOCNIVEAU}{1ère A}
\newcommand{\DOCMODULE}{Chapter 1 : Family and Social Life}
\newcommand{\DOCLECON}{ENA06}
\newcommand{\DOCTITRE}{Lesson 6 — Vocabulary: Words and expressions related to taking part in communal work}
% OnBuch+ — préambule « cours signés » ENRICHI (compilé avec tectonic / XeLaTeX)
% Système visuel « Pop » d'OnBuch+ : fond crème (jamais blanc), bordures encre
% épaisses, ombres « dures » décalées, titres Archivo Black en capitales.
% Version MATHÉMATIQUES : graphes (pgfplots/tikz), boîtes pédagogiques
% (définition, propriété, méthode, attention, le savais-tu, exercice résolu,
% exercice, TP, bilan), page de garde et signature OnBuch+ ancrée en bas.
%
% Macros attendues dans main.tex AVANT \input{preamble} :
%   \DOCMATIERE  \DOCNIVEAU  \DOCMODULE  \DOCLECON (numéro)  \DOCTITRE
\documentclass[11pt]{article}

\usepackage{fontspec}
\usepackage[english,french]{babel} % français = langue principale ; l'anglais via \eng{..} / textebox
\usepackage[a4paper,margin=2cm,top=3.4cm,bottom=2.8cm,headheight=1.4cm,footskip=1.3cm]{geometry}
\usepackage{amsmath,amssymb,amsfonts}
\usepackage{mathtools} % \xmapsto, \coloneqq... souvent utilisés par les modèles
\usepackage{cancel} % simplifications barrées (bilans de forces)
% \abs{z} (module, valeur absolue) et \norm{u} (norme), utilisés par les modèles
\providecommand{\abs}{}\let\abs\relax\DeclarePairedDelimiter{\abs}{\lvert}{\rvert}
\providecommand{\norm}{}\let\norm\relax\DeclarePairedDelimiter{\norm}{\lVert}{\rVert}
\usepackage[table]{xcolor} % [table] : fournit aussi \rowcolors (alternance de lignes)
\usepackage{pagecolor}
\usepackage{tikz}
\usetikzlibrary{arrows.meta,positioning,calc,shapes.geometric,shapes.misc,shapes.arrows,decorations.pathmorphing,decorations.pathreplacing,decorations.markings,patterns,shadows,mindmap,trees,fit,backgrounds,matrix,intersections,angles,quotes}
\usepackage{pgfplots}
\usepackage[version=4]{mhchem} % équations nucléaires \ce{^{235}_{92}U}
\usepackage[european,siunitx]{circuitikz} % schémas électriques
\pgfplotsset{compat=1.18}
\usepgfplotslibrary{fillbetween,groupplots} % aires entre courbes (calcul intégral)
\usepackage{enumitem}
\usepackage{fancyhdr}
% [most] charge toutes les bibliothèques tcolorbox (listings, minted, raster,
% documentation...) et gonfle inutilement la mémoire TeX sur un document long
% (~300 boîtes) : on ne charge que ce qui est réellement utilisé.
\usepackage[skins,breakable]{tcolorbox}
\usepackage{graphicx}
\usepackage{colortbl}
\usepackage{booktabs}
\usepackage{tabularx}
\usepackage{diagbox} % case d'en-tête diagonale des tableaux à double entrée
% Colonnes extensibles centrée / gauche / droite (C, L, R), souvent utilisées par les modèles
\newcolumntype{C}{>{\centering\arraybackslash}X}
\newcolumntype{L}{>{\raggedright\arraybackslash}X}
\newcolumntype{R}{>{\raggedleft\arraybackslash}X}
\usepackage{array}
\usepackage{multicol}
\usepackage{multirow}
\usepackage{siunitx}
% parse-numbers=false : accepte aussi \SI{2,5 \times 10^{-3}}{...}
\sisetup{locale=FR,output-decimal-marker={,},group-minimum-digits=4,parse-numbers=false}
\DeclareSIUnit\ohm{\ensuremath{\symup{\Omega}}} % U+2126 absent de la police
\DeclareSIUnit\division{div} % réglages d'oscilloscope (ms/div, V/div)
\DeclareSIUnit\tr{tr} % tours (vitesse de rotation, ex. \SI{1500}{\tr\per\minute})
\DeclareSIUnit\molar{M} % concentration molaire (ex. \SI{3}{\molar}, \SI{1}{\micro\molar})
\DeclareSIUnit\an{an} % année (le modèle confond parfois avec \year, un compteur TeX) — ex. \SI{5}{\kilo\meter\per\an}
\DeclareSIUnit\FCFA{FCFA} % franc CFA (ex. \SI{500}{\FCFA\per\kilo\gram})
\DeclareSIUnit\degreeBrix{\textdegree Brix} % teneur en sucre d'un sirop/jus (réfractométrie)
\DeclareSIUnit\month{mois} % le modèle utilise parfois \month, un compteur TeX, comme unité de durée
\usepackage{qrcode}
% \degree : commande fréquemment utilisée par les modèles pour un angle ou une
% température en dehors de \SI{}{} (non fournie par les paquets chargés).
\providecommand{\degree}{^{\circ}}
% \qed (fin de démonstration), utilisé par les modèles sans amsthm
\providecommand{\qed}{\hfill$\square$}
% \barre : double barre des valeurs interdites dans un tableau de variations (tkz-tab)
\providecommand{\barre}{\ensuremath{\|}}
% environnement proof (amsthm non chargé) utilisé par les modèles
\ifdefined\proof\else\newenvironment{proof}[1][Démonstration]{\par\noindent\textit{#1.}\ }{\qed\par}\fi
\usepackage{lastpage}
\usepackage{hyperref}
\hypersetup{hidelinks}

% ============================================================
% Palette « Pop » OnBuch+ — mêmes hex que la classe P (plus_ui.dart)
% ============================================================
\definecolor{poporange}{HTML}{F59321}
\definecolor{popdark}{HTML}{E07A0C}
\definecolor{popcream}{HTML}{FFF6E8}
\definecolor{popink}{HTML}{1C1714}
\definecolor{poppurple}{HTML}{7A5AE0}
\definecolor{popgreen}{HTML}{2E9E6A}
\definecolor{poppink}{HTML}{E0457B}
\definecolor{popblue}{HTML}{2F7DE1}
\definecolor{popgold}{HTML}{C98A12}
\definecolor{popmuted}{HTML}{8A7F76}
\definecolor{popline}{HTML}{EADFD2}
\definecolor{popgray}{HTML}{8A7F76}
\colorlet{popgrey}{popgray}
% Alias tolérants (le modèle nomme parfois les couleurs différemment)
\colorlet{popwhite}{white}
\colorlet{popblack}{popink}
\colorlet{popred}{poppink}
\colorlet{popyellow}{popgold}
% Teintes claires (fonds de tableaux/encadrés) — le modèle nomme parfois
% les couleurs avec un suffixe L (ex. popblueL) sans les avoir définies.
\colorlet{poporangeL}{poporange!20}
\colorlet{popdarkL}{popdark!20}
\colorlet{poppurpleL}{poppurple!20}
\colorlet{popgreenL}{popgreen!20}
\colorlet{poppinkL}{poppink!20}
\colorlet{popblueL}{popblue!20}
\colorlet{popgoldL}{popgold!20}
\colorlet{popredL}{popred!20}
\colorlet{popgrayL}{popgray!20}
% Teintes claires pour fonds de tableaux / zones
\colorlet{popblueL}{popblue!10!white}
\colorlet{poporangeL}{poporange!14!white}
\colorlet{popgreenL}{popgreen!12!white}
\colorlet{poppurpleL}{poppurple!12!white}
\colorlet{poppinkL}{poppink!12!white}
\colorlet{popgoldL}{popgold!14!white}

\pagecolor{popcream}

% ============================================================
% Typographie
% ============================================================
\defaultfontfeatures{Ligatures=TeX}
\setmainfont{PlusJakartaSans-Medium.ttf}[Path=../../../fonts/,BoldFont=PlusJakartaSans-Bold.ttf]
% Symboles, API et emojis absents de la police principale : repli sur DejaVu Sans (caractères actifs)
\newfontface\symfont{DejaVuSans.ttf}[Path=../../../fonts/]
\makeatletter
\newcommand\symactive[1]{\catcode#1=13 \begingroup\lccode`\~=#1 \lowercase{\endgroup\def~}{{\symfont\char#1}}}
\newcommand\symblank[1]{\catcode#1=13 \begingroup\lccode`\~=#1 \lowercase{\endgroup\def~}{}}
\newcommand\symhyph[1]{\catcode#1=13 \begingroup\lccode`\~=#1 \lowercase{\endgroup\def~}{\mbox{-}}}
\newcount\symc
\newcommand\symrange[2]{\symc=#1 \@whilenum\symc<#2 \do{\expandafter\symactive\expandafter{\the\symc}\advance\symc by 1 }}
\newcommand\symblankrange[2]{\symc=#1 \@whilenum\symc<#2 \do{\expandafter\symblank\expandafter{\the\symc}\advance\symc by 1 }}
\makeatother
\symrange{"0250}{"02FF}\symactive{"2713}\symactive{"2714}\symactive{"2717}\symactive{"2718}\symactive{"2610}\symactive{"2611}\symactive{"2612}
\symhyph{"2011}\symhyph{"2E1A}\symblankrange{"1F300}{"1FAFF}\symblank{"FE0F}
% Maths en sans-serif (Fira Math) avec chiffres et lettres droites de Plus
% Jakarta, pour que les formules se fondent dans le texte.
\usepackage[math-style=ISO,bold-style=ISO]{unicode-math}
\setmathfont{FiraMath-Regular.otf}[Path=../../../fonts/]
\setmathfont{PlusJakartaSans-Medium.ttf}[Path=../../../fonts/,range={up/num,up/{Latin,latin},\mathup}]
\usepackage{icomma} % virgule décimale sans espace en mode math
% Caractères mathématiques Unicode absents des polices de texte (Plus Jakarta,
% Archivo Black) : rendus par leur équivalent mathématique, en texte comme en
% formule — sinon ils s'affichent en carré vide (ex. « Arithmétique dans ℤ »).
\usepackage{newunicodechar}
\newunicodechar{ℝ}{\ensuremath{\mathbb{R}}}
\newunicodechar{ℤ}{\ensuremath{\mathbb{Z}}}
\newunicodechar{ℂ}{\ensuremath{\mathbb{C}}}
\newunicodechar{ℕ}{\ensuremath{\mathbb{N}}}
\newunicodechar{ℚ}{\ensuremath{\mathbb{Q}}}
\newunicodechar{𝔻}{\ensuremath{\mathbb{D}}}
\newunicodechar{α}{\ensuremath{\alpha}}
\newunicodechar{β}{\ensuremath{\beta}}
\newunicodechar{γ}{\ensuremath{\gamma}}
\newunicodechar{δ}{\ensuremath{\delta}}
\newunicodechar{ε}{\ensuremath{\varepsilon}}
\newunicodechar{θ}{\ensuremath{\theta}}
\newunicodechar{λ}{\ensuremath{\lambda}}
\newunicodechar{μ}{\ensuremath{\mu}}
\newunicodechar{σ}{\ensuremath{\sigma}}
\newunicodechar{τ}{\ensuremath{\tau}}
\newunicodechar{φ}{\ensuremath{\varphi}}
\newunicodechar{ψ}{\ensuremath{\psi}}
\newunicodechar{ω}{\ensuremath{\omega}}
\newunicodechar{Σ}{\ensuremath{\Sigma}}
% majuscules produites par \MakeUppercase dans les titres (α-aminés -> Α)
\newunicodechar{Α}{\ensuremath{\alpha}}
\newunicodechar{Β}{\ensuremath{\beta}}
\newunicodechar{Γ}{\ensuremath{\Gamma}}
\newunicodechar{Δ}{\ensuremath{\Delta}}
\newunicodechar{Ε}{\ensuremath{\varepsilon}}
\newunicodechar{Θ}{\ensuremath{\Theta}}
\newunicodechar{Λ}{\ensuremath{\Lambda}}
\newunicodechar{Μ}{\ensuremath{\mu}}
\newunicodechar{Τ}{\ensuremath{\tau}}
\newunicodechar{Φ}{\ensuremath{\Phi}}
\newunicodechar{Ψ}{\ensuremath{\Psi}}
\newunicodechar{Ω}{\ensuremath{\Omega}}
\newunicodechar{↦}{\ensuremath{\mapsto}}
\newunicodechar{∈}{\ensuremath{\in}}
\newunicodechar{∉}{\ensuremath{\notin}}
\newunicodechar{∧}{\ensuremath{\wedge}}
\newunicodechar{⇔}{\ensuremath{\Leftrightarrow}}
\newunicodechar{⟺}{\ensuremath{\Longleftrightarrow}}
\newunicodechar{⇒}{\ensuremath{\Rightarrow}}
\newunicodechar{⟹}{\ensuremath{\Longrightarrow}}
\newunicodechar{∘}{\ensuremath{\circ}}
\newunicodechar{∪}{\ensuremath{\cup}}
\newunicodechar{∩}{\ensuremath{\cap}}
\newunicodechar{≡}{\ensuremath{\equiv}}
\newunicodechar{⊂}{\ensuremath{\subset}}
\newunicodechar{⊥}{\ensuremath{\perp}}
\newunicodechar{∃}{\ensuremath{\exists}}
\newunicodechar{∀}{\ensuremath{\forall}}
\newunicodechar{∛}{\ensuremath{\sqrt[3]{\,}}}
\newunicodechar{∜}{\ensuremath{\sqrt[4]{\,}}}
\newunicodechar{⃗}{\ensuremath{{}^{\to}}}
\newunicodechar{ⁿ}{\ifmmode^{n}\else\textsuperscript{\ensuremath{n}}\fi}
\newunicodechar{ˣ}{\ifmmode^{x}\else\textsuperscript{\ensuremath{x}}\fi}
\newunicodechar{ᵇ}{\ifmmode^{b}\else\textsuperscript{\ensuremath{b}}\fi}
\newunicodechar{ʳ}{\ifmmode^{r}\else\textsuperscript{\ensuremath{r}}\fi}
\newunicodechar{ᵘ}{\ifmmode^{u}\else\textsuperscript{\ensuremath{u}}\fi}
\newunicodechar{ʸ}{\ifmmode^{y}\else\textsuperscript{\ensuremath{y}}\fi}
\newunicodechar{ᵃ}{\ifmmode^{a}\else\textsuperscript{\ensuremath{a}}\fi}
\newunicodechar{ᵉ}{\ifmmode^{e}\else\textsuperscript{\ensuremath{e}}\fi}
\newunicodechar{ᵏ}{\ifmmode^{k}\else\textsuperscript{\ensuremath{k}}\fi}
\newunicodechar{ᵖ}{\ifmmode^{p}\else\textsuperscript{\ensuremath{p}}\fi}
\newunicodechar{ᴺ}{\ifmmode^{N}\else\textsuperscript{\ensuremath{N}}\fi}
\newunicodechar{ᶜ}{\ifmmode^{c}\else\textsuperscript{\ensuremath{c}}\fi}
\newunicodechar{ʰ}{\ifmmode^{h}\else\textsuperscript{\ensuremath{h}}\fi}
\newunicodechar{ᵠ}{\ifmmode^{q}\else\textsuperscript{\ensuremath{q}}\fi}
\newunicodechar{ₙ}{\ifmmode_{n}\else\textsubscript{\ensuremath{n}}\fi}
\newunicodechar{ₐ}{\ifmmode_{a}\else\textsubscript{\ensuremath{a}}\fi}
\newunicodechar{ᵢ}{\ifmmode_{i}\else\textsubscript{\ensuremath{i}}\fi}
\newunicodechar{ᵣ}{\ifmmode_{r}\else\textsubscript{\ensuremath{r}}\fi}
\newunicodechar{ₖ}{\ifmmode_{k}\else\textsubscript{\ensuremath{k}}\fi}
\newunicodechar{ᵦ}{\ifmmode_{\beta}\else\textsubscript{\ensuremath{\beta}}\fi}
\newunicodechar{ₓ}{\ifmmode_{x}\else\textsubscript{\ensuremath{x}}\fi}
\newfontfamily\popdisplay{ArchivoBlack-Regular.ttf}[Path=../../../fonts/]
\newfontfamily\poplogo{SpaceGrotesk-Bold.ttf}[Path=../../../fonts/]
\newfontfamily\popsemi{PlusJakartaSans-SemiBold.ttf}[Path=../../../fonts/]
\newfontfamily\popxbold{PlusJakartaSans-ExtraBold.ttf}[Path=../../../fonts/]
\newfontfamily\popxboldls{PlusJakartaSans-ExtraBold.ttf}[Path=../../../fonts/,LetterSpace=3.0]

\setlist[enumerate]{leftmargin=1.7em,itemsep=5pt,topsep=4pt}
\setlist[itemize]{leftmargin=1.7em,itemsep=4pt,topsep=4pt,label={\color{poporange}\textbullet}}
\setlength{\parindent}{0pt}
\setlength{\parskip}{5pt}
\renewcommand{\arraystretch}{1.25}

% pgfplots : style Pop par défaut
\pgfplotsset{
  every axis/.append style={axis line style={popink,line width=1pt},tick style={popink},
    grid style={popline,line width=0.5pt},label style={font=\small},tick label style={font=\footnotesize}},
  popcurve/.style={poporange,line width=1.8pt,no markers,smooth},
}
% Même style disponible dans un \draw TikZ simple (hors pgfplots)
\tikzset{popcurve/.style={poporange,line width=1.8pt}}
\tikzset{popgrid/.style={popline,very thin}}
\colorlet{popcurve}{poporange} % le nom du style est parfois utilisé comme couleur

% ============================================================
% Pill / squiggle
% ============================================================
\newcommand{\poppill}[3]{%
  \tikz[baseline=-0.55ex]\node[draw=popink,line width=1.2pt,rounded corners=2.6pt,%
    fill=#2,inner xsep=6.5pt,inner ysep=3pt]{%
    \popxboldls\fontsize{7}{7}\selectfont\color{#3}\MakeUppercase{#1}};%
}
\newcommand{\popsquiggle}[1]{%
  \tikz[baseline=-0.4ex]{
    \draw[#1,line width=2.6pt,line cap=round]
      plot [smooth] coordinates {(0,0.09) (0.34,0.01) (0.68,0.21) (1.02,0.01) (1.36,0.10)};
  }%
}
% Mot-clé surligné (marqueur orange sous le texte)
\newcommand{\cle}[1]{{\popxbold\color{popdark}#1}}

% ============================================================
% En-tête / pied
% ============================================================
\pagestyle{fancy}
\fancyhf{}
\renewcommand{\headrulewidth}{0pt}
\renewcommand{\footrulewidth}{0pt}
\lhead{\raisebox{-0.15ex}{\poplogo\fontsize{13}{13}\selectfont\color{popink}OnBuch\color{poporange}+}}
\rhead{\poppill{\DOCMATIERE\ \textbullet\ \DOCNIVEAU\ \textbullet\ Leçon \DOCLECON}{white}{popink}}
\lfoot{\popxbold\fontsize{7.6}{7.6}\selectfont\color{popmuted}\MakeUppercase{Cours signé OnBuch+}\ \textbullet\ {\popsemi\DOCTITRE}}
\rfoot{\tikz[baseline=-0.6ex]\node[rounded corners=8.5pt,fill=popink,minimum height=17pt,minimum width=17pt,inner xsep=5pt,inner sep=0]{\popxbold\fontsize{8}{8}\selectfont\color{popcream}\thepage\,/\,\pageref*{LastPage}};}

% ============================================================
% Titres
% ============================================================
\newcommand{\chaptitre}[2]{%
  \begin{tcolorbox}[enhanced,colback=poporange,colframe=popink,boxrule=2.6pt,arc=11pt,
      drop shadow=popink,left=15pt,right=15pt,top=12pt,bottom=11pt,before skip=3pt,after skip=18pt]
    \poppill{#2}{popink}{popcream}\par\vspace{9pt}
    {\raggedright\hyphenpenalty=10000\exhyphenpenalty=10000\popdisplay\fontsize{22}{24}\selectfont\color{popcream}\MakeUppercase{#1}\par}
  \end{tcolorbox}
}
% Section numérotée : pastille orange avec le numéro + titre Archivo
\newcounter{popsec}
\newcommand{\coursec}[1]{%
  \stepcounter{popsec}\setcounter{popsub}{0}%
  \par\vspace{16pt}\noindent
  \begin{minipage}{\linewidth}
  \tikz[baseline=-0.7ex]\node[draw=popink,line width=1.6pt,fill=poporange,rounded corners=4pt,minimum size=22pt,inner sep=0]{\popdisplay\fontsize{12}{12}\selectfont\color{popcream}\thepopsec};%
  \hspace{8pt}{\popdisplay\fontsize{15}{17}\selectfont\color{popink}\MakeUppercase{#1}}\par
  \vspace{4pt}\noindent{\color{poporange}\rule{36pt}{3.6pt}}
  \end{minipage}\par\nopagebreak\vspace{8pt}
}
\newcounter{popsub}
\newcommand{\courssub}[1]{%
  \stepcounter{popsub}%
  \par\vspace{9pt}\noindent{\popxbold\fontsize{11.5}{13}\selectfont\color{popink}{\color{poporange}\thepopsec.\thepopsub}\hspace{6pt}#1}\par\nopagebreak\vspace{4pt}
}

% ============================================================
% Boîtes pédagogiques — même recette : fond blanc, bordure épaisse colorée,
% ombre dure encre, badge plein. Titre optionnel [#1].
% ============================================================
\tcbset{popbase/.style={breakable,enhanced,colback=white,boxrule=2.3pt,arc=9pt,
  shadow={3pt}{-3pt}{0pt}{popink},left=14pt,right=14pt,top=11pt,bottom=11pt,before skip=11pt,after skip=15pt,
  fontupper=\popsemi\fontsize{10.6}{14.5}\selectfont\color{popink},
  fontlower=\popsemi\fontsize{10.6}{14.5}\selectfont\color{popink}}}
% Coche dessinée (le glyphe ✓ n'existe pas dans les polices du document)
\newcommand{\popcheck}[1]{\tikz[baseline=-0.5ex]\draw[#1,line width=1.6pt,line cap=round,line join=round](-0.13,0.0)--(-0.04,-0.09)--(0.13,0.1);}
\providecommand{\coche}{\popcheck{popgreen}}
\providecommand{\resultat}[1]{\par\smallskip\noindent\fcolorbox{popgreen}{popgreenL}{\parbox{\dimexpr\linewidth-2\fboxsep-2\fboxrule\relax}{\textbf{Résultat.} #1}}\par\smallskip}
\newcommand{\popboxhead}[3]{\poppill{#1}{#2}{white}\ifx\relax#3\relax\else\hspace{6pt}{\popxbold\fontsize{10.6}{12}\selectfont\color{popink}#3}\fi\par\vspace{7pt}}

\newtcolorbox{aretenir}[1][]{popbase,colframe=popgreen,before upper={\popboxhead{À retenir}{popgreen}{#1}}}
% Texte en anglais (document de lecture, dialogue, extrait) : cadre
% d'étude. Un titre optionnel (source, auteur) s'affiche dans l'en-tête.
\newtcolorbox{textebox}[1][]{popbase,colframe=popblue,before upper={\popboxhead{Text}{popblue}{#1}\begin{otherlanguage}{english}},after upper={\end{otherlanguage}}}
% Marques ✓ / ✗ (forme correcte / fautive) souvent utilisées par les modèles
\providecommand{\cross}{\textcolor{poppink}{\ensuremath{\times}}}
\providecommand{\xmark}{\cross}\providecommand{\crossmark}{\cross}\providecommand{\cmark}{\ensuremath{\checkmark}}
% Ligne de réponse / justification à compléter (inventée par les modèles)
\providecommand{\justification}[1]{\par\noindent\textit{Justification~:} \dotfill\par}
% \textipa (package tipa non chargé) : la police ne couvre pas l'API -> simple passage
\providecommand{\textipa}[1]{#1}
% Soulignement (ulem non chargé) : \uline, \ul
\providecommand{\uline}[1]{\underline{#1}}\providecommand{\ul}[1]{\underline{#1}}
% \glossaire{\item ...} inventée par les modèles : liste de glossaire
\providecommand{\glossaire}[1]{\begin{itemize}#1\end{itemize}}
% \alert (beamer) inventée par les modèles : texte mis en évidence
\providecommand{\alert}[1]{\textbf{\textcolor{poppink}{#1}}}
% variantes ASCII d'ordinaux écrites en macro
\providecommand{\iere}{\textsuperscript{re}}\providecommand{\ieme}{\textsuperscript{e}}\providecommand{\ere}{\textsuperscript{re}}\providecommand{\eme}{\textsuperscript{e}}\providecommand{\er}{\textsuperscript{er}}\providecommand{\ie}{\textsuperscript{e}}
% Ordinaux français écrits en macro par les modèles : 1\ière, 3\ième
\newcommand{\ière}{\textsuperscript{re}}\newcommand{\ième}{\textsuperscript{e}}
% \exemple (inventée par les modèles) : étiquette « Exemple : »
\providecommand{\exemple}{\textbf{Exemple~:}\ }
% \square utilisable aussi hors mode math (cases à cocher)
\AtBeginDocument{\let\squareMath\square\renewcommand{\square}{\ensuremath{\squareMath}}}
% Mot ou phrase en anglais à l'intérieur d'un paragraphe français
\newcommand{\eng}[1]{\ifmmode\text{\textit{#1}}\else\begingroup\selectlanguage{english}\itshape #1\endgroup\fi}
\let\esp\eng % alias toléré
% flèches d'intonation inventées par les modèles
\providecommand{\textsearrow}{}\renewcommand{\textsearrow}{\ensuremath{\searrow}}\providecommand{\textnearrow}{}\renewcommand{\textnearrow}{\ensuremath{\nearrow}}
% \fr{..} : mot français (inventé par les modèles)
\providecommand{\fr}[1]{\textit{#1}}
\newtcolorbox{exemplebox}[1][]{popbase,colframe=poporange,before upper={\popboxhead{Exemple}{poporange}{#1}}}
\newtcolorbox{definition}[1][]{popbase,colframe=popblue,before upper={\popboxhead{Définition}{popblue}{#1}}}
\newtcolorbox{propriete}[1][]{popbase,colframe=poppurple,before upper={\popboxhead{Propriété}{poppurple}{#1}}}
\newtcolorbox{methode}[1][]{popbase,colframe=popgold,before upper={\popboxhead{Méthode}{popgold}{#1}}}
\newtcolorbox{attention}[1][]{popbase,colframe=poppink,before upper={\popboxhead{Attention}{poppink}{#1}}}
\newtcolorbox{experience}[1][]{popbase,colframe=popblue,colback=popblueL,before upper={\popboxhead{Expérience}{popblue}{#1}}}
% « Le savais-tu ? » : carte violette pleine (contexte, culture, Cameroun)
\newtcolorbox{savaistu}[1][]{popbase,colback=poppurple,colframe=popink,
  fontupper=\popsemi\fontsize{10.4}{14.2}\selectfont\color{white},
  before upper={\poppill{Le savais-tu\,?}{white}{poppurple}\ifx\relax#1\relax\else\hspace{6pt}{\popxbold\color{white}#1}\fi\par\vspace{7pt}}}
% Exercice résolu : énoncé en haut, \tcblower puis la solution sur fond crème
\newcounter{popexo}\newcounter{popexr}
\newtcolorbox{exoresolu}[1][]{popbase,colframe=poporange,colbacklower=popcream,
  segmentation style={popink,line width=1.2pt,dashed},
  before upper={\stepcounter{popexr}\popboxhead{Exercice résolu \thepopexr}{poporange}{#1}},
  before lower={\poppill{Solution}{popink}{popcream}\par\vspace{6pt}}}
% Exercice (non résolu en place) — la correction est dans la partie Corrigés
\newtcolorbox{exercice}[1][]{popbase,colframe=popink,
  before upper={\stepcounter{popexo}\popboxhead{Exercice \thepopexo}{popink}{#1}}}
% Corrigé
\newtcolorbox{corrige}[1][]{popbase,colframe=popgreen,colback=popgreenL,
  before upper={\popboxhead{Corrigé}{popgreen}{#1}}}
% Objectifs / prérequis / bilan
\newtcolorbox{objectifs}{popbase,colframe=popink,colback=poporangeL,
  before upper={\popboxhead{Objectifs de la leçon}{popink}{}}}
\newtcolorbox{prerequis}{popbase,colframe=popink,colback=popblueL,
  before upper={\popboxhead{Prérequis}{popblue}{}}}
\newtcolorbox{bilan}{popbase,colframe=popink,colback=white,boxrule=2.8pt,
  before upper={\popboxhead{Fiche bilan}{poporange}{L'essentiel à connaître pour l'examen}}}
% Figure encadrée avec légende
\newtcolorbox{popfigure}[1][]{enhanced,breakable=false,colback=white,colframe=popline,boxrule=1.4pt,arc=8pt,
  left=10pt,right=10pt,top=10pt,bottom=8pt,before skip=10pt,after skip=12pt,halign=center,
  lower separated=false,halign lower=center,
  fontlower=\popsemi\fontsize{9}{11.5}\selectfont\color{popmuted},
  #1}
\newcommand{\legende}[1]{\par\vspace{6pt}{\popsemi\fontsize{9}{11.5}\selectfont\color{popmuted}#1}}

% ============================================================
% Signature OnBuch+
% ============================================================
\newcommand{\poplogoinline}[1]{{\poplogo\fontsize{#1}{#1}\selectfont\color{popink}OnBuch\color{poporange}+}}
\newcommand{\signatureonbuch}{%
  \begin{tcolorbox}[enhanced,colback=popink,colframe=popink,boxrule=2.2pt,arc=10pt,
      drop shadow=poporange,left=14pt,right=14pt,top=10pt,bottom=10pt,before skip=0pt,after skip=0pt]
    \tikz[baseline=-0.7ex]\node[circle,fill=popgreen,draw=popcream,line width=1.3pt,minimum size=22pt,inner sep=0]{\popcheck{white}};%
    \hspace{9pt}\begin{minipage}[c]{0.62\linewidth}
      {\popxbold\fontsize{12}{13}\selectfont\color{popcream}Signé\ \textbullet\ {\poplogo OnBuch\color{poporange}+}}\par\vspace{2pt}
      {\raggedright\popsemi\fontsize{8}{10.5}\selectfont\color{popline}\MakeUppercase{Équipe pédagogique OnBuch+}\\\MakeUppercase{Conforme au programme officiel MINESEC}\par}
    \end{minipage}\hfill
    {\poplogo\fontsize{22}{22}\selectfont\color{popcream}OnBuch\color{poporange}+}
  \end{tcolorbox}%
}

% ============================================================
% Page de garde
% #1 titre · #2 matière · #3 niveau · #4 module · #5 n° leçon · #6 durée ·
% #7 description / objectifs (texte ou itemize)
% ============================================================
\newcommand{\pagegarde}[7]{%
  \thispagestyle{empty}
  \newgeometry{a4paper,margin=1.7cm,top=1.6cm,bottom=1.4cm}%
  \begin{tikzpicture}[remember picture,overlay]
    \fill[poporange!18] ([xshift=-3.2cm,yshift=-3.6cm]current page.north east) circle (4.6cm);
    \fill[poppurple!14] ([xshift=2.4cm,yshift=7.6cm]current page.south west) circle (3.8cm);
  \end{tikzpicture}%
  {\poplogo\fontsize{30}{30}\selectfont\color{popink}OnBuch\color{poporange}+}\hfill
  \raisebox{0.6ex}{\poppill{Cours signé \textbullet\ Premium}{popink}{popcream}}\par
  \vspace{1pt}\popsquiggle{poppurple}\par
  \vspace{8pt}
  \begin{tcolorbox}[enhanced,colback=poporange,colframe=popink,boxrule=2.8pt,arc=13pt,
      drop shadow=popink,left=18pt,right=18pt,top=12pt,bottom=13pt,after skip=10pt]
    \poppill{#2 \textbullet\ #3}{popink}{popcream}\hspace{5pt}\poppill{#4}{white}{popink}\par\vspace{8pt}
    {\popdisplay\fontsize{14}{14}\selectfont\color{popink}LEÇON #5}\par\vspace{4pt}
    {\raggedright\hyphenpenalty=10000\exhyphenpenalty=10000\popdisplay\fontsize{25}{27}\selectfont\color{popcream}\MakeUppercase{#1}\par}
    \vspace{8pt}
    {\popxbold\fontsize{9.5}{9.5}\selectfont\color{popink}\MakeUppercase{Durée conseillée : #6}}
  \end{tcolorbox}
  \begin{tcolorbox}[enhanced,colback=white,colframe=popink,boxrule=2.2pt,arc=10pt,
      drop shadow=popink,left=16pt,right=16pt,top=10pt,bottom=10pt,after skip=8pt]
    \poppill{Dans cette leçon}{poppurple}{white}\par\vspace{5pt}
    \popsemi\fontsize{9.3}{12.6}\selectfont\color{popink}#7
  \end{tcolorbox}
  \noindent\begin{minipage}[t]{0.487\linewidth}
    \begin{tcolorbox}[enhanced,colback=popgreen,colframe=popink,boxrule=2.2pt,arc=10pt,
        drop shadow=popink,left=11pt,right=11pt,top=10pt,bottom=10pt,height=3.1cm,valign=center]
      \tcbox[colback=white,colframe=popink,boxrule=1.3pt,arc=5pt,boxsep=0pt,nobeforeafter,
        left=3pt,right=3pt,top=3pt,bottom=3pt,valign=center]{\qrcode[height=1.9cm]{https://play.google.com/store/apps/details?id=com.luvvix.onbuch&hl=fr}}%
      \hspace{8pt}\begin{minipage}[c]{0.5\linewidth}
        {\popdisplay\fontsize{11}{12.5}\selectfont\color{white}\MakeUppercase{Télécharge OnBuch}}\par\vspace{4pt}
        {\raggedright\popsemi\fontsize{8.4}{11}\selectfont\color{white}Gratuit \textbullet\ Google Play\\Tuteur IA, annales, concours, résultats\par}
      \end{minipage}
    \end{tcolorbox}
  \end{minipage}\hfill
  \begin{minipage}[t]{0.487\linewidth}
    \begin{tcolorbox}[enhanced,colback=poppurple,colframe=popink,boxrule=2.2pt,arc=10pt,
        drop shadow=popink,left=11pt,right=11pt,top=10pt,bottom=10pt,height=3.1cm,valign=center]
      \tikz[baseline=-0.7ex]\node[circle,fill=white,draw=popink,line width=1.3pt,minimum size=1.6cm,inner sep=0]{\popdisplay\fontsize{24}{24}\selectfont\color{poppurple}+};%
      \hspace{8pt}\begin{minipage}[c]{0.6\linewidth}
        {\popdisplay\fontsize{11}{12.5}\selectfont\color{white}\MakeUppercase{Passe à OnBuch+}}\par\vspace{4pt}
        {\raggedright\popsemi\fontsize{8.4}{11}\selectfont\color{white}Tous les cours signés, Tuteur IA illimité, fiches PDF — onglet OnBuch+ de l'app\par}
      \end{minipage}
    \end{tcolorbox}
  \end{minipage}
  \vspace{8pt}
  \popcontenu
  \vfill
  \signatureonbuch
  \restoregeometry
  \newpage
}

% Rangée « Ce cours contient » (page de garde)
\newcommand{\poptile}[3]{% #1 couleur #2 gros texte #3 légende
  \begin{tcolorbox}[enhanced,colback=#1,colframe=popink,boxrule=2pt,arc=9pt,shadow={2.5pt}{-2.5pt}{0pt}{popink},
      left=6pt,right=6pt,top=9pt,bottom=9pt,halign=center,valign=center,height=1.75cm,width=\linewidth,nobeforeafter]
    {\popdisplay\fontsize{12.5}{13}\selectfont\color{white}\MakeUppercase{#2}}\par\vspace{3pt}
    {\centering\popsemi\fontsize{7.8}{9.5}\selectfont\color{white}#3\par}
  \end{tcolorbox}}
\newcommand{\popcontenu}{%
  \par\noindent{\popxboldls\fontsize{7.5}{7.5}\selectfont\color{popmuted}CE COURS SIGNÉ CONTIENT}\par\vspace{7pt}
  \noindent\begin{minipage}[t]{0.235\linewidth}\poptile{popblue}{Cours}{complet, illustré, pas à pas}\end{minipage}\hfill
  \begin{minipage}[t]{0.235\linewidth}\poptile{poporange}{Méthodes}{et exercices résolus}\end{minipage}\hfill
  \begin{minipage}[t]{0.235\linewidth}\poptile{poppink}{Exercices}{type examen + corrigés}\end{minipage}\hfill
  \begin{minipage}[t]{0.235\linewidth}\poptile{popgreen}{Fiche bilan}{l'essentiel à retenir}\end{minipage}\par}

% Sommaire
\newtcolorbox{sommaire}{enhanced,colback=white,colframe=popink,boxrule=2.2pt,arc=10pt,
  drop shadow=popink,left=18pt,right=18pt,top=13pt,bottom=13pt,before skip=6pt,after skip=20pt,
  before upper={\poppill{Sommaire}{poppurple}{white}\par\vspace{9pt}\popsemi\fontsize{10.6}{14.5}\selectfont\color{popink}}}

% Signature de fin de cours — poussée tout en bas de la dernière page
\newcommand{\findecours}{%
  \par\vfill\nopagebreak
  \begin{center}\popsquiggle{poporange}\end{center}\vspace{4pt}
  \signatureonbuch
}
\AtBeginDocument{\def\llbracket{\mathopen{[\mkern-3.5mu[}}\def\rrbracket{\mathclose{]\mkern-3.5mu]}}} % intervalles d'entiers (glyphe ⟦ absent de la police)
% Glyphes absents de Fira Math : redéfinis à partir de glyphes disponibles
\AtBeginDocument{%
  \def\setminus{\mathbin{\backslash}}\let\smallsetminus\setminus
  \def\vdots{\mathord{\vbox{\baselineskip3.2pt\lineskiplimit0pt\kern4pt\hbox{.}\hbox{.}\hbox{.}}}}%
  \def\ddots{\mathinner{\mkern1mu\raise6.5pt\hbox{.}\mkern2mu\raise3.6pt\hbox{.}\mkern2mu\raise0.7pt\hbox{.}\mkern1mu}}%
  \def\triangle{\mathord{\scalebox{0.9}{$\symup{\Delta}$}}}%
  \def\nabla{\mathord{\raisebox{\depth}{\scalebox{1}[-1]{$\symup{\Delta}$}}}}%
  \def\checkmark{\popcheck{popdark}}%
  \def\star{\mathbin{\ast}}\def\bigstar{\mathord{\ast}}%
  \def\diamond{\mathbin{\scalebox{0.6}{$\Diamond$}}}%
  \def\bigcup{\mathop{\mathchoice{\raisebox{-0.3em}{\scalebox{1.7}{$\cup$}}}{\scalebox{1.25}{$\cup$}}{\cup}{\cup}}\displaylimits}%
  \def\bigcap{\mathop{\mathchoice{\raisebox{-0.3em}{\scalebox{1.7}{$\cap$}}}{\scalebox{1.25}{$\cap$}}{\cap}{\cap}}\displaylimits}%
  \def\lesseqqgtr{\mathrel{\lessgtr}}\def\bigoplus{\mathop{\oplus}\displaylimits}%
}
\providecommand{\ssi}{si et seulement si} % abréviation fréquente
\AtBeginDocument{\providecommand{\wideparen}{\overparen}} % arc de cercle

%% FIGFIX-BEGIN v3 — correctifs globaux des figures (docs/onbuch-generation/tools/figfix)
\usepackage{adjustbox}\usepackage{etoolbox}
% toute figure TikZ/pgfplots plus large que la ligne est réduite à la largeur disponible
\BeforeBeginEnvironment{tikzpicture}{\begin{adjustbox}{max width=\linewidth}}
\AfterEndEnvironment{tikzpicture}{\end{adjustbox}}
% légendes pgfplots lisibles par-dessus les courbes
\pgfplotsset{every axis/.append style={legend style={fill=white,fill opacity=0.92,text opacity=1,draw=black!15,inner sep=3pt}}}
% étiquettes d'arêtes (syntaxe edge["..."]) avec fond blanc
\tikzset{every edge quotes/.append style={fill=white,inner sep=1.2pt}}
% bulles de cartes mentales : texte à la ligne dans la bulle, bulle un peu plus grande
\tikzset{every concept/.append style={text width=2.0cm,align=center,minimum size=2.3cm,inner sep=2pt}}
%% FIGFIX-END
\begin{document}

\pagegarde{Lesson 6 — Vocabulary: Words and expressions related to taking part in communal work}{Anglais}{1ère A}{Chapter 1 : Family and Social Life}{ENA06}{2 h}{Cette leçon de vocabulaire présente les mots, collocations et expressions liés au travail communautaire (travaux collectifs, volontariat, entraide). Elle permet aux élèves de décrire, raconter et promouvoir une action communautaire à l'oral comme à l'écrit, en évitant les calques du français.
\vspace{4pt}
\begin{multicols}{2}\small
\begin{itemize}
\item À la fin de cette leçon, je sais nommer les principales activités de travail communautaire (cleaning, farming, building, fundraising).
\item À la fin de cette leçon, je sais employer le vocabulaire du volontariat et de la participation (to volunteer, to take part in, to contribute).
\item À la fin de cette leçon, je sais utiliser les collocations correctes (to carry out a project, to lend a hand, community service).
\item À la fin de cette leçon, je sais éviter les faux amis et les calques du français (to assist ≠ assister à, communal ≠ commun).
\item À la fin de cette leçon, je sais former des mots de la même famille (volunteer, voluntary, voluntarily ; participate, participation, participant).
\item À la fin de cette leçon, je sais prononcer et orthographier correctement le lexique du thème.
\end{itemize}\end{multicols}}

\begin{sommaire}
\begin{multicols}{2}
\begin{enumerate}
\item Découverte du thème : communal work around the world
\item Champ lexical 1 : les activités communautaires
\item Champ lexical 2 : participer, aider, s'organiser
\item Collocations et expressions idiomatiques
\item Faux amis et pièges des francophones
\item Familles de mots, prononciation et orthographe
\item Réemploi en contexte : décrire et promouvoir une action
\item Activité d'intégration
\item Méthodes et exercices résolus
\item Exercices
\item Corrigés des exercices
\item Fiche bilan
\end{enumerate}
\end{multicols}
\end{sommaire}

\begin{objectifs}
\begin{itemize}
\item À la fin de cette leçon, je sais nommer les principales activités de travail communautaire (cleaning, farming, building, fundraising).
\item À la fin de cette leçon, je sais employer le vocabulaire du volontariat et de la participation (to volunteer, to take part in, to contribute).
\item À la fin de cette leçon, je sais utiliser les collocations correctes (to carry out a project, to lend a hand, community service).
\item À la fin de cette leçon, je sais éviter les faux amis et les calques du français (to assist ≠ assister à, communal ≠ commun).
\item À la fin de cette leçon, je sais former des mots de la même famille (volunteer, voluntary, voluntarily ; participate, participation, participant).
\item À la fin de cette leçon, je sais prononcer et orthographier correctement le lexique du thème.
\item À la fin de cette leçon, je sais décrire une journée de travail communautaire dans un court texte ou un dialogue.
\item À la fin de cette leçon, je sais rédiger une affiche ou une annonce invitant la population à participer à un travail collectif.
\end{itemize}
\end{objectifs}

\begin{prerequis}
\begin{itemize}
\item Vocabulaire général de la vie sociale et de la famille (Seconde, Family and Social Life).
\item Présent simple et présent be + -ing pour décrire des habitudes et des actions en cours.
\item Prétérit simple pour raconter un événement passé.
\item Verbes suivis de l'infinitif ou du gérondif (want to, enjoy -ing).
\item Structures de base de la proposition et de l'invitation (Let's..., Would you like to...?).
\end{itemize}
\end{prerequis}

\newpage

\coursec{Découverte du thème : communal work around the world}

\courssub{Situation d'accroche et problématique}

\begin{textebox}[Hook situation]
Every Saturday morning in many Cameroonian neighbourhoods, the whole quarter comes out with machetes, brooms and wheelbarrows for \eng{keep-clean} day. In the North-West and South-West regions, people call this tradition \eng{njangi work} or community labour, while in Britain volunteers join a \eng{clean-up campaign} and in the USA neighbours gather for a \eng{barn raising} or a \eng{community service day}. How do we talk in correct English about working together for the good of the community? This lesson gives you all the words and expressions you need.
\end{textebox}

En français : chaque samedi matin, dans de nombreux quartiers camerounais, tout le voisinage sort avec machettes, balais et brouettes pour la journée « \eng{keep-clean} ». Dans les régions du Nord-Ouest et du Sud-Ouest, on appelle cette tradition le travail de \eng{njangi} ou le travail communautaire ; en Grande-Bretagne, les volontaires participent à une \eng{clean-up campaign} (campagne de nettoyage) et aux États-Unis, les voisins se réunissent pour un \eng{barn raising} (construction collective d'une grange) ou une journée de \eng{community service}. Comment parler en anglais correct du travail accompli ensemble pour le bien de la communauté ? Cette leçon vous donne tous les mots et toutes les expressions nécessaires.

\textbf{Problématique.} Quel vocabulaire anglais permet de décrire, d'organiser et de promouvoir le travail communautaire ? Pour y répondre, nous allons d'abord lire un texte authentique, repérer les mots nouveaux, deviner leur sens grâce au contexte, puis comparer cette tradition à travers le monde anglophone.

\courssub{Lecture du texte support : Keep-Clean Day in Buea}

Lisez le texte suivant une première fois sans dictionnaire. Soulignez mentalement les mots qui vous semblent liés au thème du travail en commun.

\begin{textebox}[Keep-Clean Day in Buea]
Last Saturday, the people of Molyko took part in a keep-clean exercise. Volunteers swept the streets, cleared the gutters and collected rubbish. The council provided wheelbarrows and gloves. Everybody lent a hand, and by noon the quarter was spotless.

The day started early. At six o'clock, the quarter head rang a bell, and men, women and children came out of their houses with brooms, machetes and buckets. Some residents cut the grass along the main road, while others removed the sand and plastic bags that were blocking the gutters before the rainy season. A group of students from the university carried the rubbish to a lorry sent by the council.

Mrs Eposi, a shopkeeper, offered water and fried groundnuts to the workers. “When we work together, the work becomes light,” she said with a smile. “A clean quarter means fewer mosquitoes and less malaria for our children.”

At the end of the morning, the quarter head thanked everybody and announced the next communal work: the repair of the small bridge near the market. “Nobody will pay you,” he said, “but everybody will benefit. That is the spirit of communal work.”
\end{textebox}

\textbf{Glossaire.}

\begin{center}
\begin{tabularx}{\linewidth}{|l|l|X|}
\hline
\rowcolor{popblueL} \textbf{English word} & \textbf{Nature} & \textbf{Français} \\
\hline
keep-clean exercise & nom & journée d'assainissement (nettoyage collectif) \\
\hline
volunteer & nom & volontaire, bénévole \\
\hline
to sweep (swept, swept) & verbe & balayer \\
\hline
gutter & nom & caniveau, rigole d'évacuation des eaux \\
\hline
rubbish & nom (indénombrable) & ordures, déchets \\
\hline
council & nom & municipalité, conseil municipal \\
\hline
wheelbarrow & nom & brouette \\
\hline
to lend a hand & expression & donner un coup de main \\
\hline
spotless & adjectif & impeccable, d'une propreté parfaite \\
\hline
quarter head & nom & chef de quartier \\
\hline
to block & verbe & boucher, obstruer \\
\hline
lorry & nom (GB) & camion \\
\hline
to benefit & verbe & bénéficier, profiter (de) \\
\hline
\end{tabularx}
\end{center}

\textbf{Questions de compréhension.} Répondez en anglais, par des phrases complètes.

\begin{enumerate}
\item What did the people of Molyko do last Saturday?
\item Name three tools that the volunteers used.
\item Who offered food and drink to the workers?
\item Why is a clean quarter good for children's health, according to Mrs Eposi?
\item What will the next communal work be?
\end{enumerate}

\courssub{Définition de la notion : communal work}

\begin{definition}[Communal work]
\cle{Communal work} (also called \cle{community work} or \cle{communal labour}) is work done together by the members of a community for the benefit of all, usually without payment. — Travail accompli ensemble par les membres d'une communauté pour le bien de tous, généralement sans rémunération.
\end{definition}

Observons la phrase finale du texte : \eng{“Nobody will pay you, but everybody will benefit. That is the spirit of communal work.”} (« Personne ne vous paiera, mais tout le monde en profitera. C'est l'esprit du travail communautaire. ») Cette phrase contient les deux caractéristiques essentielles de la notion : l'absence de paiement (\eng{without payment}) et le bénéfice collectif (\eng{for the benefit of all}).

\begin{aretenir}[Les mots outils du thème]
\begin{itemize}
\item \cle{community} (nom) : la communauté ; \eng{the local community} = la communauté locale.
\item \cle{communal} (adjectif) : communautaire, collectif ; \eng{communal work}, \eng{communal land} (terrain communautaire).
\item \cle{voluntary} (adjectif) : bénévole, volontaire ; \eng{voluntary work} = travail bénévole.
\item \cle{volunteer} : nom (\eng{a volunteer}, un bénévole) et verbe (\eng{to volunteer}, se porter volontaire).
\end{itemize}
\end{aretenir}

\begin{attention}[Faux ami : community et communal]
Le français « communal » évoque souvent la mairie ou l'administration (la salle communale, le stade communal). En anglais, \eng{communal} signifie « partagé par le groupe, collectif » : \eng{communal work} = travaux collectifs du quartier ou du village. De même, \eng{community} ne se traduit pas toujours par « communauté » au sens administratif : \eng{the community} désigne simplement « les gens du quartier, les habitants ». Enfin, attention : \eng{rubbish} est indénombrable — on dit \eng{some rubbish}, jamais \eng{a rubbish}.
\end{attention}

\courssub{Deviner le sens d'un mot nouveau grâce au contexte}

Avant de consulter le glossaire, un bon lecteur essaie de deviner le sens des mots inconnus. C'est une compétence essentielle pour l'épreuve de compréhension de l'écrit.

\begin{methode}[Guessing meaning from context : deviner le sens à partir du contexte]
\begin{enumerate}
\item \textbf{Lisez la phrase entière}, pas seulement le mot inconnu : le sens se trouve souvent avant ou après.
\item \textbf{Identifiez la nature du mot} : est-ce un nom, un verbe, un adjectif ? La place dans la phrase et les terminaisons (\eng{-ed}, \eng{-ly}, \eng{-tion}) vous aident.
\item \textbf{Cherchez les indices du contexte} : mots de la même famille, synonymes, exemples, contrastes (\eng{but}, \eng{while}), ou conséquences logiques.
\item \textbf{Formulez une hypothèse} en français ou en anglais simple.
\item \textbf{Vérifiez ensuite dans le glossaire} ou le dictionnaire, et corrigez votre hypothèse si nécessaire.
\end{enumerate}
\end{methode}

\begin{exoresolu}[Deviner le sens de “spotless”]
Énoncé. Dans la phrase \eng{Everybody lent a hand, and by noon the quarter was spotless}, devinez le sens de l'adjectif \eng{spotless} sans dictionnaire, puis vérifiez.
\tcblower
Solution détaillée.
\begin{enumerate}
\item Je lis la phrase entière : tout le monde a donné un coup de main, et à midi le quartier était \eng{spotless}.
\item Nature du mot : \eng{spotless} suit le verbe \eng{was} et décrit \eng{the quarter} ; c'est un adjectif. Il est formé de \eng{spot} (tache) + suffixe \eng{-less} (sans), comme \eng{hopeless} (sans espoir) ou \eng{careless} (sans attention).
\item Indices du contexte : les volontaires ont balayé, débouché les caniveaux et ramassé les ordures ; la conséquence logique est que le quartier est devenu très propre.
\item Hypothèse : \eng{spotless} = « sans tache », donc « très propre, impeccable ».
\item Vérification dans le glossaire : \eng{spotless} = impeccable, d'une propreté parfaite. Mon hypothèse était correcte.
\end{enumerate}
\end{exoresolu}

\begin{exercice}[À vous de deviner]
En appliquant la méthode en cinq étapes, devinez le sens des mots suivants tels qu'ils apparaissent dans le texte, puis vérifiez dans le glossaire : 1. \eng{wheelbarrow} ; 2. \eng{to lend a hand} ; 3. \eng{gutters} ; 4. \eng{quarter head}.
\end{exercice}

\begin{corrige}[Exercice : À vous de deviner]
1. \eng{Wheelbarrow} : nom d'objet ; le contexte (\eng{the council provided wheelbarrows and gloves}) indique du matériel de travail fourni par la municipalité ; il sert à transporter le sable et les ordures. Sens : brouette. 2. \eng{To lend a hand} : expression verbale ; le contexte (\eng{everybody}) montre que chacun a participé. Sens : donner un coup de main, aider. 3. \eng{Gutters} : nom pluriel ; ils sont \eng{blocked} par le sable et les sacs plastiques \eng{before the rainy season} ; il s'agit donc de canaux d'évacuation de l'eau : les caniveaux. 4. \eng{Quarter head} : nom de personne ; c'est lui qui sonne la cloche, remercie et annonce le prochain chantier : le chef de quartier.
\end{corrige}

\courssub{Repérage et classification des mots nouveaux liés au thème}

Relevons dans le texte les mots et expressions qui appartiennent au champ lexical du travail communautaire, et classons-les par catégories.

\begin{center}
\begin{tabularx}{\linewidth}{|X|X|X|}
\hline
\rowcolor{popblueL} \textbf{English} & \textbf{Français} & \textbf{Exemple tiré ou construit} \\
\hline
volunteers & volontaires, bénévoles & \eng{Volunteers swept the streets.} « Les volontaires ont balayé les rues. » \\
\hline
clean-up campaign & campagne de nettoyage & \eng{Our school launched a clean-up campaign.} « Notre école a lancé une campagne de nettoyage. » \\
\hline
to sweep (swept, swept) & balayer & \eng{They swept the market square.} « Ils ont balayé la place du marché. » \\
\hline
to clear the gutters & déboucher les caniveaux & \eng{We cleared the gutters before the rains.} « Nous avons débouché les caniveaux avant les pluies. » \\
\hline
to collect rubbish & ramasser les ordures & \eng{The children collected rubbish in bags.} « Les enfants ont ramassé les ordures dans des sacs. » \\
\hline
to lend a hand & donner un coup de main & \eng{Everybody lent a hand.} « Tout le monde a donné un coup de main. » \\
\hline
to take part in & participer à & \eng{The people took part in the exercise.} « Les gens ont participé à l'opération. » \\
\hline
to benefit & bénéficier, profiter & \eng{Everybody will benefit.} « Tout le monde en profitera. » \\
\hline
\end{tabularx}
\end{center}

\begin{propriete}[Deux verbes irréguliers à connaître]
\begin{center}
\begin{tabular}{|l|l|l|l|}
\hline
\rowcolor{popgreenL} \textbf{Base verbale} & \textbf{Prétérit} & \textbf{Participe passé} & \textbf{Français} \\
\hline
sweep & swept & swept & balayer \\
\hline
lend & lent & lent & prêter \\
\hline
\end{tabular}
\end{center}
\eng{They swept the streets.} « Ils ont balayé les rues. » — \eng{Everybody lent a hand.} « Tout le monde a prêté main-forte. » Attention : \eng{sweeped} et \eng{lended} n'existent pas.
\end{propriete}

\courssub{Champ lexical 2 : participer, aider, s'organiser}

Au-delà des actions de nettoyage, le travail communautaire suppose des verbes d'organisation et de participation.

\begin{center}
\begin{tabularx}{\linewidth}{|X|X|X|}
\hline
\rowcolor{popblueL} \textbf{English} & \textbf{Français} & \textbf{Collocation / Exemple} \\
\hline
to organise / to coordinate & organiser / coordonner & \eng{The quarter head organised the work.} \\
\hline
to mobilise & mobiliser & \eng{They mobilised the youth for the bridge repair.} \\
\hline
to contribute & contribuer, cotiser & \eng{Everyone contributed time or money.} \\
\hline
to raise funds / money & collecter des fonds & \eng{The njangi group raised funds for the well.} \\
\hline
to supervise & superviser & \eng{The engineer supervised the construction.} \\
\hline
to participate / to take part in & participer & \eng{We participated in the clean-up.} \\
\hline
to help (sb) with (sth) & aider qn à qc & \eng{He helped with the cooking.} \\
\hline
to join forces & unir ses forces & \eng{The two villages joined forces.} \\
\hline
\end{tabularx}
\end{center}

\begin{attention}[Pièges de prépositions et calques]
\begin{itemize}
\item \eng{To take part \textbf{in}} (et non \eng{to} ou \eng{at}) : \eng{He took part in the meeting.}
\item \eng{To benefit \textbf{from}} : \eng{The village benefited from the new well.}
\item \eng{To help \textbf{with}} (quelque chose) / \eng{to help (sb) \textbf{to} do / do (sth)} : \eng{She helped with the cleaning.} / \eng{She helped me carry the buckets.}
\item \eng{To participate \textbf{in}} (formel) vs \eng{to take part \textbf{in}} (courant). Évitez \eng{*participate to}.
\item \eng{To volunteer \textbf{for}} (une tâche) / \eng{to volunteer \textbf{to} do} : \eng{He volunteered for the job.} / \eng{She volunteered to cook.}
\end{itemize}
\end{attention}

\courssub{Collocations et expressions idiomatiques}

Les collocations (associations de mots figées) donnent un anglais naturel.

\begin{center}
\begin{tabularx}{\linewidth}{|X|X|X|}
\hline
\rowcolor{popblueL} \textbf{Collocation / Idiom} & \textbf{Sens} & \textbf{Exemple} \\
\hline
\eng{to lend a hand} & donner un coup de main & \eng{Can you lend a hand with the painting?} \\
\hline
\eng{many hands make light work} & à plusieurs, le travail est léger (proverbe) & \eng{Let's do it together: many hands make light work.} \\
\hline
\eng{to do one's bit / share} & faire sa part & \eng{Everyone did their bit for the community.} \\
\hline
\eng{to pull together} & serrer les coudes, unir ses efforts & \eng{In times of crisis, the community pulls together.} \\
\hline
\eng{to give back to the community} & rendre à la communauté & \eng{He volunteers to give back to his community.} \\
\hline
\eng{community spirit} & esprit communautaire, solidarité & \eng{The keep-clean day shows great community spirit.} \\
\hline
\eng{unpaid / voluntary work} & travail non rémunéré / bénévole & \eng{Voluntary work enriches your CV.} \\
\hline
\end{tabularx}
\end{center}

\courssub{Faux amis et pièges des francophones : focus}

\begin{attention}[Faux amis et erreurs fréquentes]
\begin{itemize}
\item \textbf{Communitary} \eng{\neq} \eng{communal} / \eng{community}. \eng{Communitary} n'est pas un adjectif courant en anglais standard ; on dit \eng{communal work}, \eng{community project}.
\item \textbf{Volontaire} (adj) $\rightarrow$ \eng{voluntary} (ex. \eng{voluntary work}). \textbf{Volontaire} (nom) $\rightarrow$ \eng{volunteer}. Ne dites pas \eng{*a voluntary} pour « un bénévole ».
\item \textbf{Bénévole} (adj) $\rightarrow$ \eng{voluntary} / \eng{unpaid}. \eng{Benevolent} existe mais signifie « bienveillant, charitable » (souvent pour une organisation : \eng{a benevolent society}).
\item \textbf{Assainissement} $\rightarrow$ \eng{sanitation} (système) ou \eng{clean-up} / \eng{keep-clean exercise} (action). Pas \eng{*assainissement} $\neq$ \eng{assassination} (meurtre) !
\item \textbf{Participer} $\rightarrow$ \eng{take part in}. \eng{Participate} est plus formel. Évitez \eng{*assist to} (qui signifie « assister à » = être spectateur).
\item \textbf{Profiter} $\rightarrow$ \eng{benefit from}. \eng{Profit} est un nom (bénéfice financier) ou verbe intransitif (tirer profit), mais \eng{profit from} est moins courant que \eng{benefit from} pour « bénéficier de ».
\end{itemize}
\end{attention}

\courssub{Familles de mots, prononciation et orthographe}

Travailler sur les familles de mots permet d'enrichir son vocabulaire de manière exponentielle.

\begin{center}
\begin{tabularx}{\linewidth}{|l|l|l|l|X|}
\hline
\rowcolor{popblueL} \textbf{Verbe} & \textbf{Nom (action/état)} & \textbf{Nom (agent)} & \textbf{Adjectif} & \textbf{Exemple / Prononciation (approx. fr.)} \\
\hline
organise & organisation & organiser & organised & \eng{The organisation was perfect.} (or-ga-ni-za-tion) \\
\hline
participate & participation & participant & participative & \eng{Active participation is required.} (par-ti-ci-pa-tion) \\
\hline
volunteer & volunteering & volunteer & voluntary & \eng{Voluntary work.} (vo-lon-te-ri) \\
\hline
contribute & contribution & contributor & contributory & \eng{A major contribution.} (kon-tri-biu-chon) \\
\hline
mobilise & mobilisation & mobiliser & mobilised & \eng{General mobilisation.} (mo-bi-li-za-tion) \\
\hline
benefit & benefit / beneficiary & beneficiary & beneficial & \eng{Beneficial effects.} (be-ne-fi-ciel) \\
\hline
clean & clean-up / cleaning & cleaner & clean & \eng{A clean-up campaign.} \\
\hline
\end{tabularx}
\end{center}

\begin{propriete}[Règles d'orthographe et de prononciation clés]
\begin{itemize}
\item \textbf{-ise / -ize} : En anglais britannique (norme camerounaise), on écrit \eng{organise}, \eng{mobilise}, \eng{realise}. La forme \eng{-ize} est américaine (acceptée mais moins usitée au Cameroun).
\item \textbf{Double consonne} : \eng{participate} (un seul \eng{t}), mais \eng{contribution} (double \eng{t} ? Non, \eng{contribution} un seul \eng{t}). Attention à \eng{organise} $\rightarrow$ \eng{organisation} (apparition du \eng{t}).
\item \textbf{Accentuation} : \eng{volunteer} (verbe/nom) \eng{[vɒlənˈtɪə]} $\rightarrow$ accent sur la dernière syllabe. \eng{voluntary} (adj) \eng{[ˈvɒləntri]} $\rightarrow$ accent sur la première, \eng{a} muet. \eng{beneficial} \eng{[benɪˈfɪʃl]} $\rightarrow$ accent sur \eng{fi}.
\item \textbf{Indénombrables} : \eng{rubbish}, \eng{litter}, \eng{work} (au sens « ouvrage »), \eng{equipment}, \eng{furniture}. Jamais de \eng{a/an} devant, pas de pluriel en \eng{-s}.
\end{itemize}
\end{propriete}

\begin{exercice}[Famille de mots : complétez le tableau]
Complétez les cases vides. Utilisez un dictionnaire si nécessaire.

\begin{center}
\begin{tabular}{|l|l|l|l|l|}
\hline
\rowcolor{popgreenL} \textbf{Verbe} & \textbf{Nom (action)} & \textbf{Nom (personne)} & \textbf{Adjectif} & \textbf{Adverbe} \\
\hline
organise & \ldots & organiser & \ldots & \ldots \\
\hline
\ldots & participation & \ldots & participative & \ldots \\
\hline
volunteer & \ldots & \ldots & voluntary & \ldots \\
\hline
\ldots & contribution & contributor & \ldots & \ldots \\
\hline
benefit & \ldots & beneficiary & \ldots & \ldots \\
\hline
\end{tabular}
\end{center}
\end{exercice}

\begin{corrige}[Exercice : Famille de mots]
\begin{center}
\begin{tabular}{|l|l|l|l|l|}
\hline
\rowcolor{popgreenL} \textbf{Verbe} & \textbf{Nom (action)} & \textbf{Nom (personne)} & \textbf{Adjectif} & \textbf{Adverbe} \\
\hline
organise & organisation & organiser & organised & \ldots \\
\hline
participate & participation & participant & participative & participatively \\
\hline
volunteer & volunteering & volunteer & voluntary & voluntarily \\
\hline
contribute & contribution & contributor & contributory & \ldots \\
\hline
benefit & benefit & beneficiary & beneficial & beneficially \\
\hline
\end{tabular}
\end{center}
\end{corrige}

\courssub{Carte mentale du thème}

La carte mentale ci-dessous organise le champ lexical du travail communautaire en cinq grandes branches d'activités. Vous la compléterez au fil de la leçon.

\begin{popfigure}
\begin{tikzpicture}[mindmap, grow cyclic, every node/.style={concept, execute at begin node=\hskip0pt}, concept color=popblue!40, root concept/.append style={concept color=poporange!60, font=\bfseries}, level 1/.append style={level distance=3.5cm, sibling angle=72}, text=popdark, font=\small]
\node[root concept]{COMMUNAL WORK}
  child[concept color=popblue!30]{ node[align=center] {Cleaning\\ \tiny sweep, clear gutters,\\ collect rubbish, litter picking} }
  child[concept color=popgreen!35]{ node[align=center] {Building \& Repair\\ \tiny repair bridge, build school,\\ maintain roads, paint walls} }
  child[concept color=popgold!45]{ node[align=center] {Farming\\ \tiny harvest help, clear fields,\\ plant trees, njangi farm} }
  child[concept color=poppurple!30]{ node[align=center] {Fundraising\\ \tiny raise money, njangi,\\ contributions, donations} }
  child[concept color=poppink!35]{ node[align=center] {Helping the Needy\\ \tiny visit sick, share food,\\ support orphans, care} };
\end{tikzpicture}
\legende{Carte mentale : les cinq grandes familles d'activités du travail communautaire.}
\end{popfigure}

\courssub{Brainstorming : communal work in your quarter or village}

\begin{experience}[Brainstorming : What communal work exists in your area?]
En groupes de quatre, listez en deux minutes tous les travaux communautaires qui existent dans votre quartier ou votre village, puis complétez les phrases suivantes en anglais :
\begin{enumerate}
\item \eng{In my quarter, people take part in\ldots} (road repair — réparation des routes)
\item \eng{Every year, the community organises\ldots} (well digging — creusement d'un puits)
\item \eng{The parents helped with\ldots} (school construction — construction d'une école)
\item \eng{During the season, neighbours offer\ldots} (harvest help — aide à la récolte)
\end{enumerate}
Chaque groupe présente ensuite une phrase à la classe : \eng{In our village, young people repair the road after the rainy season.} « Dans notre village, les jeunes réparent la route après la saison des pluies. »
\end{experience}

Vocabulaire utile pour ce brainstorming :

\begin{center}
\begin{tabularx}{\linewidth}{|X|X|}
\hline
\rowcolor{popblueL} \textbf{English} & \textbf{Français} \\
\hline
road repair / maintenance & réparation / entretien des routes \\
\hline
well digging & creusement d'un puits \\
\hline
school construction & construction d'une école \\
\hline
harvest help & aide à la récolte \\
\hline
bridge repair & réparation d'un pont \\
\hline
tree planting & plantation d'arbres \\
\hline
market cleaning & nettoyage du marché \\
\hline
drainage clearing & curage des caniveaux \\
\hline
\end{tabularx}
\end{center}

\courssub{Comparaison culturelle : communal work around the world}

Le travail communautaire existe dans toutes les cultures, mais chaque pays lui a donné un nom et une forme particulière.

\begin{center}
\begin{tabularx}{\linewidth}{|l|X|X|}
\hline
\rowcolor{popblueL} \textbf{Pays / Région} & \textbf{Expression} & \textbf{Signification} \\
\hline
Cameroun & keep-clean day, njangi work, community labour & journée d'assainissement ; travail collectif du quartier ou du village ; épargne solidaire \\
\hline
Royaume-Uni & clean-up campaign, volunteering & campagne de nettoyage organisée par des bénévoles \\
\hline
États-Unis & community service, barn raising & service communautaire (souvent obligatoire pour lycéens) ; construction collective d'une grange (tradition rurale) \\
\hline
Kenya & harambee & « tirons ensemble » : collecte et travail communautaires \\
\hline
Indonésie & gotong royong & « porter ensemble » : entraide collective du village \\
\hline
\end{tabularx}
\end{center}

\begin{savaistu}[Harambee : un mot sur les armoiries du Kenya]
Au Kenya, le mot \cle{harambee}, qui signifie littéralement « tirons ensemble » en swahili, est inscrit sur les armoiries du pays. Il désigne les collectes communautaires organisées pour construire une école, payer des frais médicaux ou financer un projet de village. Le premier président du Kenya, Jomo Kenyatta, en a fait la devise nationale : tout le peuple tire dans la même direction. Au Cameroun, la tradition du \eng{njangi} — l'épargne et l'entraide en groupe — repose sur exactement le même esprit : \eng{many hands make light work}, « à plusieurs, le travail devient léger ».
\end{savaistu}

\begin{exemplebox}[Comparer les traditions en anglais correct]
\eng{In Cameroon, communal work is often called keep-clean day, while in Kenya people say harambee.} « Au Cameroun, le travail communautaire s'appelle souvent keep-clean day, tandis qu'au Kenya on dit harambee. »

\eng{In the USA, high school students often do community service to graduate.} « Aux États-Unis, les lycéens font souvent du community service pour obtenir leur diplôme. »

\eng{Community service is unpaid work for the good of the community.} « Le community service est un travail non rémunéré pour le bien de la communauté. »
\end{exemplebox}

\courssub{Réemploi en contexte : décrire et promouvoir une action}

\courssub{Méthode : Rédiger un appel à la participation (Notice / Poster)}

\begin{methode}[Comment rédiger un avis d'appel au travail communautaire (Notice)]
\begin{enumerate}
\item \textbf{Titre accrocheur} : \eng{COMMUNAL WORK NOTICE} ou \eng{KEEP-CLEAN DAY : CALL TO ACTION}.
\item \textbf{Informations essentielles (WH-questions)} : \eng{What?} (activité), \eng{When?} (date, heure), \eng{Where?} (lieu de rendez-vous), \eng{Who?} (public ciblé : \eng{All residents}, \eng{Youths}, \eng{Parents}).
\item \textbf{Verbes d'action à l'impératif} : \eng{Come}, \eng{Bring}, \eng{Join}, \eng{Help}, \eng{Participate}.
\item \textbf{Motivation} : \eng{Many hands make light work.} / \eng{For a cleaner, healthier quarter.} / \eng{Your community needs you.}
\item \textbf{Contact / Signature} : \eng{Quarter Head}, \eng{Youth President}, \eng{Date}.
\end{enumerate}
\end{methode}

\begin{exemplebox}[Modèle d'avis (Notice) — Keep-Clean Day]
\begin{textebox}[Notice Board — Molyko Quarter]
\textbf{COMMUNAL WORK NOTICE}

\textbf{KEEP-CLEAN EXERCISE — NEXT SATURDAY}

\textbf{What:} General cleaning of streets, gutters and market square.
\textbf{When:} Saturday, 6:00 a.m. sharp.
\textbf{Where:} Meet at the Quarter Head's office.
\textbf{Who:} All residents (men, women, youths, children).
\textbf{Bring:} Brooms, machetes, wheelbarrows, gloves, buckets.
\textbf{Provided:} Water, snacks, rubbish bags, council lorry.

\textit{Many hands make light work. A clean quarter is a healthy quarter.}
\textbf{COME ONE, COME ALL!}

\eng{The Quarter Head} \\
\eng{Molyko, Buea}
\end{textebox}
\end{exemplebox}

\begin{exercice}[Production écrite : Votre avis]
Rédigez un avis (\eng{Notice}) pour mobiliser votre quartier ou votre école pour l'une des actions suivantes (choisissez) :
\begin{enumerate}
\item Réparation du pont près du marché (\eng{Bridge Repair}).
\item Plantation d'arbres dans la cour de l'école (\eng{Tree Planting Campaign}).
\item Creusement d'un puits pour le village (\eng{Well Digging Project}).
\end{enumerate}
Respectez la structure de la méthode (Titre, What/When/Where/Who/Bring, Motivation, Signature). 80–100 mots.
\end{exercice}

\begin{corrige}[Exercice : Production écrite — Modèle Bridge Repair]
\begin{textebox}[Notice Board — Your Quarter]
\textbf{COMMUNAL WORK NOTICE}

\textbf{BRIDGE REPAIR — COMMUNITY ACTION}

\textbf{What:} Repair of the wooden bridge near the main market.
\textbf{When:} Saturday 15th June, 7:00 a.m.
\textbf{Where:} Meeting point at the bridge site.
\textbf{Who:} All able-bodied men and women, youths welcome.
\textbf{Bring:} Hammers, nails, saws, planks, sand, cement bags.
\textbf{Provided:} Technical supervision by the council engineer, refreshments.

\textit{Many hands make light work. Safe passage for all!}
\textbf{YOUR PRESENCE IS VITAL.}

\eng{The Quarter Head} \\
\eng{Date: 10th June}
\end{textebox}
\end{corrige}

\courssub{Expression orale : Promouvoir une action (Role-play)}

\begin{experience}[Role-play : Convaincre un voisin réticent]
En binôme. \textbf{Élève A} : Vous êtes le président des jeunes du quartier. Vous voulez que tout le monde participe au \eng{keep-clean} de samedi. \textbf{Élève B} : Vous êtes un habitant fatigué qui préfère dormir ou regarder le match. Vous trouvez des excuses (\eng{I'm tired}, \eng{It's not my job}, \eng{The council should pay people}).
\textbf{Objectif A} : Utiliser \eng{lend a hand}, \eng{benefit from}, \eng{community spirit}, \eng{many hands make light work}, \eng{take part in} pour convaincre B.
\textbf{Objectif B} : Répondre poliment mais fermement, puis finalement accepter.
\textbf{Exemple de début :} \eng{A: "Good morning Papa Paul. We need you for the keep-clean on Saturday."} \eng{B: "Oh, my back hurts... and anyway, I pay my taxes."}
\end{experience}

\begin{aretenir}[Bilan complet de la leçon 6]
\begin{itemize}
\item \textbf{Notion clé} : \cle{Communal work} = travail collectif, non rémunéré, au profit de tous (\eng{community work}, \eng{communal labour}).
\item \textbf{Verbes d'action (irréguliers)} : \eng{sweep / swept / swept}, \eng{lend / lent / lent}, \eng{build / built / built}, \eng{dig / dug / dug}.
\item \textbf{Collocations indispensables} : \eng{take part in}, \eng{lend a hand}, \eng{clear the gutters}, \eng{collect rubbish}, \eng{raise funds}, \eng{join forces}, \eng{many hands make light work}.
\item \textbf{Familles de mots} : \eng{organise/organisation/organiser}, \eng{participate/participation/participant}, \eng{volunteer/volunteering/voluntary}, \eng{contribute/contribution/contributor}, \eng{benefit/beneficiary/beneficial}.
\item \textbf{Faux amis} : \eng{communitary} (non), \eng{voluntary} (adj) vs \eng{volunteer} (n/v), \eng{assist} (assister à) vs \eng{take part in} (participer), \eng{rubbish} (indénombrable).
\item \textbf{Culture} : \eng{keep-clean day / njangi} (Cameroun), \eng{harambee} (Kenya), \eng{gotong royong} (Indonésie), \eng{community service / barn raising} (USA), \eng{clean-up campaign} (UK).
\item \textbf{Production} : Savoir rédiger un \eng{Notice} (avis d'appel) et argumenter à l'oral pour mobiliser.
\end{itemize}
\end{aretenir}

\coursec{Champ lexical 1 : les activités communautaires}

Dans la section précédente, nous avons découvert que le travail communautaire — \eng{communal work} — existe dans toutes les cultures : du \eng{gotong royong} indonésien aux \eng{barn raisings} américains, en passant par nos propres journées de \eng{clean-up} à Yaoundé ou à Bamenda. Mais pour décrire ces actions en anglais, il faut d'abord maîtriser le vocabulaire de base. C'est l'objet de cette section : nommer les activités, les actions concrètes, les outils et les lieux du travail collectif.

\courssub{Observer d'abord : un texte pour découvrir le lexique}

Lisons d'abord ce court article de presse locale. Soulignez mentalement tous les mots qui désignent une activité, un outil ou un lieu.

\begin{textebox}[A clean-up day in Mfoundi — community news]
Last Saturday, the inhabitants of the Nkolndongo neighbourhood organised a big \textbf{clean-up campaign} before the rainy season. From six o'clock in the morning, dozens of volunteers gathered on the \textbf{market square} with their \textbf{brooms}, \textbf{shovels}, \textbf{wheelbarrows} and \textbf{buckets}.

The youths formed small teams. One group swept the streets while another cleared the \textbf{gutters}, which were full of plastic bags and sand. A third group collected rubbish and carried it away in wheelbarrows. Near the \textbf{health centre}, women painted the fence and repaired the benches in front of the \textbf{village hall}. A few men, wearing thick \textbf{gloves}, dug around the public \textbf{water tap} so that the waste water could flow away.

The quarter head explained that this \textbf{work party} was only the beginning. Next month, the community will start a \textbf{building project}: a new classroom for the primary school. There will also be a \textbf{fundraising event} to buy cement and iron sheets. “When we work together,” he said, “no task is too heavy. Last year we even repaired the small bridge on the road to Emana, and two years ago we finished the \textbf{well digging} near the chief's palace.”

By midday, the streets were clean, the gutters were clear, and everybody shared a meal of \emph{eru} and \emph{fufu} offered by the women's association. The next \textbf{working bee} is already planned: in April, volunteers will plant trees along the main road and harvest crops in the community farm.
\end{textebox}

\textbf{Glossaire}

\begin{center}
\begin{tabularx}{\linewidth}{|l|l|X|}\hline
\rowcolor{popblueL} \textbf{Mot anglais} & \textbf{Nature} & \textbf{Traduction} \\ \hline
a clean-up campaign & nom & une campagne de nettoyage \\ \hline
a gutter & nom & un caniveau \\ \hline
a work party & nom & une corvée collective \\ \hline
a fence & nom & une clôture \\ \hline
waste water & nom & les eaux usées \\ \hline
to flow away & verbe & s'écouler \\ \hline
a task & nom & une tâche \\ \hline
well digging & nom & le creusement d'un puits \\ \hline
\end{tabularx}
\end{center}

\textbf{Questions de compréhension}
\begin{enumerate}
\item Why did the inhabitants organise a clean-up campaign?
\item Name three tools the volunteers brought with them.
\item What will the community build next month?
\item What did the volunteers share at midday?
\end{enumerate}

\courssub{Les grandes activités communautaires}

\begin{definition}[Community activities — les activités collectives]
Une \cle{activité communautaire} est un travail réalisé ensemble, gratuitement, pour le bien de tout le quartier ou du village. En anglais, ces activités ont des noms précis qu'il faut apprendre comme des blocs lexicaux.
\end{definition}

\begin{center}
\begin{tabularx}{\linewidth}{|X|X|X|}\hline
\rowcolor{popblueL} \textbf{English} & \textbf{Français} & \textbf{Remarque} \\ \hline
a clean-up campaign & une campagne de nettoyage & \eng{clean-up} s'écrit avec un trait d'union \\ \hline
a working bee / a work party & une corvée collective & \eng{working bee} : surtout américain et australien \\ \hline
communal labour & le travail communautaire & terme général, soutenu (orthographe GB : \eng{labour}) \\ \hline
a fundraising event & une collecte de fonds & \eng{to raise funds} = collecter des fonds \\ \hline
a building project & un chantier de construction & \eng{a building site} = le chantier (le lieu) \\ \hline
road maintenance & l'entretien des routes & \eng{maintenance} se prononce « méin-te-nens » \\ \hline
well digging & le creusement d'un puits & \eng{a well} = un puits ; \eng{to dig} = creuser \\ \hline
\end{tabularx}
\end{center}

\begin{exemplebox}[Ces activités en contexte]
\eng{Our school organised a clean-up campaign last term.} « Notre école a organisé une campagne de nettoyage le trimestre dernier. »

\eng{The village held a fundraising event to buy a new generator.} « Le village a organisé une collecte de fonds pour acheter un nouveau groupe électrogène. »

\eng{Road maintenance is everyone's business in our subdivision.} « L'entretien des routes est l'affaire de tous dans notre arrondissement. »
\end{exemplebox}

\courssub{Les actions concrètes : les verbes du travail collectif}

\begin{propriete}[Verbes d'action + complément]
En anglais comme en français, chaque action s'exprime par un \cle{verbe} suivi de son \cle{complément}. Il faut mémoriser le verbe AVEC son complément habituel (collocation), car on ne peut pas toujours traduire mot à mot.
\end{propriete}

\begin{center}
\begin{tabularx}{\linewidth}{|X|X|}\hline
\rowcolor{popgreenL} \textbf{English} & \textbf{Français} \\ \hline
to sweep the streets & balayer les rues \\ \hline
to clear the gutters & déboucher les caniveaux \\ \hline
to collect rubbish & ramasser les ordures \\ \hline
to plant trees & planter des arbres \\ \hline
to paint the fence & peindre la clôture \\ \hline
to repair the bridge & réparer le pont \\ \hline
to harvest crops & récolter (les cultures) \\ \hline
to carry water & transporter de l'eau \\ \hline
to dig a well & creuser un puits \\ \hline
\end{tabularx}
\end{center}

\begin{exemplebox}[Exemples traduits]
\eng{The youths cleared the gutters before the rainy season.} « Les jeunes ont débouché les caniveaux avant la saison des pluies. »

\eng{Women carried water in buckets to the building site.} « Les femmes ont transporté de l'eau dans des seaux jusqu'au chantier. »

\eng{Every Saturday, pupils sweep the streets around the market.} « Chaque samedi, les élèves balayent les rues autour du marché. »

\eng{The farmers harvested the crops together in November.} « Les cultivateurs ont récolté ensemble en novembre. »
\end{exemplebox}

\begin{attention}[to clear, pas « to clean » pour les caniveaux]
Pour les caniveaux, l'anglais dit \eng{to clear the gutters} (déboucher, dégager), rarement \eng{to clean}. De même, on dit \eng{to collect rubbish} et non « to pick dirts » : \eng{rubbish} (GB) ou \eng{garbage/trash} (US) sont les mots corrects pour « ordures ». Attention aussi : \eng{to repair} (réparer) ne se confond pas avec \eng{to repaint} (repeindre).
\end{attention}

\courssub{Les outils et le matériel}

\begin{definition}[Tools — les outils]
Un \cle{outil} se dit \eng{a tool} ; le matériel se dit \eng{equipment} (mot \emph{indénombrable} : jamais de « s » au pluriel !).
\end{definition}

\begin{center}
\begin{tabularx}{\linewidth}{|X|X|X|}\hline
\rowcolor{poporangeL} \textbf{English} & \textbf{Français} & \textbf{Prononciation} \\ \hline
a broom & un balai & « broum » \\ \hline
a machete & une machette & « ma-ché-ti » \\ \hline
a wheelbarrow & une brouette & « ouil-ba-ro » \\ \hline
a shovel & une pelle & « che-veul » \\ \hline
a rake & un râteau & « réik » \\ \hline
gloves (pl.) & des gants & « gleuvz » \\ \hline
a bucket & un seau & « beu-kèt » \\ \hline
\end{tabularx}
\end{center}

\begin{methode}[Mémoriser par association outil-action]
Pour retenir durablement ce vocabulaire, associez toujours chaque outil à l'action correspondante, comme une flèche mentale :
\begin{enumerate}
\item \eng{a broom} $\rightarrow$ \eng{to sweep} (un balai sert à balayer)
\item \eng{a shovel} $\rightarrow$ \eng{to dig} (une pelle sert à creuser)
\item \eng{a rake} $\rightarrow$ \eng{to rake the leaves} (un râteau sert à ratisser)
\item \eng{a wheelbarrow} $\rightarrow$ \eng{to carry / to transport} (une brouette sert à transporter)
\item \eng{a bucket} $\rightarrow$ \eng{to carry water} (un seau sert à porter l'eau)
\item \eng{a machete} $\rightarrow$ \eng{to cut the grass} (une machette sert à couper l'herbe)
\end{enumerate}
Récitez ces paires à voix haute : la mémoire orale est la plus solide.
\end{methode}

\courssub{Les lieux concernés}

Le travail communautaire se déroule dans des lieux publics précis :

\begin{center}
\begin{tabularx}{\linewidth}{|X|X|}\hline
\rowcolor{poppurpleL} \textbf{English} & \textbf{Français} \\ \hline
the market square & la place du marché \\ \hline
the health centre & le centre de santé \\ \hline
the water tap & le robinet public \\ \hline
the village hall & la salle des fêtes / la mairie de village \\ \hline
the chief's palace & le palais du chef \\ \hline
the community farm & la ferme communautaire \\ \hline
\end{tabularx}
\end{center}

\begin{attention}[Orthographe britannique : centre, litre, metre]
En anglais britannique (celui du programme camerounais), on écrit \eng{centre} et non « center » (américain). De même : \eng{theatre}, \eng{metre}. Attention aussi à \eng{hall} (se prononce « hol ») qui ne veut pas dire « halle » au sens de marché couvert : \eng{the village hall} est la salle des fêtes, le lieu de réunion du village.
\end{attention}

\courssub{Tableau de synthèse : activité, outil, lieu}

\begin{popfigure}
\begin{tikzpicture}[font=\small]
\fill[popblueL, rounded corners=3pt] (0,0) rectangle (4.6,1);
\node[align=center] at (2.3,0.5) {\textbf{Activity}};
\fill[poporangeL, rounded corners=3pt] (4.8,0) rectangle (9.4,1);
\node[align=center] at (7.1,0.5) {\textbf{Tools}};
\fill[popgreenL, rounded corners=3pt] (9.6,0) rectangle (14.2,1);
\node[align=center] at (11.9,0.5) {\textbf{Place}};
\foreach \y/\a/\b/\c in {
  -1.1/{sweep the streets}/{broom, wheelbarrow}/{market square},
  -2.3/{clear the gutters}/{shovel, gloves}/{the main road},
  -3.5/{paint the fence}/{brush, bucket}/{health centre},
  -4.7/{dig a well}/{shovel, bucket}/{near the water tap},
  -5.9/{plant trees}/{machete, rake}/{community farm}
}{
  \node[align=center, anchor=north] at (2.3,\y) {\a};
  \node[align=center, anchor=north] at (7.1,\y) {\b};
  \node[align=center, anchor=north] at (11.9,\y) {\c};
  \draw[popline] (0,\y+0.35) -- (14.2,\y+0.35);
}
\draw[popline] (0,0) -- (14.2,0);
\end{tikzpicture}
\legende{Chaque activité s'associe à des outils et à un lieu : mémorisez les lignes entières.}
\end{popfigure}

\courssub{Prononciation guidée}

Répétez après le professeur, en insistant sur les sons difficiles pour un francophone :

\begin{itemize}
\item \eng{wheelbarrow} : « ouil-ba-ro » — trois syllabes, le \eng{wh} se prononce « ou », le \eng{-ow} final se dit « o ».
\item \eng{shovel} : « che-veul » — le \eng{sh} se prononce « ch » comme dans « chemin ».
\item \eng{gutter} : « geu-teur » — le \eng{u} se prononce « eu » bref, jamais « ou ».
\item \eng{broom} : « broum » — le \eng{oo} se prononce « ou » long, comme dans « loup ».
\item \eng{rubbish} : « reu-bich » — accent sur la première syllabe.
\item \eng{fundraising} : « fonde-rei-zing » — le \eng{ai} se prononce « éi ».
\end{itemize}

\begin{savaistu}[Pourquoi « a working bee » ?]
L'expression \eng{a working bee} est typique de l'anglais américain et australien : les voisins se réunissent pour travailler ensemble, affairés comme des abeilles (\eng{bees}) dans une ruche. On trouve la même image dans \eng{a spelling bee}, un concours d'orthographe collectif très populaire aux États-Unis. Au Cameroun anglophone, on parle plus simplement de \eng{community work} ou de \eng{communal labour}.
\end{savaistu}

\begin{attention}[Le faux ami « corvée »]
Le mot français « corvée » (au sens de travail collectif du village) ne se traduit JAMAIS par « corvee » en anglais : ce mot n'existe pas dans l'anglais courant. Dites \eng{a work party}, \eng{a working bee} ou \eng{communal labour}. De même, « une campagne » au sens d'opération organisée se traduit bien par \eng{a campaign} (se prononce « kèm-péin » — le \eng{g} est muet !), mais « la campagne » au sens de la campagne rurale se dit \eng{the countryside}.
\end{attention}

\courssub{Exercice résolu et entraînement}

\begin{exoresolu}[Complétez avec le mot qui convient]
Énoncé — Complétez chaque phrase avec le mot anglais correct : \emph{broom, gutters, wheelbarrow, fundraising, water tap, harvest}.
\begin{enumerate}
\item The volunteers cleared the \ldots\ldots\ldots before the rains.
\item We need a \ldots\ldots\ldots to carry the sand to the building site.
\item The community organised a \ldots\ldots\ldots event to buy a new roof for the village hall.
\item She swept the classroom with a long \ldots\ldots\ldots .
\item The children fetch water at the public \ldots\ldots\ldots .
\item In December, the farmers \ldots\ldots\ldots their cocoa.
\end{enumerate}
\tcblower
Solution détaillée :
\begin{enumerate}
\item \eng{gutters} — on « débouche les caniveaux » : \eng{to clear the gutters}.
\item \eng{wheelbarrow} — la brouette sert à transporter le sable.
\item \eng{fundraising} — \eng{a fundraising event} = une collecte de fonds.
\item \eng{broom} — on balaie avec un balai : \eng{to sweep with a broom}.
\item \eng{water tap} — le robinet public où l'on puise l'eau.
\item \eng{harvest} — \eng{to harvest crops / cocoa} = récolter.
\end{enumerate}
\end{exoresolu}

\begin{exercice}[À vous de jouer]
Traduisez en anglais :
\begin{enumerate}
\item Les jeunes balaient la place du marché tous les samedis.
\item Nous avons planté des arbres devant le centre de santé.
\item Elles ont transporté l'eau dans des seaux jusqu'au chantier.
\item Le village a organisé une collecte de fonds pour réparer le pont.
\end{enumerate}
\end{exercice}

\begin{corrige}[Exercice « À vous de jouer »]
\begin{enumerate}
\item \eng{The youths sweep the market square every Saturday.}
\item \eng{We planted trees in front of the health centre.}
\item \eng{They carried water in buckets to the building site.}
\item \eng{The village organised a fundraising event to repair the bridge.}
\end{enumerate}
\end{corrige}

\begin{aretenir}[L'essentiel du champ lexical 1]
\begin{itemize}
\item Activités : \eng{a clean-up campaign, a work party / a working bee, a fundraising event, a building project, road maintenance, well digging}.
\item Actions : \eng{to sweep, to clear, to collect, to plant, to paint, to repair, to harvest, to dig, to carry}.
\item Outils : \eng{a broom, a machete, a wheelbarrow, a shovel, a rake, gloves, a bucket}.
\item Lieux : \eng{the market square, the health centre, the water tap, the village hall}.
\item Pièges : « corvée » $\neq$ « corvee » ; \eng{centre} (GB) $\neq$ « center » (US) ; \eng{equipment} est indénombrable.
\end{itemize}
\end{aretenir}

\coursec{Champ lexical 2 : participer, aider, s'organiser}

Dans la section précédente, nous avons découvert le vocabulaire des \emph{activités} communautaires : \eng{a clean-up exercise}, \eng{a fundraising event}, \eng{community work}. Mais une activité n'existe pas sans \emph{personnes} qui décident d'agir ensemble. Comment dit-on, en anglais, « participer », « donner un coup de main », « collecter des fonds » ? C'est précisément l'objet de cette section : le vocabulaire de la \cle{participation}, de l'\cle{aide} et de l'\cle{organisation}. Ces mots vous permettront de raconter, à l'oral comme à l'écrit, comment une communauté s'organise pour réussir un projet commun.

\courssub{Observer d'abord : un texte pour découvrir le lexique}

Lisons d'abord un court article de presse locale. Soulignez mentalement tous les mots qui expriment la participation ou l'organisation.

\begin{textebox}[The Buea Road Project — a local newspaper report]
Last Saturday, more than two hundred inhabitants of Buea \textbf{took part in} a communal work session on the old mountain road. The event was \textbf{organised} by the \textbf{quarter head}, Mr Efange, with the help of a \textbf{community leader}, Mrs Lyonga. She \textbf{mobilised} the youth and \textbf{coordinated} the different teams.

Many \textbf{volunteers} \textbf{joined in} the work early in the morning. Some students from Government High School \textbf{volunteered} to carry stones, while a group of traders \textbf{donated} food and drinks. The council \textbf{provided} wheelbarrows, shovels and drinking water. A local businessman even helped to \textbf{raise funds} to buy cement.

Everybody \textbf{contributed} something: the women cooked, the young men dug, and the children brought water. “When people \textbf{cooperate}, everything becomes easier,” said Mr Efange. “This project shows our \textbf{team spirit} and our \textbf{sense of belonging}.” The work was tiring but \textbf{rewarding}, and at the end of the day, the road was clean and safe again. It was a \textbf{successful} day of \textbf{collective} effort and \textbf{solidarity}.
\end{textebox}

\textbf{Glossaire}

\begin{center}
\begin{tabularx}{\linewidth}{|X|X|}\hline
\rowcolor{popblueL} \textbf{English} & \textbf{Français} \\ \hline
a wheelbarrow (n.) & une brouette \\ \hline
a shovel (n.) & une pelle \\ \hline
to dig (v.) & creuser \\ \hline
cement (n.) & le ciment \\ \hline
a trader (n.) & un commerçant \\ \hline
\end{tabularx}
\end{center}

\textbf{Questions de compréhension}
\begin{enumerate}
\item Who organised the communal work session?
\item What did the council provide?
\item What did the local businessman do?
\item Which two abstract qualities does Mr Efange mention at the end?
\end{enumerate}

\courssub{Les verbes de participation}

\begin{propriete}[Les verbes de participation et leurs prépositions]
\eng{To take part in}, \eng{to participate in} et \eng{to contribute to} sont \textbf{toujours suivis d'une préposition} (\eng{in} ou \eng{to}) avant le nom. On ne peut jamais dire \eng{to take part a meeting} ni \eng{to participate a project}.
\begin{itemize}
\item \eng{to take part in} + nom : participer à (le plus courant à l'oral).
\item \eng{to participate in} + nom : participer à (plus soutenu, plus formel).
\item \eng{to contribute to} + nom : contribuer à (attention : \eng{to}, pas \eng{in}).
\item \eng{to join in} + nom (ou sans complément) : se joindre à, participer activement.
\item \eng{to volunteer to} + infinitif / \eng{to volunteer for} + nom : se porter volontaire.
\item \eng{to help out} : donner un coup de main (familier, souvent sans complément).
\item \eng{to cooperate with} + personne : coopérer avec quelqu'un.
\end{itemize}
\end{propriete}

\begin{definition}[To lend a hand]
\eng{To lend a hand} (littéralement « prêter une main ») signifie \emph{to help someone, to give assistance} : prêter main-forte, donner un coup de main. C'est une expression idiomatique très fréquente.
\eng{Can you lend me a hand with these buckets?} « Peux-tu me donner un coup de main avec ces seaux ? »
\end{definition}

\begin{exemplebox}[Exemples traduits]
\eng{Many volunteers joined in the clean-up exercise.} « Beaucoup de volontaires se sont joints à l'opération de nettoyage. »

\eng{All the students took part in the tree-planting campaign.} « Tous les élèves ont participé à la campagne de reboisement. »

\eng{She volunteered to cook for the workers.} « Elle s'est portée volontaire pour cuisiner pour les travailleurs. »

\eng{Everyone contributed to the success of the project.} « Chacun a contribué au succès du projet. »

\eng{Our neighbours always help out when there is communal work.} « Nos voisins donnent toujours un coup de main quand il y a un travail communautaire. »

\eng{The two villages cooperated to build the bridge.} « Les deux villages ont coopéré pour construire le pont. »
\end{exemplebox}

\begin{attention}[Ne confondez pas les prépositions]
Les francophones traduisent souvent mot à mot et oublient la préposition anglaise. Retenez : on participe \emph{à} quelque chose, donc \eng{take part \textbf{in}}, \eng{participate \textbf{in}} ; mais on contribue \emph{à} avec \eng{contribute \textbf{to}}. Et attention : \eng{join} seul prend un complément direct (\eng{join a club} = adhérer à un club), tandis que \eng{join in} exprime la participation à une activité en cours.
\end{attention}

\courssub{Les verbes d'organisation}

Avant de participer, il faut organiser. Voici les verbes essentiels de l'organisation d'une action communautaire.

\begin{exemplebox}[Exemples traduits]
\eng{The quarter head organised a meeting to plan the work.} « Le chef de quartier a organisé une réunion pour planifier le travail. »

\eng{The community leader mobilised the youth of the village.} « Le chef de communauté a mobilisé les jeunes du village. »

\eng{Mrs Lyonga coordinated the different teams.} « Mme Lyonga a coordonné les différentes équipes. »

\eng{The council provided tools and drinking water.} « La mairie a fourni des outils et de l'eau potable. »

\eng{The traders donated food for the volunteers.} « Les commerçants ont fait don de nourriture pour les bénévoles. »

\eng{They raised funds to buy cement for the bridge.} « Ils ont collecté des fonds pour acheter du ciment pour le pont. »
\end{exemplebox}

Notez la construction de \eng{to provide} : \eng{provide + chose} ou \eng{provide somebody \textbf{with} something}. \eng{The council provided the workers with gloves.} « La mairie a fourni des gants aux travailleurs. »

\courssub{Les personnes, les noms abstraits et les adjectifs}

\begin{center}
\begin{tabularx}{\linewidth}{|X|X|X|}\hline
\rowcolor{popblueL} \textbf{English} & \textbf{Français} & \textbf{Exemple} \\ \hline
a volunteer & un volontaire, un bénévole & \eng{The volunteers arrived at 7 a.m.} \\ \hline
a participant & un participant & \eng{Each participant received a certificate.} \\ \hline
a community leader & un chef de communauté & \eng{The community leader gave a speech.} \\ \hline
a quarter head & un chef de quartier & \eng{The quarter head called a meeting.} \\ \hline
an organiser & un organisateur & \eng{The organisers thanked everybody.} \\ \hline
a donor & un donateur & \eng{An anonymous donor paid for the cement.} \\ \hline
\end{tabularx}
\end{center}

\begin{center}
\begin{tabularx}{\linewidth}{|X|X|}\hline
\rowcolor{popgreenL} \textbf{Noms abstraits} & \textbf{Français} \\ \hline
solidarity & la solidarité \\ \hline
cooperation & la coopération \\ \hline
team spirit & l'esprit d'équipe \\ \hline
mutual aid & l'entraide \\ \hline
commitment & l'engagement \\ \hline
a sense of belonging & un sentiment d'appartenance \\ \hline
\end{tabularx}
\end{center}

\begin{center}
\begin{tabularx}{\linewidth}{|X|X|}\hline
\rowcolor{popgoldL} \textbf{Adjectifs} & \textbf{Français} \\ \hline
voluntary & bénévole, volontaire \\ \hline
unpaid & non rémunéré \\ \hline
collective & collectif \\ \hline
joint & conjoint, commun \\ \hline
successful & réussi, couronné de succès \\ \hline
tiring but rewarding & fatigant mais gratifiant \\ \hline
\end{tabularx}
\end{center}

\eng{It was a joint effort by the whole village.} « C'était un effort conjoint de tout le village. »
\eng{Voluntary work is unpaid, but it is rewarding.} « Le travail bénévole n'est pas rémunéré, mais il est gratifiant. »

\begin{attention}[Prononciation : volunteer]
\eng{Volunteer} se prononce « vo-leun-\textbf{tir} » avec l'accent tonique sur la \textbf{dernière} syllabe. Ne dites pas « vo-lon-taire » comme en français ! De même, \eng{participant} se prononce « par-\textbf{ti}-si-peunt » et \eng{organiser} « \textbf{o}-gueu-naï-zeur ».
\end{attention}

\courssub{Schéma : les étapes d'un projet réussi}

\begin{popfigure}
\begin{center}
\begin{tikzpicture}[
  box/.style={draw=popblue, fill=popblueL, rounded corners, align=center, minimum width=2.6cm, minimum height=0.9cm, font=\small},
  main/.style={draw=poporange, fill=poporangeL, rounded corners, align=center, minimum width=4.2cm, minimum height=1cm, font=\bfseries\small},
  arr/.style={-{Stealth[length=3mm]}, thick, popdark}
]
\node[main] (p) at (0,0) {A successful project};
\node[box, align=center] (s1) at (-4.5,-1.8){plan\\ (planifier)};
\node[box, align=center] (s2) at (-1.5,-1.8){mobilise\\ (mobiliser)};
\node[box, align=center] (s3) at (1.5,-1.8){carry out\\ (réaliser)};
\node[box, align=center] (s4) at (4.5,-1.8){celebrate\\ (célébrer)};
\draw[arr] (p) -- (s1);
\draw[arr] (p) -- (s2);
\draw[arr] (p) -- (s3);
\draw[arr] (p) -- (s4);
\draw[arr, popgreen] (s1) -- (s2);
\draw[arr, popgreen] (s2) -- (s3);
\draw[arr, popgreen] (s3) -- (s4);
\end{tikzpicture}
\end{center}
\legende{Les quatre étapes d'un projet communautaire réussi : on planifie, on mobilise, on réalise, puis on célèbre.}
\end{popfigure}

Remarquez le verbe \eng{to carry out} : « réaliser, mener à bien ». \eng{The youth carried out the project in two days.} « Les jeunes ont réalisé le projet en deux jours. » Attention, \eng{to realise} signifie « se rendre compte », pas « réaliser » au sens de « accomplir » : c'est un faux ami que nous reverrons en détail dans la section 5.

\courssub{Exercice résolu et entraînement}

\begin{exoresolu}[Complétez avec le bon verbe]
Complétez les phrases avec : \eng{volunteered}, \eng{raised funds}, \eng{provided}, \eng{took part in}, \eng{lend a hand}.
\begin{enumerate}
\item The pupils \ldots\ the clean-up exercise last Friday.
\item The council \ldots\ shovels and gloves for the workers.
\item My uncle \ldots\ to drive the lorry full of stones.
\item The women's association \ldots\ to buy paint for the school.
\item Come and \ldots\ ! We need more people to carry water.
\end{enumerate}
\tcblower
\textbf{Solution détaillée}
\begin{enumerate}
\item \eng{took part in} — suivi de \eng{in} + nom (\eng{the clean-up exercise}) : « Les élèves ont participé à l'opération de nettoyage vendredi dernier. »
\item \eng{provided} — \eng{provide + chose} : « La mairie a fourni des pelles et des gants aux travailleurs. »
\item \eng{volunteered} — \eng{volunteer to} + infinitif (\eng{to drive}) : « Mon oncle s'est porté volontaire pour conduire le camion plein de pierres. »
\item \eng{raised funds} — \eng{raise funds to} + infinitif : « L'association des femmes a collecté des fonds pour acheter de la peinture pour l'école. »
\item \eng{lend a hand} — impératif, expression idiomatique : « Viens prêter main-forte ! Nous avons besoin de plus de monde pour porter l'eau. »
\end{enumerate}
\end{exoresolu}

\begin{exercice}[À vous de jouer]
Traduisez en anglais :
\begin{enumerate}
\item Beaucoup de jeunes se sont joints au travail communautaire samedi dernier.
\item Le chef de quartier a mobilisé tout le village.
\item Chacun a contribué au succès du projet.
\item Les bénévoles ont collecté des fonds pour réparer le pont.
\item Le travail était fatigant mais gratifiant.
\end{enumerate}
\end{exercice}

\begin{corrige}[Exercice « À vous de jouer »]
\begin{enumerate}
\item \eng{Many young people joined in the communal work last Saturday.}
\item \eng{The quarter head mobilised the whole village.}
\item \eng{Everybody contributed to the success of the project.}
\item \eng{The volunteers raised funds to repair the bridge.}
\item \eng{The work was tiring but rewarding.}
\end{enumerate}
Vérifiez bien les prépositions : \eng{joined \textbf{in}}, \eng{contributed \textbf{to}}. Et n'oubliez pas : \eng{the whole village} (« tout le village »), jamais \eng{all the village} dans ce sens.
\end{corrige}

\begin{aretenir}[L'essentiel de cette section]
\begin{itemize}
\item Participer : \eng{take part in}, \eng{participate in}, \eng{join in} — toujours avec la préposition.
\item Contribuer : \eng{contribute \textbf{to}} ; coopérer : \eng{cooperate \textbf{with}}.
\item Se porter volontaire : \eng{volunteer to} + infinitif ; donner un coup de main : \eng{help out}, \eng{lend a hand}.
\item Organiser : \eng{organise}, \eng{plan}, \eng{mobilise}, \eng{coordinate}, \eng{provide}, \eng{donate}, \eng{raise funds}.
\item Les personnes : \eng{a volunteer}, \eng{a participant}, \eng{a quarter head}, \eng{a donor}.
\item Les valeurs : \eng{solidarity}, \eng{team spirit}, \eng{mutual aid}, \eng{a sense of belonging}.
\item Prononciation : \eng{volunteer} = « vo-leun-\textbf{tir} », accent sur la dernière syllabe.
\end{itemize}
\end{aretenir}

Vous savez maintenant dire qui fait quoi dans une action communautaire. Dans la prochaine section, nous assemblerons ces mots en \emph{collocations} et \emph{expressions idiomatiques} pour parler comme un vrai anglophone.

\coursec{Collocations et expressions idiomatiques}

Dans la section précédente, nous avons appris les verbes essentiels pour parler de la participation : \eng{to take part in}, \eng{to help}, \eng{to volunteer}, \eng{to organise}. Mais un mot isolé ne suffit pas pour parler un anglais naturel. En anglais comme en français, les mots ont des « compagnons habituels » : on dit en français \eng{faire sa part} et non \eng{réaliser sa part}. Ces associations de mots s'appellent des \cle{collocations}, et les expressions figées comme \eng{donner un coup de main} sont des \cle{expressions idiomatiques}. C'est exactement ce que nous allons étudier maintenant, à travers le thème du travail communautaire.

\courssub{Observation : un texte pour découvrir les collocations}

\begin{textebox}[A successful clean-up campaign in Bamenda]
Last month, the youth association of Nkwen quarter decided to \textbf{carry out} an ambitious project: cleaning the stream that crosses the neighbourhood. The leaders \textbf{launched} a campaign on the radio and in churches to invite everybody. “\textbf{Let's join hands} and clean our quarter!” they announced. “Everyone is expected to \textbf{do their bit}.”

On Saturday morning, more than two hundred volunteers \textbf{rolled up their sleeves} and got to work. Some residents came to \textbf{lend a hand} with spades and wheelbarrows; others \textbf{raised funds} to buy gloves, boots and rubbish bags. A local businessman \textbf{made a generous contribution}: he paid for a lorry to carry the waste away. The \textbf{teamwork} was remarkable. Thanks to this \textbf{hard work}, the stream was cleared in a single weekend. \textbf{Many hands make light work}, as the saying goes!

The association now wants to \textbf{achieve} a new \textbf{goal}: building a small bridge. They know that unpaid \textbf{voluntary work} can transform a community when people \textbf{pull together}.
\end{textebox}

\textbf{Glossaire :}

\begin{center}
\begin{tabularx}{\linewidth}{|X|X|}\hline
\rowcolor{popblueL} English & Français \\ \hline
to carry out a project & mener à bien un projet \\ \hline
to launch a campaign & lancer une campagne \\ \hline
to do one's bit & faire sa part \\ \hline
to roll up one's sleeves & retrousser ses manches, se mettre au travail \\ \hline
to lend a hand & donner un coup de main \\ \hline
to raise funds & collecter des fonds \\ \hline
to make a contribution & apporter une contribution \\ \hline
to pull together & se serrer les coudes \\ \hline
\end{tabularx}
\end{center}

\textbf{Questions de compréhension :}
\begin{enumerate}
\item What project did the youth association carry out?
\item How did the businessman make his contribution?
\item Which proverb in the text means that work is easier when many people help?
\end{enumerate}

\textbf{Réponses attendues :} 1. \eng{They carried out a project to clean the stream that crosses the neighbourhood.} 2. \eng{He paid for a lorry to carry the waste away.} 3. \eng{Many hands make light work.}

\courssub{Les collocations avec « work »}

\begin{definition}[Qu'est-ce qu'une collocation ?]
Une \cle{collocation} est une association habituelle et naturelle de deux mots ou plus : \eng{hard work} (travail dur), \eng{make a contribution} (apporter une contribution). On ne peut pas remplacer un des mots par un synonyme sans produire une phrase bizarre ou incorrecte : on dit \eng{hard work}, jamais \eng{solid work}.
\end{definition}

Le mot \eng{work} est au cœur du thème du travail communautaire. Voici ses collocations principales :

\begin{center}
\begin{tabularx}{\linewidth}{|X|X|X|}\hline
\rowcolor{popblueL} Collocation & Traduction & Exemple \\ \hline
\eng{communal work} & travail communautaire & \eng{Communal work unites the village.} \\ \hline
\eng{voluntary work} & travail bénévole & \eng{She does voluntary work at the hospital.} \\ \hline
\eng{teamwork} & travail d'équipe & \eng{Teamwork made the project a success.} \\ \hline
\eng{hard work} & travail dur, gros efforts & \eng{The bridge required hard work.} \\ \hline
\eng{unpaid work} & travail non rémunéré & \eng{Community service is unpaid work.} \\ \hline
\eng{to carry out work} & effectuer des travaux & \eng{They carried out repair work on the road.} \\ \hline
\end{tabularx}
\end{center}

\begin{exemplebox}[Exemples traduits]
\eng{Many hands make light work: the bridge was repaired in two days.} « À plusieurs, le travail est léger : le pont a été réparé en deux jours. »

\eng{Everyone rolled up their sleeves and did their bit.} « Chacun s'est mis au travail et a fait sa part. »

\eng{Voluntary work is unpaid, but it is very rewarding.} « Le travail bénévole n'est pas rémunéré, mais il est très gratifiant. »
\end{exemplebox}

\begin{attention}[« Work » est indénombrable]
En anglais, \eng{work} est un nom \cle{indénombrable} : on ne dit jamais \eng{a work} ni \eng{works} pour parler de travail en général. On dit \eng{a lot of work} (beaucoup de travail), \eng{a piece of work} (un travail, une tâche). Attention aussi : \eng{works} au pluriel signifie « les œuvres » (d'un écrivain) ou « les travaux publics » dans \eng{road works}. Ne confondez pas non plus \eng{work} (le travail) et \eng{job} (un emploi, dénombrable) : \eng{She is looking for a job}, mais \eng{She has a lot of work to do}.
\end{attention}

\courssub{Collocations verbe + nom : les paires à connaître}

Pour agir dans une communauté, il faut maîtriser les couples verbe + nom. Chaque verbe a son partenaire privilégié :

\begin{center}
\begin{tabularx}{\linewidth}{|X|X|X|}\hline
\rowcolor{popblueL} Collocation & Traduction & Exemple en contexte \\ \hline
\eng{to carry out a project} & mener un projet & \eng{The students carried out a project to plant trees.} \\ \hline
\eng{to launch a campaign} & lancer une campagne & \eng{The mayor launched a campaign against litter.} \\ \hline
\eng{to raise money / funds} & collecter des fonds & \eng{We raised money to buy books for the school.} \\ \hline
\eng{to make a contribution} & apporter une contribution & \eng{Every family made a contribution of 500 francs.} \\ \hline
\eng{to do community service} & faire des travaux d'intérêt général & \eng{He does community service at the health centre.} \\ \hline
\eng{to achieve a goal} & atteindre un objectif & \eng{Together, we achieved our goal in one week.} \\ \hline
\end{tabularx}
\end{center}

\begin{propriete}[Règle : le verbe est fixé par l'usage]
Dans une collocation verbe + nom, le choix du verbe ne se traduit pas mot à mot depuis le français. Le français « collecter des fonds » devient \eng{to raise funds} (et non \eng{to collect funds}, qui existe mais est moins courant dans ce sens) ; « atteindre un objectif » devient \eng{to achieve a goal} ou \eng{to reach a goal}. Il faut donc mémoriser chaque paire comme un tout.
\end{propriete}

\begin{attention}[Ne pas inverser les verbes]
On dit \eng{to \textbf{do} one's bit} et \eng{to \textbf{make} a contribution}. Ne mélangez pas les deux : \eng{make one's bit} et \eng{do a contribution} sont des erreurs typiques des francophones. Astuce : \eng{do} s'emploie pour les actions et les efforts (\eng{do one's bit}, \eng{do community service}, \eng{do voluntary work}) ; \eng{make} s'emploie pour ce qu'on produit ou apporte (\eng{make a contribution}, \eng{make an effort}, \eng{make a donation}).
\end{attention}

\begin{methode}[Comment retenir les collocations]
\begin{enumerate}
\item N'apprenez jamais un mot isolément : apprenez-le \textbf{par paire} (\eng{launch + campaign}, \eng{raise + funds}).
\item Classez les collocations par verbe dans votre cahier : une page pour \eng{do}, une pour \eng{make}, une pour \eng{carry out}.
\item Réemployez chaque paire dans une \textbf{phrase personnelle} liée à votre vie : \eng{Our class raised funds for the football tournament.}
\item Relisez vos phrases une semaine plus tard et testez-vous en cachant la colonne anglaise.
\end{enumerate}
\end{methode}

\courssub{Expressions idiomatiques de l'entraide}

\begin{definition}[Expression idiomatique]
Une \cle{expression idiomatique} est un groupe de mots figé dont le sens global ne se devine pas à partir des mots pris un à un. \eng{To lend a hand} ne veut pas dire « prêter une main » mais « donner un coup de main ».
\end{definition}

\begin{center}
\begin{tabularx}{\linewidth}{|X|X|X|}\hline
\rowcolor{popblueL} Expression & Équivalent français & Exemple \\ \hline
\eng{to lend a hand} & donner un coup de main & \eng{Come and lend a hand this Saturday!} \\ \hline
\eng{many hands make light work} & à plusieurs, le travail est léger & \eng{Fifty volunteers came: many hands make light work.} \\ \hline
\eng{to pull together} & se serrer les coudes & \eng{The villagers pulled together after the storm.} \\ \hline
\eng{to do one's bit / one's share} & faire sa part & \eng{Everybody is expected to do their bit.} \\ \hline
\eng{to roll up one's sleeves} & retrousser ses manches, se mettre au travail & \eng{We rolled up our sleeves and cleaned the market.} \\ \hline
\end{tabularx}
\end{center}

\begin{exemplebox}[Expressions pour inviter]
\eng{Let's join hands and clean our quarter!} « Unissons nos efforts et nettoyons notre quartier ! »

\eng{Everybody is expected to do their bit.} « Chacun est invité à faire sa part. »

\eng{Come and lend a hand this Saturday!} « Viens donner un coup de main ce samedi ! »
\end{exemplebox}

\begin{attention}[Ne traduisez pas mot à mot]
Les expressions idiomatiques ne se traduisent jamais littéralement. \eng{To roll up one's sleeves} ne parle pas de manches : il faut chercher l'équivalent français (« se mettre au travail »). De même, \eng{many hands make light work} est un proverbe figé : on ne change ni l'ordre ni les mots (jamais \eng{many hands make easy work}). Notez aussi que \eng{lend} est irrégulier : \eng{lend -- lent -- lent} ; ne confondez pas \eng{to lend} (prêter) et \eng{to borrow} (emprunter).
\end{attention}

\begin{popfigure}
\begin{tikzpicture}[mindmap, grow cyclic, every node/.style={concept, align=center}, concept color=popblue, root concept/.append style={concept color=poporange, font=\bfseries}, level 1/.append style={level distance=3.4cm, sibling angle=90}]
\node[root concept, align=center]{LEND A HAND\\ donner un coup de main}
  child[concept color=popgreen] { node[align=center]{pull together\\ se serrer\\ les coudes} }
  child[concept color=poppurple] { node[align=center]{do one's bit\\ faire sa part} }
  child[concept color=poppink] { node[align=center]{roll up one's\\ sleeves\\ se mettre au travail} }
  child[concept color=popgold] { node[align=center]{many hands\\ make light work\\ à plusieurs,\\ tout est léger} };
\end{tikzpicture}
\legende{Carte mentale : les expressions idiomatiques de l'entraide autour de \eng{lend a hand}.}
\end{popfigure}

\courssub{Prononciation et orthographe : les pièges à éviter}

\begin{propriete}[Prononciation des mots clés]
\begin{itemize}
\item \eng{voluntary} se prononce « vo-len-teu-ri » (accent sur la première syllabe) ; \eng{volunteer} se prononce « vo-len-tière » (accent sur la dernière syllabe).
\item \eng{contribution} se prononce « kon-tri-biou-cheune » ; le verbe \eng{contribute} se prononce « keun-tri-bioute ».
\item \eng{campaign} se prononce « kam-péïne » : le \eng{-gn} est muet, comme dans \eng{sign} et \eng{foreign}.
\item \eng{neighbourhood} se prononce « né-beur-houde » : le \eng{-gh-} est muet. Orthographe britannique avec \eng{-our} ; l'orthographe américaine est \eng{neighborhood}.
\end{itemize}
\end{propriete}

\begin{attention}[Faux amis et calques du français]
\begin{itemize}
\item \eng{To assist} signifie « aider » (soutenir), pas « assister à » au sens d'« être présent » : « assister à une réunion » se dit \eng{to attend a meeting}.
\item « Participer à » se dit \eng{to take part in} ou \eng{to participate in} : n'oubliez jamais la préposition \eng{in}.
\item « Aider quelqu'un à faire » se dit \eng{to help somebody (to) do} : \eng{She helped me carry the bags.} — jamais de calque avec \eng{for} ou \eng{at}.
\item \eng{Actually} signifie « en fait, en réalité », pas « actuellement » : « actuellement » se dit \eng{currently} ou \eng{at present}.
\end{itemize}
\end{attention}

\courssub{Mini-dialogue : inviter, accepter, refuser poliment}

\begin{textebox}[Inviting a friend to communal work]
\textbf{Ngong:} Hi, Manka! Our quarter is organising a clean-up day this Saturday. Come and lend a hand!

\textbf{Manka:} That sounds great. What time does it start?

\textbf{Ngong:} At seven o'clock, in front of the chief's palace. Everybody is expected to do their bit.

\textbf{Manka:} Count me in! I'll bring my spade and a wheelbarrow.

\textbf{Ngong:} Wonderful. What about your brother? Can he join us?

\textbf{Manka:} I'm afraid he can't. He has an exam on Monday, so he has to revise. But he promises to make a contribution: he will buy gloves for the team.

\textbf{Ngong:} That's very kind of him. You see, many hands make light work — and every contribution counts!
\end{textebox}

\textbf{Formules utiles à retenir :}
\begin{itemize}
\item Inviter : \eng{Come and lend a hand!} / \eng{Can you join us?} / \eng{Let's join hands!}
\item Accepter : \eng{Count me in!} (je suis de la partie !) / \eng{That sounds great.} / \eng{I'll be there.}
\item Refuser poliment : \eng{I'm afraid I can't.} (je crains de ne pas pouvoir) / \eng{I'd love to, but I have to revise.}
\end{itemize}

\begin{exoresolu}[Compléter avec la bonne collocation]
Énoncé. Complétez chaque phrase avec la collocation correcte : \eng{raise funds}, \eng{do my bit}, \eng{make a contribution}, \eng{lend a hand}, \eng{achieve our goal}, \eng{launched a campaign}.
\begin{enumerate}
\item The association \ldots\ to clean the river last month.
\item We need to \ldots\ to buy new wheelbarrows.
\item Everyone must \ldots\ ; nobody should stay at home.
\item If we all pull together, we shall \ldots\ before the rainy season.
\item My uncle cannot come, but he will \ldots\ of 10,000 francs.
\item Come and \ldots\ this Saturday at the market square!
\end{enumerate}
\tcblower
Solution détaillée.
\begin{enumerate}
\item \eng{launched a campaign} — c'est le verbe \eng{launch} qui accompagne \eng{campaign} (lancer une campagne).
\item \eng{raise funds} — pour l'argent, la collocation est \eng{raise funds/money} (collecter des fonds).
\item \eng{do my bit} — attention : \eng{do} et non \eng{make} avec \eng{one's bit} (faire sa part).
\item \eng{achieve our goal} — \eng{achieve} accompagne \eng{goal} (atteindre un objectif).
\item \eng{make a contribution} — ici c'est \eng{make} : on « produit » une contribution (apporter une contribution).
\item \eng{lend a hand} — expression idiomatique pour inviter quelqu'un (donner un coup de main).
\end{enumerate}
\end{exoresolu}

\begin{savaistu}[« Community service » : bénévolat ou peine de tribunal ?]
En Grande-Bretagne et aux États-Unis, \eng{community service} a deux sens. Il désigne d'abord le travail bénévole au service de la communauté. Mais il désigne aussi une \textbf{peine de travaux d'intérêt général} prononcée par un tribunal : au lieu d'aller en prison, une personne condamnée pour un délit mineur peut devoir nettoyer des parcs ou aider dans des centres sociaux pendant un certain nombre d'heures. Le contexte permet de savoir de quel sens il s'agit !
\end{savaistu}

\begin{exercice}[À vous de jouer]
\begin{enumerate}
\item Traduisez en anglais : « Tout le monde a retroussé ses manches et le marché a été nettoyé en une matinée. »
\item Trouvez l'erreur et corrigez : \eng{The students made their bit and did a contribution.}
\item Rédigez trois phrases pour inviter vos camarades à un travail communautaire, en utilisant \eng{lend a hand}, \eng{do your bit} et \eng{pull together}.
\end{enumerate}
\end{exercice}

\begin{corrige}[Exercice « À vous de jouer »]
\begin{enumerate}
\item \eng{Everyone rolled up their sleeves and the market was cleaned in one morning.} (expression figée : \eng{roll up one's sleeves} ; voix passive naturelle en anglais.)
\item On inverse les verbes : \eng{The students \textbf{did} their bit and \textbf{made} a contribution.}
\item Exemples de réponses : \eng{Come and lend a hand this Saturday!} / \eng{Everybody is expected to do their bit for our school.} / \eng{If we pull together, we will finish before noon.}
\end{enumerate}
\end{corrige}

\begin{aretenir}[L'essentiel de la section]
\begin{itemize}
\item Une \cle{collocation} est une paire de mots fixée par l'usage : \eng{hard work}, \eng{raise funds}, \eng{achieve a goal}.
\item \eng{Work} est indénombrable : jamais \eng{a work} pour « un travail » ; un emploi se dit \eng{a job}.
\item \eng{Do} pour les efforts (\eng{do one's bit}, \eng{do community service}) ; \eng{make} pour ce qu'on apporte (\eng{make a contribution}).
\item Les expressions idiomatiques (\eng{lend a hand}, \eng{pull together}, \eng{roll up one's sleeves}) se mémorisent en bloc et ne se traduisent jamais mot à mot.
\item Pour inviter : \eng{Come and lend a hand!} ; pour refuser poliment : \eng{I'm afraid I can't.}
\item Faux amis : \eng{attend a meeting} (assister à), \eng{take part in} (participer à), \eng{actually} (en fait).
\end{itemize}
\end{aretenir}

\coursec{Faux amis et pièges des francophones}

Dans la section précédente, nous avons appris à combiner correctement les mots du travail communautaire grâce aux collocations (\eng{to take part in}, \eng{to lend a hand}, \eng{to pull together}). Mais il reste un danger plus sourd pour le francophone : les mots anglais qui \emph{ressemblent} à des mots français sans avoir le même sens. Ce sont les \cle{faux amis} (\eng{false friends} ou \eng{false cognates}). Le thème du travail communautaire en regorge : \eng{assist}, \eng{communal}, \eng{volunteer}, \eng{demand}, \eng{event}… Autant de pièges que cette section va démonter un à un.

\courssub{Observer avant de conclure : un texte piège}

Lisons d'abord ce court article de presse locale. Il contient plusieurs mots « dangereux » pour un francophone. Soulignez-les mentalement avant de lire le glossaire.

\begin{textebox}[The Buea Mountain Clean-Up — a report]
Last Saturday, more than two hundred residents of Buea took part in a communal clean-up campaign around the market and the motor park. The event was organised by the neighbourhood committee, with the support of the council.

“We did not demand anything from the population,” explains Mrs Eposi, the committee leader. “We simply asked people to give two hours of their time, and the response was wonderful.”

Among the volunteers was Mr Tabi, a retired teacher. “I attend every community meeting,” he says, “but this time I wanted to do some real work with my hands.” He helped carry sacks of rubbish to the collection lorry, assisted by two strong-willed students from the Government High School.

The council eventually provided wheelbarrows, gloves and drinking water. A possible second campaign is already being discussed for next month. “If everyone lends a hand,” concludes Mrs Eposi, “our town will be cleaner and healthier for all of us.”
\end{textebox}

\textbf{Glossaire}

\begin{center}
\begin{tabularx}{\linewidth}{|l|l|X|}
\hline
\rowcolor{popblueL} Mot anglais & Nature & Traduction française \\
\hline
communal (clean-up) & adjectif & collectif, communautaire \\
\hline
event & nom & événement, manifestation \\
\hline
to demand & verbe & exiger \\
\hline
to ask & verbe & demander \\
\hline
to attend (a meeting) & verbe & assister à, être présent à \\
\hline
volunteer & nom & bénévole, volontaire \\
\hline
strong-willed & adjectif & volontaire, déterminé \\
\hline
eventually & adverbe & finalement, en fin de compte \\
\hline
possible & adjectif & éventuel, possible \\
\hline
\end{tabularx}
\end{center}

\textbf{Questions de compréhension}

\begin{enumerate}
\item Who organised the clean-up campaign?
\item Did the committee demand anything from the population?
\item What does Mr Tabi usually attend?
\item What did the council eventually provide?
\item Is a second campaign certain or only possible?
\end{enumerate}

\courssub{Faux ami 1 : \eng{to assist} n'est pas « assister à »}

\begin{definition}[To assist / to attend]
\eng{To assist} signifie \cle{aider} quelqu'un (\eng{to assist somebody}), souvent dans un registre soutenu. En revanche, « assister à » au sens d'\emph{être présent à} se dit \eng{to attend} : \eng{to attend a meeting}, \eng{to attend a ceremony}.
\end{definition}

\eng{Two students assisted Mr Tabi.} « Deux élèves ont aidé M. Tabi. »

\eng{I attended the village meeting.} « J'ai assisté à la réunion du village. »

\begin{attention}[L'erreur classique]
La phrase \eng{I assisted to the clean-up} est \textbf{doublement incorrecte} : \eng{assist} ne signifie pas « assister à », et il est de toute façon suivi d'un complément direct, sans \eng{to}. Pour dire « j'ai participé au nettoyage », dites \eng{I took part in the clean-up} ou \eng{I joined the clean-up}. Pour dire « j'ai assisté à la réunion », dites \eng{I attended the meeting}.
\end{attention}

\courssub{Faux ami 2 : \eng{communal} — sens et prononciation}

\begin{definition}[Communal]
\eng{Communal} signifie \cle{collectif}, \cle{communautaire}, « partagé par tous » : \eng{communal work}, \eng{a communal water tap}, \eng{communal land}. Il se prononce « \textbf{ko-myou-neul} » (accent sur la deuxième syllabe \emph{myou}), et non « ko-mu-nal » comme en français.
\end{definition}

Attention également au sens : le français « communal » renvoie souvent à la \emph{commune} (administration), tandis que l'anglais \eng{communal} insiste sur le \emph{partage} et la \emph{collectivité}.

\eng{The villagers dug a communal well.} « Les villageois ont creusé un puits collectif. »

\eng{We share a communal kitchen in the students' residence.} « Nous partageons une cuisine commune à la résidence des étudiants. »

\courssub{Faux ami 3 : \eng{a volunteer} n'est pas « volontaire » (adjectif)}

\begin{definition}[Volunteer / strong-willed]
\eng{A volunteer} est un \cle{bénévole}, une personne qui travaille gratuitement. L'adjectif français « volontaire » au sens de « qui a de la volonté, déterminé » se traduit par \eng{strong-willed} ou \eng{determined}.
\end{definition}

\eng{Fifty volunteers cleaned the hospital yard.} « Cinquante bénévoles ont nettoyé la cour de l'hôpital. »

\eng{She is a strong-willed girl; she never gives up.} « C'est une fille volontaire ; elle n'abandonne jamais. »

\begin{exemplebox}[Exemples traduits]
\eng{She volunteered to help the old people.} « Elle s'est portée volontaire pour aider les personnes âgées. »

\eng{They asked the council for wheelbarrows.} « Ils ont demandé des brouettes à la mairie. »
\end{exemplebox}

\courssub{Faux ami 4 : \eng{to demand} est plus fort que « demander »}

\begin{definition}[To demand / to ask (for)]
\eng{To demand} signifie \cle{exiger}, réclamer avec force. Le simple « demander » se dit \eng{to ask} ; « demander quelque chose » se dit \eng{to ask for something} (préposition \eng{for} obligatoire devant la chose demandée).
\end{definition}

\eng{They asked for tools and gloves.} « Ils ont demandé des outils et des gants. » (demande normale)

\eng{The workers demanded better equipment.} « Les ouvriers ont exigé un meilleur équipement. » (ton plus fort, presque un ultimatum)

\eng{She asked me a question.} « Elle m'a posé une question. » (personne interrogée : pas de \eng{for})

\courssub{Faux ami 5 : \eng{event}, « éventuel » et \eng{eventually}}

\begin{definition}[Event / possible / eventually]
\eng{An event} est un \cle{événement}, une manifestation. « Éventuel » (adjectif) se dit \eng{possible} ou \eng{potential} ; « éventuellement » (adverbe) se dit \eng{possibly} ou \eng{perhaps}. Quant à \eng{eventually}, il signifie \cle{finalement}, « en fin de compte » — c'est l'inverse du piège !
\end{definition}

\eng{The clean-up was a great event.} « Le nettoyage a été un grand événement. »

\eng{A possible second campaign will take place in June.} « Une éventuelle seconde campagne aura lieu en juin. »

\eng{It rained all morning, but the work eventually started at noon.} « Il a plu toute la matinée, mais le travail a finalement commencé à midi. »

\courssub{Calques et constructions à corriger}

Au-delà des mots isolés, le francophone calque souvent des structures entières. Voici les trois plus fréquentes dans ce thème.

\begin{propriete}[Les bons verbes et les bonnes prépositions]
\begin{itemize}
\item On ne dit pas \eng{to make a work} mais \eng{to do some work} (\eng{work} est indénombrable : jamais \eng{a work} au sens de « travail »).
\item \eng{To participate} se construit avec \eng{in} : \eng{to participate in a project}, jamais \eng{to participate to}.
\item \eng{To aid} est un verbe soutenu ou journalistique ; dans la langue courante, on dit \eng{to help someone}.
\end{itemize}
\end{propriete}

\begin{propriete}[La construction de \eng{help}]
\eng{Help} est suivi d'un complément de personne puis de l'infinitif \textbf{avec ou sans} \eng{to} : \eng{Help me (to) carry this bucket.} « Aide-moi à porter ce seau. » Les deux formes sont correctes ; la forme sans \eng{to} est plus fréquente à l'oral.
\end{propriete}

\eng{They helped the women (to) build a new market shed.} « Ils ont aidé les femmes à construire un nouveau hangar au marché. »

\begin{attention}[Résumé des pièges]
Retenez : \eng{assist} = aider (pas « assister à ») ; \eng{demand} = exiger ; \eng{eventually} = finalement ; \eng{participate in} ; \eng{do some work}. Un mot anglais qui ressemble trop à un mot français doit toujours éveiller votre méfiance.
\end{attention}

\courssub{Tableaux récapitulatifs}

\begin{center}
\begin{tabularx}{\linewidth}{|X|X|X|}
\hline
\rowcolor{popblueL} Français & English & Remarque \\
\hline
assister à (une réunion) & to attend (a meeting) & \eng{assist} = aider \\
\hline
aider quelqu'un & to help / to assist somebody & \eng{help} au quotidien \\
\hline
demander quelque chose & to ask for something & préposition \eng{for} \\
\hline
exiger & to demand & sens fort \\
\hline
un bénévole & a volunteer & « volontaire » (adj.) = \eng{strong-willed} \\
\hline
éventuel(le) & possible / potential & \eng{eventual} n'existe pas dans ce sens \\
\hline
éventuellement & possibly / perhaps & \eng{eventually} = finalement \\
\hline
participer à & to participate in / to take part in & jamais \eng{participate to} \\
\hline
faire du travail & to do some work & jamais \eng{make a work} \\
\hline
\end{tabularx}
\end{center}

\begin{popfigure}
\begin{tikzpicture}
\node[draw=popdark, fill=popblueL, rounded corners, font=\bfseries] (titre) at (0,0) {Pièges du francophone};
\node[draw=poporange, fill=poporangeL, rounded corners, align=center] (a) at (-4.5,-1.8) {assister à\\ \eng{to attend}};
\node[draw=poporange, fill=poporangeL, rounded corners, align=center] (b) at (0,-1.8) {demander\\ \eng{to ask for}};
\node[draw=poporange, fill=poporangeL, rounded corners, align=center] (c) at (4.5,-1.8) {éventuellement\\ \eng{possibly}};
\node[draw=popgreen, fill=popgreenL, rounded corners, align=center] (d) at (-4.5,-3.6) {aider\\ \eng{to assist / to help}};
\node[draw=popgreen, fill=popgreenL, rounded corners, align=center] (e) at (0,-3.6) {exiger\\ \eng{to demand}};
\node[draw=popgreen, fill=popgreenL, rounded corners, align=center] (f) at (4.5,-3.6) {finalement\\ \eng{eventually}};
\draw[-{Stealth}, thick, popmuted] (titre) -- (a);
\draw[-{Stealth}, thick, popmuted] (titre) -- (b);
\draw[-{Stealth}, thick, popmuted] (titre) -- (c);
\draw[-{Stealth}, thick, popmuted] (titre) -- (d);
\draw[-{Stealth}, thick, popmuted] (titre) -- (e);
\draw[-{Stealth}, thick, popmuted] (titre) -- (f);
\end{tikzpicture}
\legende{French trap / Correct English : les paires à ne jamais confondre.}
\end{popfigure}

\courssub{Exercice résolu}

\begin{exoresolu}[Corrige les erreurs]
Énoncé — Chaque phrase contient un faux ami ou un calque du français. Corrige-la.

1. \eng{I assisted to the communal meeting yesterday.}

2. \eng{The villagers demanded for clean water.}

3. \eng{We will eventually organise a clean-up, if the mayor agrees.} (sens voulu : « éventuellement »)

4. \eng{She is a volunteer girl who never abandons her projects.} (sens voulu : « volontaire »)

5. \eng{They participated to the work and made a good work.}

\tcblower

Solution détaillée —

1. \eng{I attended the communal meeting yesterday.} « Assister à » (être présent) = \eng{attend} ; \eng{assist} signifie « aider » et ne prend jamais la préposition \eng{to}.

2. \eng{The villagers demanded clean water.} \eng{Demand} (= exiger) se construit avec un complément direct, sans \eng{for}. Si le sens est simplement « demander », on dira : \eng{The villagers asked for clean water.}

3. \eng{We will possibly organise a clean-up, if the mayor agrees.} « Éventuellement » = \eng{possibly} ; \eng{eventually} signifie « finalement » et changerait le sens de la phrase.

4. \eng{She is a strong-willed girl who never abandons her projects.} \eng{Volunteer} désigne un bénévole ; l'adjectif « volontaire » (déterminé) = \eng{strong-willed}.

5. \eng{They participated in the work and did a good job.} \eng{Participate} exige \eng{in} ; on dit \eng{to do work}, et « faire du bon travail » se rend naturellement par \eng{to do a good job}.
\end{exoresolu}

\begin{methode}[Comment éviter les faux amis]
\begin{enumerate}
\item \textbf{Repère} tout mot anglais qui ressemble fortement à un mot français (\eng{assist}, \eng{demand}, \eng{eventual}, \eng{communal}…).
\item \textbf{Vérifie} son sens dans un dictionnaire anglais-anglais (ou bilingue fiable) avant de l'utiliser : la ressemblance est un signal d'alarme, pas une garantie.
\item \textbf{Mémorise} le mot avec sa construction complète : \eng{attend a meeting}, \eng{ask for tools}, \eng{participate in a project} — jamais le mot seul.
\item \textbf{Crée} ta propre phrase en contexte camerounais (\eng{We attended the clean-up in Bamenda.}) pour ancrer le bon usage.
\item \textbf{Relis} tes productions écrites en traquant spécifiquement ces cinq faux amis.
\end{enumerate}
\end{methode}

\begin{exercice}[À toi de jouer]
Traduis en anglais correct :

1. J'ai assisté à la réunion du comité de quartier.

2. Les bénévoles ont demandé des brouettes à la mairie.

3. Le nettoyage a finalement commencé à huit heures.

4. Une éventuelle seconde campagne est prévue en décembre.

5. Aide-moi à porter ce sac de sable, s'il te plaît.
\end{exercice}

\begin{corrige}[Exercice — À toi de jouer]
1. \eng{I attended the neighbourhood committee meeting.}

2. \eng{The volunteers asked the council for wheelbarrows.}

3. \eng{The clean-up eventually started at eight o'clock.}

4. \eng{A possible second campaign is planned for December.}

5. \eng{Help me (to) carry this bag of sand, please.}
\end{corrige}

\begin{savaistu}[Pourquoi tant de faux amis ?]
Après la conquête normande de l'Angleterre en 1066, des milliers de mots français sont entrés dans la langue anglaise. Mais au fil des siècles, les deux langues ont fait évoluer ces mots différemment : le français « demander » est resté poli, tandis que l'anglais \eng{demand} a durci son sens ; « assister à » a gardé l'idée de présence, alors que \eng{assist} a gardé celle d'aide. Les faux amis sont donc des cousins historiques… qui ont pris des chemins opposés.
\end{savaistu}

\coursec{Familles de mots, prononciation et orthographe}

Dans la section précédente, nous avons démasqué les faux amis et les calques du français qui guettent le thème du travail communautaire. Nous allons maintenant franchir une étape décisive : au lieu d'apprendre les mots un par un, nous allons les apprendre \cle{par familles}. En anglais comme en français, un même mot de base se transforme en verbe, en nom de personne, en nom d'action ou en adjectif grâce à des \cle{suffixes}. Comprendre ce mécanisme, c'est multiplier votre vocabulaire par quatre ou cinq sans effort supplémentaire. Nous terminerons par la prononciation et l'orthographe, deux terrains minés pour les francophones.

\courssub{Observer avant de mémoriser : un texte témoin}

\begin{textebox}[The Njangi Group Cleans Up — reading text]
Last Saturday, the \textbf{volunteers} of the Hope Njangi Group met at six o'clock in the morning in the \textbf{neighbourhood} of Mokolo, in Yaoundé. Their \textbf{participation} was impressive: more than forty \textbf{participants}, young and old, came with brooms, wheelbarrows and buckets. The \textbf{organizers} had \textbf{organized} everything perfectly the day before, with the \textbf{cooperation} of the local council.

Mrs Etoundi, a regular \textbf{contributor}, made a generous \textbf{contribution}: she offered bread and tea to everyone. A local businessman, who prefers to remain an anonymous \textbf{donor}, had sent a \textbf{donation} of fifty pairs of gloves. ``We all \textbf{cooperate} because we love our community,'' said Paul, the youngest \textbf{volunteer}. ``I \textbf{volunteer} every month. It is not an obligation; I do it \textbf{voluntarily}, and I am proud of it.''

By noon, the gutters were clean, the rubbish had been collected, and the market square looked brand new. The \textbf{organizers} thanked all the \textbf{participants} warmly. ``Without your \textbf{voluntary} work,'' said the president, ``our neighbourhood would not be so clean. \textbf{Solidarity} is our strength.'' Everyone agreed to \textbf{participate} again the following month, and a \textbf{cooperative} spirit filled the air as the group shared a well-deserved meal.

\textbf{Glossary} : \emph{njangi} (n., tontine, groupe d'entraide camerounais) ; \emph{wheelbarrow} (n., brouette) ; \emph{gutter} (n., caniveau) ; \emph{rubbish} (n., ordures) ; \emph{council} (n., conseil municipal) ; \emph{gloves} (n., gants) ; \emph{brand new} (adj., tout neuf) ; \emph{well-deserved} (adj., bien mérité).

\textbf{Comprehension questions} :
\begin{enumerate}
\item How many participants came to the clean-up campaign?
\item What did Mrs Etoundi contribute?
\item What did the anonymous donor send?
\item Does Paul work voluntarily or by obligation? Quote the text to justify your answer.
\item Find in the text three words built from the verb \emph{to organize}.
\end{enumerate}
\end{textebox}

Observez le texte : \eng{volunteers}, \eng{volunteer}, \eng{voluntarily}, \eng{voluntary} apparaissent tous les quatre. Ce sont quatre formes d'une même famille. Repérons-les et classons-les.

\courssub{Les cinq grandes familles du travail communautaire}

\begin{definition}[Word family — famille de mots]
Une \cle{famille de mots} (\eng{word family}) est un ensemble de mots construits sur la même racine et qui remplissent des fonctions différentes dans la phrase : verbe (l'action), nom de personne (celui qui agit), nom d'action (le fait d'agir), adjectif (la qualité) et adverbe (la manière). En anglais, ces transformations se font surtout par \cle{suffixation} : on ajoute un élément à la fin du mot (\eng{-er}, \eng{-or}, \eng{-tion}, \eng{-ary}, \eng{-ly}...).
\end{definition}

\textbf{Famille 1 : to volunteer.} Le verbe \eng{to volunteer} (« vo-leun-\textbf{tir} ») signifie « se porter volontaire, travailler bénévolement ». Il donne le nom de personne \eng{a volunteer} (un volontaire, un bénévole), le nom d'action \eng{volunteering} (le bénévolat), l'adjectif \eng{voluntary} (bénévole, volontaire) et l'adverbe \eng{voluntarily} (bénévolement, de son plein gré).

\eng{Paul volunteers at the health centre every week.} « Paul travaille bénévolement au centre de santé chaque semaine. »

\eng{She is a volunteer in the Red Cross.} « Elle est bénévole à la Croix-Rouge. »

\eng{Volunteering builds strong communities.} « Le bénévolat construit des communautés solides. »

\eng{This is voluntary work; nobody is paid.} « C'est un travail bénévole ; personne n'est payé. »

\eng{He came voluntarily; nobody forced him.} « Il est venu de son plein gré ; personne ne l'a forcé. »

\textbf{Famille 2 : to participate.} Le verbe \eng{to participate} (participer) se construit avec la préposition \cle{in} : \eng{to participate in a meeting}. Il donne le nom d'action \eng{participation} (la participation) et le nom de personne \eng{a participant} (un participant). L'adjectif est \eng{participatory} (participatif).

\eng{All the pupils participated in the clean-up campaign.} « Tous les élèves ont participé à la campagne de nettoyage. »

\eng{Your participation is very important to us.} « Votre participation est très importante pour nous. »

\eng{Each participant received a certificate.} « Chaque participant a reçu un certificat. »

\eng{A participatory approach ensures everyone is heard.} « Une approche participative garantit que chacun est entendu. »

\textbf{Famille 3 : to organize.} Le verbe \eng{to organize} (organiser) donne le nom d'action \eng{organization} (organisation), le nom de personne \eng{an organizer} (un organisateur) et l'adjectif \eng{organized} (organisé). Notez l'orthographe britannique en \eng{-ise} : \eng{organise}, \eng{organisation} ; les deux formes (\eng{-ize} / \eng{-ise}) sont correctes.

\eng{The students organized a fundraising event.} « Les élèves ont organisé une collecte de fonds. »

\eng{This organisation helps orphans in Douala.} « Cette organisation aide les orphelins à Douala. »

\eng{The organizers thanked all the participants.} « Les organisateurs ont remercié tous les participants. »

\eng{She is a very organized person.} « C'est une personne très organisée. »

\textbf{Famille 4 : to contribute.} Le verbe \eng{to contribute} (contribuer) se construit avec \cle{to} : \eng{to contribute to a project}. Il donne \eng{a contribution} (une contribution) et \eng{a contributor} (un contributeur). L'adjectif est \eng{contributory} (contributif).

\eng{Everybody contributed money to buy the wheelbarrow.} « Tout le monde a contribué de l'argent pour acheter la brouette. »

\eng{Her voluntary contribution helped the project.} « Sa contribution bénévole a aidé le projet. »

\eng{He is a regular contributor to the community fund.} « Il contribue régulièrement à la caisse communautaire. »

\textbf{Famille 5 : to cooperate et to donate.} Le verbe \eng{to cooperate} (coopérer) se construit avec \cle{with} : \eng{to cooperate with the council}. Il donne \eng{cooperation} (coopération) et l'adjectif \eng{cooperative} (coopératif ; aussi nom : une coopérative). Le verbe \eng{to donate} (faire don de) donne \eng{a donation} (un don) et \eng{a donor} (un donateur, un donneur). L'adjectif morphologique est \eng{donative} (rare) ; on utilise souvent \eng{charitable} (caritatif) par lien sémantique.

\eng{The villagers cooperate with the NGO.} « Les villageois coopèrent avec l'ONG. »

\eng{Thank you for your cooperation.} « Merci de votre coopération. »

\eng{They formed a farming cooperative.} « Ils ont formé une coopérative agricole. »

\eng{She donated books to the school library.} « Elle a fait don de livres à la bibliothèque de l'école. »

\eng{The donation arrived on Monday.} « Le don est arrivé lundi. » / \eng{The donors prefer to remain anonymous.} « Les donateurs préfèrent rester anonymes. »

\begin{propriete}[Les suffixes qui créent les familles de mots]
Deux grands mécanismes permettent de deviner le sens et la nature d'un mot anglais inconnu :
\begin{itemize}
\item Les suffixes \cle{-tion} et \cle{-sion} forment des \textbf{noms d'action} : \eng{organize} $\rightarrow$ \eng{organization} (organisation), \eng{participate} $\rightarrow$ \eng{participation}, \eng{contribute} $\rightarrow$ \eng{contribution}, \eng{donate} $\rightarrow$ \eng{donation}, \eng{cooperate} $\rightarrow$ \eng{cooperation}. Remarquez que le \eng{-e} final du verbe disparaît : \eng{donate} $\rightarrow$ \eng{donation}.
\item Les suffixes \cle{-er} et \cle{-or} désignent la \textbf{personne qui fait l'action} : \eng{organize} $\rightarrow$ \eng{organizer}, \eng{contribute} $\rightarrow$ \eng{contributor}, \eng{donate} $\rightarrow$ \eng{donor}, \eng{work} $\rightarrow$ \eng{worker}, \eng{lead} $\rightarrow$ \eng{leader}.
\item Le suffixe \cle{-ly} transforme un adjectif en \textbf{adverbe} : \eng{voluntary} $\rightarrow$ \eng{voluntarily} (attention : \eng{-y} devient \eng{-i} devant \eng{-ly}).
\end{itemize}
Bonne nouvelle pour les francophones : les noms en \eng{-tion} se prononcent « cheun » en anglais, jamais « sion » comme en français.
\end{propriete}

\begin{center}
\begin{tabularx}{\linewidth}{|l|X|X|X|l|}
\hline
\rowcolor{popblueL} \textbf{Verbe} & \textbf{Personne} & \textbf{Action / chose} & \textbf{Adjectif / adverbe} & \textbf{Français} \\
\hline
to volunteer & a volunteer & volunteering / voluntary work & voluntary / voluntarily & se porter volontaire \\
\hline
to participate (in) & a participant & participation & participatory & participer (à) \\
\hline
to organize & an organizer & organization & organized & organiser \\
\hline
to contribute (to) & a contributor & contribution & contributory & contribuer (à) \\
\hline
to cooperate (with) & a cooperator & cooperation & cooperative / cooperatively & coopérer (avec) \\
\hline
to donate & a donor & a donation & donative / charitable & faire un don \\
\hline
\end{tabularx}
\end{center}

\begin{popfigure}
\begin{center}
\begin{tikzpicture}[
  grow=down,
  level distance=1.7cm,
  level 1/.style={sibling distance=4.4cm},
  edge from parent/.style={draw=popblue, thick, -{Stealth[length=2.5mm]}},
  every node/.style={align=center, font=\small}
]
\node[draw=popdark, fill=popblue, text=white, rounded corners, font=\bfseries, align=center]{to participate\\(verbe : participer)}
  child { node[draw=poporange, fill=poporangeL, rounded corners, align=center]{a participant\\(personne : participant)} edge from parent node[left, font=\scriptsize\itshape, text=popmuted]{-ant} }
  child { node[draw=popgreen, fill=popgreenL, rounded corners, align=center]{participation\\(action : participation)} edge from parent node[right, font=\scriptsize\itshape, text=popmuted]{-ation} }
  child { node[draw=poppurple, fill=poppurpleL, rounded corners, align=center]{participatory\\(adjectif : participatif)} edge from parent node[right, font=\scriptsize\itshape, text=popmuted]{-atory} };
\end{tikzpicture}
\end{center}
\legende{Arbre de dérivation : un verbe, trois mots nouveaux. Le suffixe indique la fonction du mot dans la phrase.}
\end{popfigure}

\begin{exemplebox}[Les familles en action dans une même phrase]
\eng{Her voluntary contribution helped the project.} « Sa contribution bénévole a aidé le projet. » (adjectif + nom d'action)

\eng{The organizers thanked all the participants.} « Les organisateurs ont remercié tous les participants. » (deux noms de personnes)

\eng{The volunteers participated voluntarily in the organization of the event.} « Les bénévoles ont participé de leur plein gré à l'organisation de la manifestation. » (quatre mots de trois familles différentes !)

\eng{Thanks to the donors' donations, the cooperative bought new tools.} « Grâce aux dons des donateurs, la coopérative a acheté de nouveaux outils. »
\end{exemplebox}

\courssub{Prononciation : placer l'accent sur la bonne syllabe}

En français, l'accent tombe presque toujours sur la dernière syllabe. En anglais, chaque mot a une \cle{syllabe accentuée} propre, qu'il faut apprendre avec le mot. Si vous accentuez la mauvaise syllabe, votre interlocuteur anglophone peut ne pas vous comprendre.

\begin{center}
\begin{tabularx}{\linewidth}{|l|l|X|}
\hline
\rowcolor{popblueL} \textbf{Mot} & \textbf{Prononciation (lettres françaises)} & \textbf{Syllabe accentuée} \\
\hline
volunteer & « vo-leun-tir » & volun\textbf{TEER} (3\textsuperscript{e} syllabe) \\
\hline
voluntary & « \textbf{vo}-leun-te-ri » & \textbf{VOL}untary (1\textsuperscript{re} syllabe) \\
\hline
voluntarily & « vo-leun-\textbf{te}-ri-li » & volun\textbf{TE}rily (3\textsuperscript{e} syllabe) \\
\hline
participation & « par-ti-si-pé-cheun » & partici\textbf{PA}tion (4\textsuperscript{e} syllabe) \\
\hline
contribution & « kon-tri-biou-cheun » & contri\textbf{BU}tion (3\textsuperscript{e} syllabe) \\
\hline
solidarity & « so-li-da-ri-ti » & soli\textbf{DA}rity (3\textsuperscript{e} syllabe) \\
\hline
community & « ke-miou-ni-ti » & com\textbf{MU}nity (2\textsuperscript{e} syllabe) \\
\hline
organization & « or-ga-naï-zé-cheun » & organi\textbf{ZA}tion \\
\hline
\end{tabularx}
\end{center}

\begin{attention}[L'accent se déplace dans la famille !]
Piège classique : l'accent n'est pas sur la même syllabe dans tous les mots d'une famille. On dit \textbf{VOL}untary (accent au début) mais volun\textbf{TEER} (accent à la fin). De même : \textbf{PAR}ticipant mais partici\textbf{PA}tion. En français, « volontaire » et « volonté » s'accentuent pareil ; en anglais, non. Apprenez toujours le mot \emph{avec} son accent.
\end{attention}

\begin{methode}[Bien prononcer les mots du thème]
\begin{enumerate}
\item Dans votre cahier, écrivez chaque mot nouveau en \textbf{majuscules sur la syllabe accentuée} : volunTEER, parTiciPAtion, comMUnity, contriBUtion.
\item Lisez le mot à voix haute trois fois, en exagérant la syllabe forte.
\item Passez immédiatement à la \textbf{phrase complète} : on ne prononce jamais un mot isolé dans la vraie vie. Dites : \eng{She is a volunTEER.} puis \eng{Her parti-ciPAtion was remarkable.}
\item Enregistrez-vous avec votre téléphone et comparez avec la lecture du professeur.
\item Révisez par familles : dites d'affilée \eng{volunteer, voluntary, voluntarily} en marquant l'accent à chaque fois.
\end{enumerate}
\end{methode}

\courssub{Orthographe : les mots pièges du thème}

\begin{propriete}[Le doublement des lettres]
Plusieurs mots du champ lexical du travail communautaire contiennent des \cle{lettres doubles} que les francophones oublient souvent :
\begin{itemize}
\item \eng{committee} (comité) : double \textbf{m}, double \textbf{t}, double \textbf{e}. Trois doublements dans un seul mot !
\item \eng{rubbish} (ordures) : double \textbf{b}.
\item \eng{gutter} (caniveau) : double \textbf{t}.
\item \eng{community} (communauté) : double \textbf{m} — comme en français « communauté », pensez-y.
\item \eng{neighbourhood} (quartier) : double \textbf{o} à la fin (\eng{-hood}).
\end{itemize}
Astuce : découpez le mot en morceaux — \eng{com-mit-tee}, \eng{neigh-bour-hood} — et épellez-le à voix haute.
\end{propriete}

\begin{attention}[Les trois fautes d'orthographe les plus fréquentes]
\begin{itemize}
\item \eng{voluntary} s'écrit avec \textbf{-tary} et non « -tery » ni « -tory » : \eng{volun-ta-ry}. Prononcez « vo-leun-te-ri » pour sentir le \textbf{a}.
\item \eng{wheelbarrow} (brouette) s'écrit \textbf{en un seul mot} : jamais « wheel barrow » ni « weelbarrow ». C'est \eng{wheel} (roue) + \eng{barrow}.
\item \eng{neighbourhood} s'écrit avec \textbf{-our} en anglais britannique (\eng{neighbourhood}, \eng{neighbour}) et avec \textbf{-or} en anglais américain (\eng{neighborhood}, \eng{neighbor}). \textbf{Les deux sont acceptés}, mais il faut rester cohérent dans une même copie : ne mélangez pas \eng{colour} (britannique) et \eng{neighbor} (américain) dans le même paragraphe. Au Cameroun, on privilégie généralement l'anglais britannique.
\end{itemize}
\end{attention}

\begin{savaistu}[Pourquoi l'anglais a-t-il deux orthographes ?]
En 1828, l'Américain Noah Webster publia un dictionnaire qui simplifiait délibérément l'orthographe anglaise pour marquer l'indépendance culturelle des États-Unis : \eng{colour} devint \eng{color}, \eng{neighbour} devint \eng{neighbor}, \eng{organise} devint \eng{organize}. C'est pourquoi le même mot peut avoir deux graphies correctes aujourd'hui. Au Cameroun, pays officiellement bilingue hérité de la colonisation britannique et française, c'est la norme britannique qui domine dans les manuels scolaires.
\end{savaistu}

\begin{exoresolu}[Compléter avec le bon mot de la famille]
Énoncé — Complétez chaque phrase avec la forme correcte du mot entre parenthèses.
\begin{enumerate}
\item Many \ldots\ldots\ldots came to clean the market. (volunteer)
\item Her \ldots\ldots\ldots in the project was remarkable. (participate)
\item The \ldots\ldots\ldots of the event distributed T-shirts. (organize)
\item He made a generous \ldots\ldots\ldots to the school fund. (contribute)
\item The \ldots\ldots\ldots gave money anonymously. (donate)
\end{enumerate}
\tcblower
Solution détaillée :
\begin{enumerate}
\item \eng{Many \textbf{volunteers} came to clean the market.} « De nombreux bénévoles sont venus nettoyer le marché. » Il faut un \textbf{nom de personne} au pluriel (\eng{many} + pluriel).
\item \eng{Her \textbf{participation} in the project was remarkable.} « Sa participation au projet était remarquable. » Après l'adjectif possessif \eng{her}, il faut un \textbf{nom d'action} : suffixe \eng{-ation}.
\item \eng{The \textbf{organizers} of the event distributed T-shirts.} « Les organisateurs de la manifestation ont distribué des T-shirts. » Il faut la \textbf{personne} : suffixe \eng{-er}, au pluriel.
\item \eng{He made a generous \textbf{contribution} to the school fund.} « Il a fait une contribution généreuse à la caisse de l'école. » Après l'article \eng{a} et l'adjectif \eng{generous}, il faut un \textbf{nom} : \eng{contribute} perd son \eng{-e} et prend \eng{-tion}.
\item \eng{The \textbf{donor} gave money anonymously.} « Le donateur a donné de l'argent anonymement. » Il faut la \textbf{personne} : suffixe \eng{-or} (et non « donner », qui n'existe pas en anglais !).
\end{enumerate}
Stratégie : avant de choisir la forme, demandez-vous toujours quelle \textbf{fonction} le mot occupe dans la phrase (sujet ? complément après un article ? qualité ?). La fonction dicte le suffixe.
\end{exoresolu}

\begin{exercice}[À vous de jouer]
\begin{enumerate}
\item Complétez avec la bonne forme du mot entre parenthèses : a) The \ldots\ldots\ldots of the neighbourhood worked together. (inhabit / community) b) They work \ldots\ldots\ldots ; they receive no salary. (voluntary) c) We thank every \ldots\ldots\ldots for his or her help. (participate) d) Their \ldots\ldots\ldots with the council made the project possible. (cooperate)
\item Corrigez l'orthographe : \eng{comittee}, \eng{weelbarrow}, \eng{rubish}, \eng{voluntery}, \eng{guter}.
\item Classez ces mots par famille et indiquez leur nature (verbe, personne, action, adjectif) : \eng{donation}, \eng{organize}, \eng{contributor}, \eng{cooperative}, \eng{participant}, \eng{voluntarily}.
\item Prononciation : soulignez la syllabe accentuée dans \eng{community}, \eng{solidarity}, \eng{contribution}, \eng{voluntary}, puis lisez chaque mot dans une phrase complète.
\end{enumerate}
\end{exercice}

\begin{corrige}[Exercice « À vous de jouer »]
\begin{enumerate}
\item a) \eng{community} (nom, sujet) ; b) \eng{voluntarily} (adverbe : il modifie le verbe \eng{work}) ; c) \eng{participant} (personne, après \eng{every}) ; d) \eng{cooperation} (nom d'action, après le possessif \eng{their}).
\item \eng{committee} (double m, t, e) ; \eng{wheelbarrow} (un seul mot, \eng{wheel} avec h) ; \eng{rubbish} (double b) ; \eng{voluntary} (avec a) ; \eng{gutter} (double t).
\item \eng{donation} : action (famille de \eng{donate}) ; \eng{organize} : verbe ; \eng{contributor} : personne (famille de \eng{contribute}) ; \eng{cooperative} : adjectif ou nom ; \eng{participant} : personne (famille de \eng{participate}) ; \eng{voluntarily} : adverbe (famille de \eng{volunteer}).
\item com\textbf{MU}nity, soli\textbf{DA}rity, contri\textbf{BU}tion, \textbf{VOL}untary. Exemples : \eng{Our comMUnity is united.} / \eng{SoliDArity is our strength.} / \eng{Her contriBUtion helped us.} / \eng{This is VOLuntary work.}
\end{enumerate}
\end{corrige}

\begin{aretenir}[L'essentiel de la section]
\begin{itemize}
\item Apprenez le vocabulaire \textbf{par familles} : verbe, personne (\eng{-er}/\eng{-or}), action (\eng{-tion}/\eng{-sion}), adjectif, adverbe (\eng{-ly}).
\item La \textbf{fonction du mot dans la phrase} vous indique quel membre de la famille choisir.
\item L'\textbf{accent tonique} change de place à l'intérieur d'une même famille : \textbf{VOL}untary mais volun\textbf{TEER}.
\item Orthographe à retenir : \eng{committee}, \eng{rubbish}, \eng{gutter} (lettres doubles) ; \eng{wheelbarrow} en un seul mot ; \eng{voluntary} avec un \textbf{a} ; \eng{neighbourhood} (GB) / \eng{neighborhood} (US), sans mélanger les deux normes.
\end{itemize}
\end{aretenir}

Votre lexique est désormais structuré, bien prononcé et bien orthographié. Il ne reste plus qu'une étape pour le rendre vraiment vôtre : le réemployer dans des situations réelles. C'est l'objet de la section suivante, « Réemploi en contexte : décrire et promouvoir une action », où vous utiliserez toutes ces familles de mots pour parler et écrire sur le travail communautaire.

\coursec{Réemploi en contexte : décrire et promouvoir une action}

Dans la section précédente, nous avons travaillé sur les familles de mots (\eng{volunteer / voluntary / voluntarily}), la prononciation et l'orthographe du lexique du travail communautaire. Vous possédez maintenant les \emph{outils} : il est temps de les \emph{utiliser}. Cette dernière section est consacrée au \cle{réemploi} du vocabulaire dans des situations de communication réelles : raconter une action passée, mobiliser des voisins à l'oral, et rédiger une affiche en anglais. C'est exactement ce que l'on attend de vous à l'examen : non pas réciter des listes, mais \emph{faire quelque chose} avec les mots.

\courssub{Lecture : un récit de travail communautaire au prétérit}

Commençons par observer un texte modèle. Lisez-le deux fois : une première fois pour le sens général, une deuxième fois en repérant les mots du lexique de la leçon, écrits en gras.

\begin{textebox}[A bridge for our village — a community report]
Last month, our community \textbf{carried out} a building project. The old wooden bridge over the River Mungo was dangerous, so the \textbf{quarter head} called a meeting and \textbf{mobilised} the whole village. On Saturday morning, more than fifty \textbf{volunteers} \textbf{rolled up their sleeves} and started work at six o'clock. The young men carried stones and sand, the women prepared food for the workers, and local traders \textbf{donated} cement, nails and planks. A retired carpenter from Buea \textbf{supervised} the construction and showed the youths how to fix the wooden beams. Everyone \textbf{did their bit}: even the children brought water to the workers. By sunset, the new bridge was finished. The quarter head thanked all the \textbf{participants} warmly. It was tiring but \textbf{rewarding}. Now the farmers can reach the market safely, and the whole village is proud of this \textbf{joint effort}. Next month, we plan to \textbf{clean up} the area around the market. We hope that even more people will \textbf{take part} and \textbf{lend a hand}. \textbf{Together}, we can do great things!
\end{textebox}

\textbf{Glossaire}

\begin{center}
\begin{tabularx}{\linewidth}{|l|l|X|}
\hline
\rowcolor{popblueL} Mot anglais & Nature & Traduction française \\
\hline
to carry out (a project) & verbe & mener à bien, réaliser (un projet) \\
\hline
quarter head & nom & chef de quartier \\
\hline
to mobilise & verbe & mobiliser \\
\hline
to roll up one's sleeves & expression & retrousser ses manches (se mettre au travail) \\
\hline
to donate & verbe & donner, faire don de \\
\hline
to supervise & verbe & superviser, diriger (les travaux) \\
\hline
to do one's bit & expression & apporter sa contribution, faire sa part \\
\hline
rewarding & adjectif & gratifiant, qui vaut la peine \\
\hline
joint effort & nom & effort collectif \\
\hline
to lend a hand & expression & donner un coup de main \\
\hline
\end{tabularx}
\end{center}

\textbf{Questions de compréhension}

\begin{enumerate}
\item Why did the community decide to build a new bridge?
\item Who donated the building materials?
\item What did the children do during the work?
\item What is the community planning to do next month?
\item Find in the text two expressions that mean « to help ».
\end{enumerate}

\begin{exoresolu}[Questions de compréhension — corrigé]
Répondez aux cinq questions ci-dessus en phrases complètes.
\tcblower
\begin{enumerate}
\item The community decided to build a new bridge because the old wooden bridge was dangerous.
\item Local traders donated the building materials (cement, nails and planks).
\item The children brought water to the workers.
\item Next month, the community is planning to clean up the area around the market.
\item Two expressions meaning « to help » are \eng{to do one's bit} and \eng{to lend a hand}.
\end{enumerate}
\end{exoresolu}

\begin{attention}[Le prétérit pour raconter une action passée]
Remarquez que tout le récit est au \cle{prétérit simple} : \eng{carried}, \eng{mobilised}, \eng{donated}, \eng{supervised}, \eng{brought}, \eng{thanked}. Quand vous racontez un travail communautaire déjà terminé, utilisez le prétérit, et non le présent. Erreur fréquente des francophones : mélanger les temps (\eng{Last Saturday, the volunteers come and clean the market} est faux). Dites : \eng{Last Saturday, the volunteers came and cleaned the market.} « Samedi dernier, les volontaires sont venus et ont nettoyé le marché. »
\end{attention}

\courssub{Dialogue modèle : mobiliser les habitants}

Le travail communautaire commence souvent par une conversation : quelqu'un invite, propose, organise. Voici un dialogue modèle entre un chef de quartier et une habitante de Yaoundé.

\begin{textebox}[Mr Ekane mobilises his neighbours]
\textbf{Mr Ekane:} Good morning, Mrs Biya! Our quarter is organising a clean-up campaign this Saturday. Would you like to join us?

\textbf{Mrs Biya:} Good morning! Yes, with pleasure. What time does it start?

\textbf{Mr Ekane:} We meet at the market square at 7 a.m. sharp. Please bring a broom and a pair of gloves if you have one.

\textbf{Mrs Biya:} No problem. What exactly will we do?

\textbf{Mr Ekane:} The youths will sweep the streets and clear the gutters, and the adults will collect the rubbish and burn it. My son will supervise the teams.

\textbf{Mrs Biya:} That's a good idea. Can I bring my daughter? She is only twelve.

\textbf{Mr Ekane:} Of course! Everyone is welcome. Even children can do their bit.

\textbf{Mrs Biya:} All right, count me in! Together, we can keep our quarter clean.

\textbf{Mr Ekane:} Thank you very much, Mrs Biya. See you on Saturday!
\end{textebox}

\textbf{Glossaire} : \eng{clean-up campaign} (nom) : campagne de nettoyage ; \eng{at 7 a.m. sharp} : à 7 heures précises ; \eng{gutter} (nom) : caniveau ; \eng{count me in} (expression) : je suis partant, comptez sur moi.

\begin{aretenir}[Inviter, accepter, répartir les tâches]
\begin{itemize}
\item \textbf{Inviter} : \eng{Would you like to join us?} « Voulez-vous vous joindre à nous ? » / \eng{Can you lend a hand?} « Peux-tu donner un coup de main ? »
\item \textbf{Accepter} : \eng{Yes, with pleasure.} / \eng{Count me in!} « Je suis partant ! »
\item \textbf{Refuser poliment} : \eng{I'm sorry, I can't. I have to work.} « Je suis désolé, je ne peux pas. Je dois travailler. »
\item \textbf{Répartir les tâches} : \eng{The youths will sweep\ldots{} and the adults will collect\ldots} Notez l'emploi de \cle{will} pour annoncer la répartition des rôles.
\end{itemize}
\end{aretenir}

\begin{savaistu}[Le travail communautaire au Cameroun : l'esprit « Njangi »]
Au Cameroun, le travail communautaire est une tradition vivante. On l'appelle souvent \eng{communal labour} ou \eng{self-help}. Dans les villages comme dans les quartiers urbains (Yaoundé, Douala, Buea, Bamenda), les habitants se réunissent pour : \eng{clear the bush} (débroussailler les routes), \eng{build a community hall / a school / a health centre}, organiser des \eng{clean-up campaigns} (journées de salubrité, souvent le jeudi ou le samedi), ou préparer les fêtes comme le \eng{Ngondo} à Douala. Le \textbf{quarter head} (chef de quartier) ou le \textbf{traditional ruler} (chef traditionnel) \textbf{mobilise} la population. Chaque famille \textbf{does its bit} : les hommes font les travaux lourds, les femmes préparent le \eng{food for the workers} (souvent du \eng{corn fufu} et de la sauce \eng{eru} ou \eng{ndolé}), les jeunes vont chercher l'eau. Cet \textbf{joint effort} renforce la cohésion sociale. On dit souvent : \eng{Unity is strength} — « L'union fait la force ».
\end{savaistu}

\begin{experience}[Jeu de rôle : La réunion de quartier]
\textbf{Objectif :} S'entraîner à mobiliser, inviter, accepter ou refuser, et répartir les tâches.
\textbf{Consignes :}
\begin{enumerate}
\item Formez des groupes de 4 à 5 élèves. L'un joue le \textbf{quarter head}, les autres sont des \textbf{inhabitants}.
\item Le quarter head annonce un projet concret : \eng{clean the well}, \eng{build a latrine for the school}, \eng{paint the community hall}, \eng{plant trees along the main road}.
\item Il invite : \eng{We are organising a... Would you like to join us?}
\item Les habitants réagissent : questions (\eng{When? Where? What to bring?}), acceptation (\eng{Count me in!}), refus poli (\eng{I'm sorry, I can't...}), proposition d'aide (\eng{I can bring tools / I will supervise...}).
\item Le quarter head répartit : \eng{Group A will... Group B will...}
\item Chaque groupe présente sa saynète (1 à 2 minutes) devant la classe.
\end{enumerate}
\textbf{Lexique utile :} \eng{tools} (outils), \eng{paint / brushes} (peinture / pinceaux), \eng{seedlings} (plants), \eng{well} (puits), \eng{latrines} (latrines).
\end{experience}

\courssub{La structure d'une annonce ou d'une affiche en anglais}

Pour promouvoir une action communautaire, on rédige souvent une affiche (\eng{a poster}) ou une annonce (\eng{a notice}). Une affiche anglaise efficace suit toujours la même architecture.

\begin{methode}[Rédiger une annonce en 5 questions + 1 slogan]
Votre affiche doit répondre, dans l'ordre, à cinq questions, puis se terminer par une formule d'appel :
\begin{enumerate}
\item \textbf{Who?} — Qui organise ? (\eng{The Youth Association of Nkolbisson presents\ldots})
\item \textbf{What?} — Quelle activité ? (\eng{A clean-up campaign / A bridge-building project})
\item \textbf{When?} — Quand ? (\eng{Saturday, 15th June, from 7 a.m. to noon})
\item \textbf{Where?} — Où ? (\eng{At the market square})
\item \textbf{What to bring?} — Quel matériel apporter ? (\eng{Bring: brooms, gloves, buckets})
\item \textbf{Slogan final} — une formule courte et chaleureuse : \eng{Come and lend a hand!} « Venez donner un coup de main ! »
\end{enumerate}
\end{methode}

\begin{exemplebox}[Exemples de formules d'affiche traduites]
\eng{Bring your brooms and join us this Saturday at 7 a.m.!} « Apportez vos balais et rejoignez-nous samedi à 7 heures ! »

\eng{Together, we can keep our quarter clean.} « Ensemble, nous pouvons garder notre quartier propre. »

\eng{Everyone is welcome — come and do your bit!} « Tout le monde est le bienvenu — venez apporter votre contribution ! »
\end{exemplebox}

\begin{attention}[Heure et date : ne calquez pas le français]
Dans une affiche anglaise, l'heure s'écrit \eng{7 a.m.} (et non « 7 h » ni « 7 o'clock a.m. ») et la date s'écrit \eng{Saturday, 15th June} : le \textbf{jour} d'abord, puis la \textbf{date}, puis le \textbf{mois}. Ne calquez jamais l'ordre français « samedi 15 juin » mot à mot sans adapter : en anglais, on met une virgule après le jour et le mois se place après le numéro. De même, « de 7 h à midi » devient \eng{from 7 a.m. to noon}, et non « from 7 hours to 12 hours ».
\end{attention}

\begin{popfigure}
\begin{center}
\begin{tikzpicture}[node distance=0.55cm,
box/.style={draw=popblue, fill=popblueL, rounded corners=3pt, align=center, inner sep=6pt, text width=6.5cm},
arr/.style={-{Stealth[length=3mm]}, thick, poporange}]
\node[box, fill=poppurpleL, draw=poppurple, align=center] (titre){\textbf{TITRE ACCROCHEUR}\\ \eng{CLEAN-UP DAY!}};
\node[box, below=of titre, align=center] (info){\textbf{What? Where? When?}\\ \eng{Clean-up campaign — Market square}\\ \eng{Saturday, 15th June, 7 a.m.}};
\node[box, fill=popgreenL, draw=popgreen, below=of info, align=center] (bring){\textbf{What to bring}\\ \eng{Bring: brooms, gloves, buckets}};
\node[box, fill=popgoldL, draw=popgold, below=of bring, align=center] (slogan){\textbf{SLOGAN}\\ \eng{Come and lend a hand!}};
\draw[arr] (titre) -- (info);
\draw[arr] (info) -- (bring);
\draw[arr] (bring) -- (slogan);
\end{tikzpicture}
\end{center}
\legende{Structure d'une affiche communautaire en anglais : du titre au slogan.}
\end{popfigure}

\courssub{Production guidée : du texte à trous à la rédaction libre}

Passons maintenant à votre production. On procède en deux étapes : d'abord un texte à trous (réemploi contrôlé), ensuite des phrases libres sur votre propre quartier (réemploi créatif).

\begin{exoresolu}[Texte à trous : la campagne de nettoyage de Bonabéri]
Complétez le texte avec les mots suivants : \eng{volunteers — lend a hand — did their bit — clean-up — rewarding — donated — mobilised}.

« Last Saturday, the quarter head of Bonabéri \ldots\ the inhabitants for a big \ldots\ campaign. About forty \ldots\ came with brooms and gloves. A local shopkeeper \ldots\ plastic bags and drinking water. Everyone \ldots\ : the youths swept the streets and the adults cleared the gutters. At the end of the day, the quarter head thanked everybody and said: “Next month, come again and \ldots\ !” It was a tiring but \ldots\ day. »
\tcblower
« Last Saturday, the quarter head of Bonabéri \textbf{mobilised} the inhabitants for a big \textbf{clean-up} campaign. About forty \textbf{volunteers} came with brooms and gloves. A local shopkeeper \textbf{donated} plastic bags and drinking water. Everyone \textbf{did their bit}: the youths swept the streets and the adults cleared the gutters. At the end of the day, the quarter head thanked everybody and said: “Next month, come again and \textbf{lend a hand}!” It was a tiring but \textbf{rewarding} day. »

Remarquez : \eng{did their bit} est une expression idiomatique (ne pas traduire mot à mot), et \eng{mobilised}, \eng{donated} sont au prétérit car l'action est terminée.
\end{exoresolu}

\begin{exercice}[À vous de rédiger !]
Rédigez \textbf{cinq phrases} en anglais sur un travail communautaire réel de votre quartier ou de votre village (construction d'un pont, nettoyage du marché, plantation d'arbres, réparation d'une route\ldots). Consignes :
\begin{enumerate}
\item Utilisez le prétérit pour raconter les faits.
\item Employez au moins six mots du lexique de la leçon.
\item Incluez au moins deux collocations (ex. \eng{carry out a project}, \eng{roll up one's sleeves}).
\item Terminez par une phrase d'appel pour inviter les autres (\eng{Come and lend a hand!}).
\end{enumerate}
\end{exercice}

\begin{corrige}[Exercice : production libre — modèle de réponse]
Voici un modèle possible (le vôtre sera différent, et c'est très bien) :

\eng{Last month, the inhabitants of my quarter carried out a clean-up campaign around the market. The quarter head mobilised about thirty volunteers, and we all rolled up our sleeves. A trader donated brooms and buckets, and everyone did their bit. We swept the streets and cleared the gutters from 7 a.m. to noon. It was hard work, but very rewarding. Next time, come and lend a hand!}

« Le mois dernier, les habitants de mon quartier ont mené une campagne de nettoyage autour du marché. Le chef de quartier a mobilisé une trentaine de volontaires, et nous avons tous retroussé nos manches. Un commerçant a donné des balais et des seaux, et chacun a fait sa part. Nous avons balayé les rues et dégagé les caniveaux de 7 heures à midi. C'était un travail dur, mais très gratifiant. La prochaine fois, venez donner un coup de main ! »
\end{corrige}

\courssub{Auto-évaluation : vérifier sa production}

Avant de rendre votre texte, relisez-vous avec cette grille. Cochez chaque critère seulement s'il est réellement respecté.

\begin{center}
\begin{tabularx}{\linewidth}{|X|c|}
\hline
\rowcolor{popblueL} Critère de réussite & Oui / Non \\
\hline
J'ai utilisé au moins \textbf{6 mots du lexique} de la leçon (volunteer, donate, mobilise\ldots) & \\
\hline
J'ai utilisé au moins \textbf{2 collocations} correctes (carry out a project, roll up one's sleeves\ldots) & \\
\hline
J'ai utilisé au moins \textbf{1 expression idiomatique} (do one's bit, lend a hand\ldots) & \\
\hline
Mes verbes sont au \textbf{prétérit} pour raconter les faits passés & \\
\hline
Je n'ai fait \textbf{aucun faux ami} (\eng{actually} $\neq$ « actuellement », \eng{to assist} $\neq$ « assister à ») & \\
\hline
Ma date et mon heure suivent le \textbf{modèle anglais} (\eng{Saturday, 15th June, 7 a.m.}) & \\
\hline
J'ai terminé par une \textbf{formule d'appel} (\eng{Come and lend a hand!}) & \\
\hline
\end{tabularx}
\end{center}

\begin{aretenir}[L'essentiel de la section]
\begin{itemize}
\item Pour \textbf{raconter} une action communautaire : prétérit simple + lexique de la leçon (\eng{carried out, mobilised, donated, did their bit}).
\item Pour \textbf{mobiliser} à l'oral : inviter (\eng{Would you like to join us?}), répondre (\eng{Count me in!}), répartir les tâches avec \eng{will}.
\item Pour \textbf{rédiger une affiche} : Who? What? When? Where? What to bring? + un slogan final.
\item Attention aux \textbf{faux amis} et au format anglais de la \textbf{date et de l'heure}.
\end{itemize}
Vous êtes désormais capable non seulement de reconnaître le vocabulaire du travail communautaire, mais de l'employer pour décrire, inviter et promouvoir — à l'écrit comme à l'oral. C'est la marque d'un vrai utilisateur de la langue !
\end{aretenir}

\newpage

\coursec{Activité d'intégration}

\courssub{Objectifs de l'activité}

À l'issue de cette activité d'intégration, l'élève sera capable de :

\begin{itemize}
\item réemployer à l'écrit au moins \cle{huit mots du lexique} de la leçon (\eng{communal work, community, volunteer, clean-up, tools, neighbourhood}, etc.) dans un document authentique ;
\item utiliser correctement au moins \cle{deux collocations} étudiées (par exemple \eng{to take part in}, \eng{to lend a hand}, \eng{to clean up}, \eng{to work together}) ;
\item intégrer l'expression idiomatique \eng{Many hands make light work!} dans un contexte pertinent ;
\item éviter les \cle{faux amis} et les calques du français signalés dans la leçon (\eng{actually} $\neq$ « actuellement », \eng{to assist} $\neq$ « assister à », etc.) ;
\item présenter oralement, en une minute, une affiche en utilisant un anglais clair, une prononciation correcte et des structures simples du niveau A2+/B1 ;
\item évaluer la production des autres groupes selon des critères explicites et voter de manière argumentée.
\end{itemize}

\courssub{Situation}

\textbf{Contexte.} Vous êtes membres du club d'anglais (\eng{English Club}) du lycée de Biyem-Assi, à Yaoundé. Après la saison des pluies, les rues de votre quartier sont sales : les caniveaux sont bouchés, les ordures s'accumulent et les moustiques se multiplient. Le chef du quartier, en collaboration avec la mairie, a décidé d'organiser une grande journée de travaux communautaires : \eng{Operation Clean Quarter}. Votre club d'anglais est chargé de concevoir une \cle{affiche en anglais} pour inviter toute la population à participer, puis de la présenter devant la classe, qui joue le rôle du comité de quartier.

Lisez d'abord le message que le chef du quartier a envoyé au président du club :

\begin{textebox}[Message from the Quarter Head]
Dear English Club members,

Our neighbourhood needs our help. After the heavy rains, the gutters are blocked and rubbish is everywhere. On Saturday 15th June, we are organising a big communal work day: “Operation Clean Quarter”. Everybody is invited to take part: young people, parents, shopkeepers and teachers.

We will clean up the streets, clear the gutters and plant trees near the primary school. Please bring your own tools: machetes, rakes, brooms, buckets and gloves. We will meet at 7 a.m. in front of the community hall. Food and drinks will be offered to all volunteers.

Please design an attractive poster in English to mobilise the population. Remember: many hands make light work!

The Quarter Head
\end{textebox}

\begin{definition}[Glossaire utile]
\eng{gutter} (n.) : caniveau ; \eng{rubbish} (n., indénombrable) : ordures ; \eng{to clear} (v.) : dégager, déboucher ; \eng{to plant trees} (v.) : planter des arbres ; \eng{tools} (n.) : outils, matériel ; \eng{community hall} (n.) : salle communautaire ; \eng{to mobilise} (v.) : mobiliser ; \eng{sharp} (adv.) : précises (pour l'heure : \eng{at 7 a.m. sharp} = « à 7 heures précises »).
\end{definition}

\courssub{Consignes}

\textbf{Tâche 1 — Compréhension du message (travail individuel, 5 min).} Lisez le message du chef du quartier et répondez par écrit, en anglais et en phrases complètes :
\begin{enumerate}
\item What is the problem in the neighbourhood?
\item When and where will the communal work take place?
\item What three activities are planned?
\item What must the participants bring?
\end{enumerate}

\textbf{Tâche 2 — Préparation du lexique (par groupe de 4, 10 min).} Dans votre cahier, complétez le tableau suivant avec des mots de la leçon, puis choisissez ensemble les \cle{huit mots} que vous utiliserez sur votre affiche :

\begin{center}
\begin{tabularx}{\linewidth}{|X|X|X|}
\hline
\rowcolor{popblueL} People & Activities & Tools and places \\
\hline
volunteer, \ldots & to clean up, \ldots & broom, \ldots \\
\hline
\end{tabularx}
\end{center}

Vérifiez dans la leçon l'orthographe et la prononciation de chaque mot choisi. Attention aux faux amis : n'écrivez jamais \eng{assist} pour « assister à » ; dites \eng{to take part in} ou \eng{to attend}.

\textbf{Tâche 3 — Conception de l'affiche (par groupe, 20 min).} Sur une feuille, créez votre affiche. Elle doit obligatoirement contenir :
\begin{enumerate}
\item un \cle{titre accrocheur} en anglais (court, avec un verbe d'action ou un slogan) ;
\item le nom de l'opération : \eng{Operation Clean Quarter} ;
\item la \cle{date}, l'\cle{heure} et le \cle{lieu} du rendez-vous ;
\item les \cle{activités prévues} (au moins trois verbes d'action) ;
\item le \cle{matériel à apporter} ;
\item au moins \cle{8 mots du lexique} de la leçon, soulignés ;
\item au moins \cle{2 collocations} correctes (par exemple \eng{to take part in communal work}, \eng{to lend a hand}, \eng{to work together as a team}) ;
\item l'expression idiomatique \eng{Many hands make light work!} bien visible ;
\item une phrase d'invitation chaleureuse à la population.
\end{enumerate}

\textbf{Tâche 4 — Présentation orale (1 minute par groupe).} Chaque groupe présente son affiche à la classe. Répartissez la parole entre les quatre membres. Utilisez des structures simples : \eng{We invite you to\ldots}, \eng{Come and join us!}, \eng{Don't forget to bring\ldots}. Soignez la prononciation des mots difficiles : \eng{volunteer} « vo-lon-tir », \eng{neighbourhood} « né-beur-houd », \eng{communal} « ko-miou-nal ».

\textbf{Tâche 5 — Vote et justification (classe entière, 10 min).} La classe vote pour l'affiche la plus convaincante selon la grille suivante (1 point par critère) : titre accrocheur ; informations complètes ; 8 mots du lexique bien employés ; 2 collocations correctes ; expression idiomatique présente ; aucun faux ami ; présentation orale claire. Justifiez votre vote en une phrase : \eng{I vote for Group 2 because\ldots}

\courssub{Ressources à mobiliser}

\begin{aretenir}[Boîte à outils de la leçon]
\begin{itemize}
\item \textbf{Lexique :} \eng{communal work, community, neighbourhood, volunteer, to take part, to help, to clean up, to repair, to plant, tools (broom, rake, machete, bucket, gloves), community hall, quarter head}.
\item \textbf{Collocations :} \eng{to take part in communal work} ; \eng{to lend a hand} ; \eng{to work together} ; \eng{to clean up the streets} ; \eng{to plant trees} ; \eng{to bring your own tools}.
\item \textbf{Expression idiomatique :} \eng{Many hands make light work!} (« À plusieurs, le travail devient léger »).
\item \textbf{Structures utiles :} impératifs pour inviter (\eng{Come! Join us! Bring your gloves!}) ; \eng{We will + base verbale} pour annoncer ; \eng{Don't forget to + base verbale}.
\item \textbf{Faux amis à éviter :} \eng{actually} = « en fait » (et non « actuellement » $\rightarrow$ \eng{currently}) ; \eng{to assist} = « aider » (et non « assister à » $\rightarrow$ \eng{to attend, to take part in}) ; \eng{a meeting} = « une réunion » (et non « un meeting » politique).
\end{itemize}
\end{aretenir}

\begin{attention}[Pièges fréquents dans ce type de production]
\begin{itemize}
\item Ne traduisez pas « travaux communautaires » par \eng{communitarian works} : dites \eng{communal work} ou \eng{community work}.
\item \eng{Everyone is invited}, et non \eng{All people is invited} : \eng{everyone} est singulier, et \eng{people} est un pluriel (\eng{people are}).
\item L'heure se dit \eng{at 7 a.m.}, et la date \eng{on Saturday 15th June} : attention aux prépositions \eng{at} (heure) et \eng{on} (jour).
\item Pas de « s » à \eng{equipment} ni à \eng{rubbish} : ce sont des indénombrables.
\item « Participer à » se dit \eng{to take part in} ou \eng{to participate in}, jamais \eng{to participate to}.
\end{itemize}
\end{attention}

\courssub{Proposition de production et corrigé commenté}

\begin{exoresolu}[Affiche modèle — Group 3, Biyem-Assi]
Énoncé : produire l'affiche complète répondant à toutes les contraintes, puis la présenter oralement.

\textbf{OPERATION CLEAN QUARTER!}

\textbf{Let's clean up Biyem-Assi together!}

Dear neighbours, our community needs YOU!

Take part in our big communal work day:

\textbf{When?} On Saturday 15th June, at 7 a.m. sharp.

\textbf{Where?} In front of the community hall.

\textbf{What?} We will clean up the streets, clear the gutters and plant trees near the primary school.

\textbf{Bring your own tools:} brooms, rakes, machetes, buckets and gloves.

Free food and drinks for all volunteers!

Come and lend a hand — let us work together for a clean and healthy neighbourhood.

\textbf{MANY HANDS MAKE LIGHT WORK!}

\tcblower
\textbf{Solution commentée.}

\textbf{Lexique de la leçon (8 mots minimum) :} \eng{community}, \eng{take part}, \eng{communal work}, \eng{clean up}, \eng{clear the gutters}, \eng{plant trees}, \eng{tools} (\eng{brooms, rakes, buckets, gloves}), \eng{volunteers}, \eng{neighbourhood}. L'objectif est largement atteint et chaque mot est employé dans son sens exact.

\textbf{Collocations :} \eng{to take part in} (et non \eng{to participate to}, calque du français « participer à ») ; \eng{to lend a hand} ; \eng{to work together} ; \eng{to clean up the streets}. Quatre collocations correctes, soit deux de plus que le minimum demandé.

\textbf{Expression idiomatique :} \eng{Many hands make light work!} est placée en conclusion, en majuscules, comme un slogan : c'est son emploi naturel en anglais.

\textbf{Structures :} impératifs d'invitation (\eng{Come and lend a hand}), futur d'annonce (\eng{We will clean up\ldots}), questions-réponses pour la lisibilité (\eng{When? Where? What?}), prépositions correctes (\eng{on Saturday}, \eng{at 7 a.m.}, \eng{in front of}).

\textbf{Faux amis évités :} le groupe n'utilise ni \eng{assist} ni \eng{actually} ; « quartier » est rendu par \eng{neighbourhood} et non par le faux calque \eng{quarter} employé seul (on garde \eng{Quarter} uniquement dans le nom propre de l'opération).

\textbf{Présentation orale modèle (1 minute) :} \eng{Good morning everybody! This is our poster for Operation Clean Quarter. We invite all the people of Biyem-Assi to take part in a big communal work day on Saturday 15th June, at 7 a.m., in front of the community hall. We will clean up the streets, clear the gutters and plant trees. Don't forget to bring your own tools: brooms, rakes and gloves. Come and lend a hand! Remember: many hands make light work. Thank you!}
\end{exoresolu}

\begin{corrige}[Tâche 1 — Compréhension du message]
\begin{enumerate}
\item \eng{The gutters are blocked and rubbish is everywhere in the streets.} (« Les caniveaux sont bouchés et il y a des ordures partout dans les rues. »)
\item \eng{It will take place on Saturday 15th June, at 7 a.m., in front of the community hall.} (« Elle aura lieu le samedi 15 juin, à 7 heures, devant la salle communautaire. »)
\item \eng{They will clean up the streets, clear the gutters and plant trees near the primary school.} (« Ils nettoieront les rues, déboucheront les caniveaux et planteront des arbres près de l'école primaire. »)
\item \eng{They must bring their own tools: machetes, rakes, brooms, buckets and gloves.} (« Ils doivent apporter leurs propres outils : machettes, râteaux, balais, seaux et gants. »)
\end{enumerate}
Remarque de méthode : chaque réponse reprend les termes de la question et forme une phrase complète avec sujet et verbe, comme l'exigent les épreuves de compréhension de Première.
\end{corrige}

\courssub{Bilan de l'activité}

Cette activité d'intégration a permis de mobiliser l'ensemble des savoir-faire de la leçon dans une tâche finale authentique et motivante :

\begin{itemize}
\item \textbf{Compréhension de l'écrit :} le message du chef du quartier a servi de point de départ et a guidé le contenu de l'affiche.
\item \textbf{Lexique :} les élèves ont sélectionné, orthographié et réemployé en contexte au moins huit mots du champ lexical des travaux communautaires.
\item \textbf{Collocations et idiomes :} les contraintes de production ont imposé l'emploi actif de collocations (\eng{to take part in}, \eng{to lend a hand}) et de l'expression \eng{Many hands make light work!}.
\item \textbf{Vigilance linguistique :} la grille d'évaluation a sanctionné les faux amis et les calques du français, renforçant les mises en garde de la leçon.
\item \textbf{Expression orale :} la présentation d'une minute a entraîné la prise de parole en public, la prononciation du lexique nouveau et l'usage d'impératifs d'invitation.
\item \textbf{Dimension citoyenne :} enfin, l'activité a valorisé l'esprit communautaire, réalité bien connue des élèves camerounais, et a montré que l'anglais est un outil concret de communication au service de la vie du quartier.
\end{itemize}

Pour prolonger l'activité à la maison, chaque élève peut rédiger dans son cahier un court paragraphe (six à huit lignes) intitulé \eng{Why I like communal work in my neighbourhood}, en réemployant au moins cinq mots du lexique et une collocation de la leçon.

\newpage

\coursec{Méthodes et exercices résolus}

\coursec{Lesson 6 — Vocabulary: Words and expressions related to taking part in communal work}

\courssub{Savoir-faire 1 : Identifier et employer le vocabulaire du travail communautaire}

\begin{methode}[Identifier le vocabulaire du thème dans un texte]
Pour repérer et mémoriser le lexique de \emph{communal work} dans un texte de compréhension :
\begin{enumerate}
\item \textbf{Lis d'abord le titre et le premier paragraphe} : ils annoncent le thème (\eng{community, village, neighbourhood, project}).
\item \textbf{Repère les mots-clés du champ lexical} : souligne les noms d'actions (\eng{clean-up, fund-raising, harvest, construction}), les noms de personnes (\eng{volunteers, participants, organisers, community leaders}) et les verbes d'action (\eng{to gather, to contribute, to join in, to take part in}).
\item \textbf{Classe les mots en colonnes} : personnes / actions / lieux / outils. Cette classification facilite la mémorisation.
\item \textbf{Devine le sens par le contexte} avant de chercher dans le glossaire : \eng{The villagers pooled their resources to build a bridge} → \emph{pool} = mettre en commun (contexte : ressources, construire).
\item \textbf{Note chaque mot nouveau avec sa nature et sa traduction} dans ton carnet de vocabulaire, avec une phrase d'exemple complète.
\end{enumerate}
Réflexe : un mot s'apprend toujours \textbf{dans une phrase}, jamais isolément.
\end{methode}

\begin{textebox}[Community work in Bafut — Reading text]
Every last Saturday of the month, the people of Bafut take part in communal work. Men, women and youths gather at the market square at 6 a.m. with hoes, machetes and wheelbarrows. They clean the gutters, sweep the streets and collect rubbish. The council provides gloves and bags, while local traders donate food and drinks for the volunteers. Thanks to this collective effort, the town has become cleaner and the risk of cholera has dropped.
\end{textebox}

\begin{exoresolu}[Reading comprehension — Community work in Bafut]
\textbf{Énoncé.} Read the text and answer the questions in English.

Questions :
\begin{enumerate}
\item When does communal work take place in Bafut?
\item Mention two tools the people bring.
\item Who provides gloves and bags?
\item What do local traders donate?
\item Give two benefits of this collective effort.
\end{enumerate}
\tcblower
\textbf{Solution détaillée.}

\textbf{Question 1.} Le texte dit : \eng{Every last Saturday of the month}. Réponse complète : \emph{Communal work takes place every last Saturday of the month.} — Attention à la préposition : on dit \eng{on the last Saturday} pour un samedi précis, mais ici la réponse peut reprendre la formulation du texte.

\textbf{Question 2.} Le texte énumère : \eng{hoes, machetes and wheelbarrows}. Deux suffisent : \emph{They bring hoes and machetes.} (houes et machettes). \eng{Wheelbarrow} = brouette, à retenir.

\textbf{Question 3.} \eng{The council provides gloves and bags.} → \emph{The council provides gloves and bags.} Ici \eng{council} = la municipalité / la commune (et non « conseil » au sens de conseil donné).

\textbf{Question 4.} \eng{Local traders donate food and drinks for the volunteers.} → \emph{They donate food and drinks for the volunteers.} Noter le verbe \eng{to donate} (faire don de) et la préposition \eng{for}.

\textbf{Question 5.} Le dernier énoncé donne deux bénéfices : \emph{The town has become cleaner and the risk of cholera has dropped.} Noter l'expression \eng{thanks to this collective effort} (grâce à cet effort collectif) et le présent perfect \eng{has become / has dropped} qui présente un résultat actuel d'une action passée.

\textbf{Conclusion.} Toutes les réponses sont des phrases complètes construites à partir du texte, sans copier-coller mécanique : c'est ce qu'attend l'examinateur.
\end{exoresolu}

\begin{aretenir}[Vocabulaire de base — Communal work]
\begin{center}
\begin{tabularx}{\linewidth}{|X|X|X|}
\hline
\rowcolor{popblueL} \textbf{Personnes} & \textbf{Actions / Verbes} & \textbf{Lieux / Outils / Noms} \\
\hline
\eng{volunteer} (bénévole) & \eng{to take part in} (participer à) & \eng{community} (communauté) \\
\eng{participant} (participant) & \eng{to participate in} (participer à) & \eng{neighbourhood} (quartier) \\
\eng{organiser} (organisateur) & \eng{to volunteer} (faire du bénévolat) & \eng{market square} (place du marché) \\
\eng{community leader} (chef de communauté) & \eng{to gather} (se rassembler) & \eng{hoe} (houe) \\
\eng{local trader} (commerçant local) & \eng{to contribute (to)} (contribuer à) & \eng{machete} (machette) \\
\eng{the council} (la municipalité) & \eng{to join in} (se joindre à) & \eng{wheelbarrow} (brouette) \\
& \eng{to lend a hand} (donner un coup de main) & \eng{gutters} (caniveaux) \\
& \eng{to pool resources} (mettre en commun) & \eng{rubbish / litter} (déchets) \\
& \eng{to clean up} (nettoyer) & \eng{gloves} (gants) \\
& \eng{to sweep} (balayer) & \eng{fund-raising} (collecte de fonds) \\
& \eng{to donate} (faire don de) & \eng{collective effort} (effort collectif) \\
\hline
\end{tabularx}
\end{center}
\end{aretenir}

\courssub{Savoir-faire 2 : Réemployer le lexique à l'oral et à l'écrit}

\begin{methode}[Réemployer le vocabulaire dans tes propres productions]
Pour utiliser activement le lexique du travail communautaire :
\begin{enumerate}
\item \textbf{Prépare une banque de 10 à 15 mots} du thème avant de parler ou d'écrire : \eng{to volunteer, to take part in, to lend a hand, to organise, to contribute, clean-up campaign, community project, team spirit, fund-raising, collective effort}.
\item \textbf{Structure ta production en trois temps} : présenter l'action (\eng{Our school organised a clean-up campaign}), décrire le déroulement (\eng{We gathered at 7 a.m., shared the tasks and worked together}), conclure sur le résultat (\eng{As a result, the campus is now clean and safe}).
\item \textbf{Varie les verbes} : évite de répéter \eng{to help} ; utilise \eng{to support, to assist, to contribute to, to join in, to take part in}.
\item \textbf{Utilise les connecteurs} : \eng{first, then, moreover, as a result, thanks to}.
\item \textbf{À l'oral, projette la voix et soigne la prononciation} des mots-clés : \eng{volunteer} « vo-lon-ti-euh », \eng{communal} « ko-miou-nal », \eng{contribute} « kon-tri-bioute », \eng{neighbourhood} « neï-beu-ho-d ».
\end{enumerate}
\end{methode}

\begin{exoresolu}[Guided writing — Describe a communal activity]
\textbf{Énoncé.} Your class took part in a tree-planting day at your school. Write a short paragraph (60–80 words) describing the event. Use at least five words from the lesson: \emph{take part in, volunteers, to organise, tools, team spirit, as a result}.

\tcblower
\textbf{Solution détaillée.}

\textbf{Étape 1 — Plan.} Introduction (quoi, quand, où) → déroulement (qui, tâches, outils) → résultat et sentiment.

\textbf{Étape 2 — Choix du lexique.} \eng{take part in, volunteers, organised, spades and watering cans, team spirit, as a result}.

\textbf{Étape 3 — Rédaction.}

\emph{Last Friday, our class took part in a tree-planting day at Government High School Mfandena. The event was organised by the environmental club. Twenty volunteers gathered at 7 a.m. with spades and watering cans. We dug holes, planted fifty young trees and watered them together. There was a real team spirit: everybody lent a hand. As a result, our school yard is greener and shadier, and we are proud of our collective effort.}

\textbf{Étape 4 — Vérification.} 75 mots environ ; six mots-cibles employés ; passé simple pour le récit (\eng{took part, gathered, dug, planted}) ; connecteur de conséquence \eng{as a result} ; aucun calque du français (on écrit \eng{took part in} et non *\eng{participated to}).

\textbf{Conclusion.} Une bonne production courte = un plan simple, du lexique précis, des temps corrects et un connecteur de conclusion.
\end{exoresolu}

\begin{experience}[Speaking activity — Organise a community project]
\textbf{Consigne.} En binôme, imaginez un projet communautaire pour votre quartier ou votre école (ex. : \eng{clean-up campaign, painting the classroom, creating a garden, repairing a road}).
\begin{enumerate}
\item \textbf{Student A} : Vous êtes le chef de quartier / le président du club. Présentez le projet : but, date, matériel nécessaire, appel aux bénévoles. Utilisez : \eng{We are organising... We need volunteers to... Please bring...}
\item \textbf{Student B} : Vous êtes un habitant / un élève. Posez des questions : \eng{When does it take place? What tools do we need? Can I bring...?} Acceptez de participer : \eng{I'll take part in it. I'll lend a hand.}
\item \textbf{Échangez les rôles} et changez de projet.
\end{enumerate}
\textbf{Évaluation par les pairs :} Ai-je utilisé \eng{take part in, lend a hand, volunteers, organise, tools, community} ? Ma prononciation de \eng{volunteer} et \eng{neighbourhood} était-elle correcte ?
\end{experience}

\courssub{Savoir-faire 3 : Éviter les faux amis et les calques du français}

\begin{methode}[Déjouer les faux amis et les calques]
\begin{enumerate}
\item \textbf{Méfie-toi des mots qui ressemblent au français.} \eng{Actually} = en fait (et non « actuellement » → \eng{currently}) ; \eng{to assist} = aider, mais \eng{to assist at} n'existe pas pour « assister à » → \eng{to attend}.
\item \textbf{Vérifie la construction du verbe.} \eng{To participate} se construit avec \eng{in} : \eng{to participate in communal work} (jamais *\eng{participate to}). \eng{To help} est transitif direct : \eng{to help the community} (jamais *\eng{help to the community}). \eng{To help someone (to) do something} (aider quelqu'un à faire qqch).
\item \textbf{Attention aux adjectifs.} \eng{A community} = une communauté ; « communal » (qui appartient à la communauté) se dit \eng{communal} (\eng{communal work}) mais « commun » au sens d'ordinaire/fréquent se dit \eng{common}.
\item \textbf{Traduis mentalement ta phrase dans les deux sens.} Si la retraduction donne une phrase anglaise bizarre, c'est un calque : « prendre part à » → \eng{to take part in} (correct), mais « faire un effort » → \eng{to make an effort} (et non *\eng{do an effort}).
\item \textbf{Tiens une liste personnelle de faux amis} rencontrés dans les textes et relis-la avant chaque évaluation.
\end{enumerate}
\end{methode}

\begin{exoresolu}[Error correction — Find and correct the mistakes]
\textbf{Énoncé.} Each sentence contains one mistake (false friend, wrong preposition or calque from French). Correct it.

\begin{enumerate}
\item Many youths participated to the clean-up campaign last month.
\item The volunteers assisted to the meeting organised by the chief.
\item We must make an effort to keep our quarter clean — actually it is very dirty.
\item The population helped to the workers to build the bridge.
\item Communal work is very commun in Cameroonian villages.
\end{enumerate}
\tcblower
\textbf{Solution détaillée.}

\textbf{1.} *\eng{participated to} → \eng{participated in}. Le verbe \eng{to participate} exige la préposition \eng{in}. Phrase correcte : \emph{Many youths participated in the clean-up campaign last month.}

\textbf{2.} Faux ami : « assister à » (être présent) ne se traduit pas par \eng{to assist to}. \eng{To assist} = aider. Correction : \emph{The volunteers attended the meeting organised by the chief.}

\textbf{3.} Deux erreurs. Faux ami : \eng{actually} signifie « en fait », pas « actuellement » → \eng{currently} (ou \eng{at present}). Calque : « notre quartier » ne se dit pas *\eng{our quarter} (qui signifie « notre trimestre » ou « notre quartier [historique/militaire] ») mais \eng{our neighbourhood} ou \eng{our area}. Correction : \emph{We must make an effort to keep our neighbourhood clean — currently it is very dirty.}

\textbf{4.} *\eng{helped to the workers} : \eng{to help} est transitif direct. Correction : \emph{The population helped the workers (to) build the bridge.} (Le \eng{to} après \eng{help} est optionnel).

\textbf{5.} *\eng{commun} : orthographe et sens. « Commun » (fréquent, ordinaire) = \eng{common} ; \eng{communal} = collectif, de la communauté. Correction : \emph{Communal work is very common in Cameroonian villages.} — Noter que la phrase joue justement sur les deux adjectifs.

\textbf{Conclusion.} À l'examen, relis chaque phrase en te demandant : « Ce mot existe-t-il vraiment en anglais avec ce sens et cette construction ? »
\end{exoresolu}

\courssub{Savoir-faire 4 : Employer collocations et expressions idiomatiques en contexte}

\begin{methode}[Utiliser les collocations et expressions idiomatiques]
\begin{enumerate}
\item \textbf{Apprends les mots par paires fixes} (collocations), pas isolément : \eng{to take part in, to lend a hand, to carry out a project, to raise funds, to make a contribution, to join forces, community spirit, collective effort}.
\item \textbf{Respecte le verbe de la collocation} : on dit \eng{to raise funds} (et non *\eng{rise funds} — \eng{rise} est intransitif), \eng{to make a contribution} (et non *\eng{do a contribution}), \eng{to carry out a project} (et non *\eng{make a project}).
\item \textbf{Pour les expressions idiomatiques}, retiens le sens global et un exemple : \eng{to lend a hand} = donner un coup de main ; \eng{many hands make light work} = à plusieurs, le travail est plus facile ; \eng{to pull together} = se serrer les coudes / unir ses forces.
\item \textbf{Place la collocation dans une phrase complète} pour vérifier qu'elle sonne naturel : \eng{The whole village pulled together to repair the school roof.}
\item \textbf{À l'écrit, une ou deux idiomatiques suffisent} : trop en mettre rend le texte artificiel.
\end{enumerate}
\end{methode}

\begin{propriete}[Tableau des collocations clés du thème]
\begin{center}
\begin{tabularx}{\linewidth}{|l|X|X|}
\hline
\rowcolor{popblueL} \textbf{Collocation / Idiome} & \textbf{Sens français} & \textbf{Exemple} \\
\hline
\eng{to take part in} & participer à & \eng{All students took part in the clean-up.} \\
\eng{to participate in} & participer à (plus formel) & \eng{They participated in the fund-raising.} \\
\eng{to lend a hand} & donner un coup de main & \eng{Can you lend a hand with the painting?} \\
\eng{to raise funds} & collecter des fonds & \eng{We raised funds for the new well.} \\
\eng{to carry out} & mener à bien, réaliser & \eng{They carried out the project successfully.} \\
\eng{to make a contribution} & apporter une contribution & \eng{Every family made a contribution.} \\
\eng{to join forces} & unir ses forces & \eng{The two villages joined forces.} \\
\eng{community spirit} & esprit communautaire & \eng{The event showed great community spirit.} \\
\eng{collective effort} & effort collectif & \eng{Thanks to our collective effort, we finished early.} \\
\eng{many hands make light work} & à plusieurs, le travail est facile (proverbe) & \eng{Let's go! Many hands make light work.} \\
\hline
\end{tabularx}
\end{center}
\end{propriete}

\begin{exoresolu}[Collocations — Complete the sentences]
\textbf{Énoncé.} Complete each sentence with the correct word from the box: \emph{raise — lend — carry — take — make}.

\begin{enumerate}
\item The youth association wants to \ldots\ funds to buy dustbins for the market.
\item During the harvest, neighbours always \ldots\ a hand on each other's farms.
\item Our school will \ldots\ out a clean-up project next term.
\item All the pupils \ldots\ part in the sports day organised by the council.
\item Every family should \ldots\ a contribution to the village water project.
\end{enumerate}
\tcblower
\textbf{Solution détaillée.}

\textbf{1.} \eng{raise} → \emph{The youth association wants to raise funds to buy dustbins for the market.} Collocation : \eng{to raise funds} = collecter des fonds. Piège : *\eng{rise} est intransitif (le soleil se lève).

\textbf{2.} \eng{lend} → \emph{During the harvest, neighbours always lend a hand on each other's farms.} Expression idiomatique : \eng{to lend a hand} = donner un coup de main. Ne pas confondre avec \eng{to borrow} (emprunter) : \eng{lend} = prêter.

\textbf{3.} \eng{carry} → \emph{Our school will carry out a clean-up project next term.} Phrasal verb : \eng{to carry out} = mener à bien, réaliser. On ne dit pas *\eng{make a project} dans ce sens.

\textbf{4.} \eng{take} → \emph{All the pupils take part in the sports day organised by the council.} Collocation : \eng{to take part in} = participer à. Ne jamais oublier la préposition \eng{in}.

\textbf{5.} \eng{make} → \emph{Every family should make a contribution to the village water project.} Collocation : \eng{to make a contribution} = apporter une contribution. Piège du calque : « faire une contribution » existe en français, mais ici le verbe \eng{make} est bien le bon en anglais ; en revanche on dit \eng{to make an effort}, pas *\eng{do an effort}.

\textbf{Conclusion.} Les collocations s'acquièrent par l'usage : rédige une phrase personnelle avec chacune d'elles pour les mémoriser durablement.
\end{exoresolu}

\courssub{Savoir-faire 5 : Exploiter les familles de mots, la prononciation et l'orthographe}

\begin{methode}[Construire et utiliser les familles de mots]
\begin{enumerate}
\item \textbf{À partir d'un mot de base, forme la famille} : \eng{community} (nom) → \eng{communal} (adjectif) ; \eng{to volunteer} (verbe) → \eng{volunteer} (nom de personne) / \eng{volunteering} (nom d'action) → \eng{voluntary} (adjectif) ; \eng{to organise} (verbe) → \eng{organisation} (nom d'action) → \eng{organiser} (nom de personne) → \eng{organised} (adjectif).
\item \textbf{Repère les suffixes typiques} : \emph{-tion / -sion} (action : \eng{contribution, organisation, construction, participation}), \emph{-er / -or} (personne : \eng{organiser, contributor, participant}), \emph{-al / -ive / -ary} (adjectif : \eng{communal, collective, voluntary}).
\item \textbf{Soigne l'orthographe} : \eng{community} (un seul m, un seul t), \eng{neighbourhood} (avec \emph{-gh-} muet), \eng{volunteer} (double e final), \eng{organisation} (avec s en graphie britannique ; \eng{z} en américain — choisis l'une et reste cohérent).
\item \textbf{Travaille la prononciation} : \eng{community} « ko-miou-ni-ti » (accent sur \emph{ni}), \eng{volunteer} « vo-lon-ti-euh » (accent sur la dernière syllabe), \eng{organise} « o-ga-naïze », \eng{contribute} « kon-tri-bioute », \eng{neighbourhood} « neï-beu-ho-d ».
\item \textbf{En production écrite, varie les formes} pour éviter les répétitions : \eng{The organisation of the event was perfect: the organisers had planned everything.}
\end{enumerate}
\end{methode}

\begin{exoresolu}[Word formation — Complete the table and the sentences]
\textbf{Énoncé A.} Complete the table with the missing forms. Distingue le nom d'action et le nom de personne.

\begin{center}
\begin{tabular}{|l|l|l|l|}
\hline
\rowcolor{popblueL} \textbf{Verb} & \textbf{Noun (action)} & \textbf{Noun (person)} & \textbf{Adjective} \\
\hline
to organise & \ldots & \ldots & organised \\
\hline
to contribute & \ldots & \ldots & contributive \\
\hline
\ldots & community & \ldots & communal \\
\hline
to volunteer & \ldots & \ldots & voluntary \\
\hline
\end{tabular}
\end{center}

\textbf{Énoncé B.} Complete the sentences with the correct form of the word in brackets.

\begin{enumerate}
\item The \ldots\ of the clean-up day took two weeks. (organise)
\item His \ldots\ to the project was highly appreciated. (contribute)
\item The event was entirely \ldots\ : nobody was paid. (volunteer)
\item She is a very active \ldots\ in the local association. (participate)
\end{enumerate}
\tcblower
\textbf{Solution détaillée.}

\textbf{Partie A.}
\begin{center}
\begin{tabular}{|l|l|l|l|}
\hline
\rowcolor{popblueL} \textbf{Verb} & \textbf{Noun (action)} & \textbf{Noun (person)} & \textbf{Adjective} \\
\hline
to organise & \textbf{organisation} & \textbf{organiser} & organised \\
\hline
to contribute & \textbf{contribution} & \textbf{contributor} & contributive \\
\hline
\textbf{— (pas de verbe courant)} & community & \textbf{community member} & communal \\
\hline
to volunteer & \textbf{volunteering / voluntary work} & \textbf{volunteer} & voluntary \\
\hline
\end{tabular}
\end{center}
Remarques : \eng{community} n'a pas de verbe dérivé courant ; on utilise la périphrase \eng{to build a community}. Pour \eng{volunteer}, le nom d'action est \eng{volunteering} (ou \eng{voluntary work}) ; \eng{volunteer} désigne la personne. \eng{Participant} vient de \eng{to participate}.

\textbf{Partie B.}

\textbf{1.} \emph{The \textbf{organisation} of the clean-up day took two weeks.} — Après l'article \eng{the}, il faut un nom ; suffixe \emph{-tion}. Orthographe : \eng{organisation} (graphie britannique ; \eng{organization} est la graphie américaine, les deux sont admises mais ne pas mélanger dans un même texte).

\textbf{2.} \emph{His \textbf{contribution} to the project was highly appreciated.} — Après le possessif \eng{his}, il faut un nom. Noter la préposition \eng{to} : \eng{a contribution to a project}.

\textbf{3.} \emph{The event was entirely \textbf{voluntary}: nobody was paid.} — Après le verbe \eng{was} et l'adverbe \eng{entirely}, il faut un adjectif. Piège : \eng{volunteer} est un nom ou un verbe, jamais un adjectif.

\textbf{4.} \emph{She is a very active \textbf{participant} in the local association.} — Après l'article \eng{a} et l'adjectif \eng{active}, il faut un nom de personne. Suffixes \emph{-ant} / \emph{-ent}.

\textbf{Conclusion.} Pour choisir la bonne forme, demande-toi toujours : « Quelle est la fonction du mot dans la phrase ? » (nom après article ou possessif, adjectif après \eng{be}, verbe après le sujet).
\end{exoresolu}

\begin{definition}[Mots de la même famille — Récapitulatif]
\begin{itemize}
\item \cle{Community} (nom) : communauté. \eng{Community work} = travail communautaire. \eng{Community leader} = chef de communauté.
\item \cle{Communal} (adjectif) : communal, collectif (qui appartient à la communauté). \eng{Communal work, communal land, communal effort}. Ne pas confondre avec \eng{common} (courant, fréquent, partagé).
\item \cle{Volunteer} (verbe / nom personne) : faire du bénévolat / bénévole. \eng{To volunteer for a task}. \eng{A volunteer}.
\item \cle{Voluntary} (adjectif) : volontaire, bénévole. \eng{Voluntary work, a voluntary basis}.
\item \cle{Organise} (verbe) : organiser. \eng{Organisation} (nom action), \eng{Organiser} (nom personne), \eng{Organised} (adjectif).
\item \cle{Contribute} (verbe) : contribuer. \eng{Contribution} (nom action), \eng{Contributor} (nom personne), \eng{Contributive} (adjectif, rare), souvent remplacé par \eng{valuable contribution}.
\item \cle{Participate} (verbe) : participer. \eng{Participation} (nom action), \eng{Participant} (nom personne), \eng{Participatory} (adjectif).
\end{itemize}
\end{definition}

\begin{savaistu}[Culture \& Contexte : Le travail communautaire au Cameroun et ailleurs]
Au Cameroun, le travail communautaire prend des formes variées selon les régions :
\begin{itemize}
\item \textbf{Le \eng{Njangi} (ou \eng{Meeting}) :} association d'épargne et de crédit rotative très répandue (surtout chez les femmes). Les membres cotisent chaque semaine/mois et le pot est donné à tour de rôle à l'un d'eux. C'est une forme de \eng{collective effort} financier.
\item \textbf{Les \eng{Country Sundays} / \eng{Clean-up Saturdays} :} dans les zones anglophones (Nord-Ouest, Sud-Ouest), le dernier samedi du mois est souvent consacré au \eng{communal labour} : nettoyage des marchés, entretien des routes, construction d'écoles. Le \eng{Fon} (chef traditionnel) ou le \eng{Council} (municipalité) mobilise la population.
\item \textbf{Le \eng{Call to serve} :} dans certaines écoles et universités, les étudiants effectuent des \eng{community service hours} (heures de service communautaire) obligatoires (nettoyage du campus, plantation d'arbres).
\end{itemize}
Dans le monde anglophone : \eng{Barn raising} (Amish/USA — construction collective d'une grange), \eng{Working bee} (Australie/NZ — journée de travail bénévole), \eng{Volunteering} (Royaume-Uni/USA — engagement associatif). L'expression \eng{"Many hands make light work"} résume cet esprit universel.
\end{savaistu}

\begin{attention}[Les 5 erreurs qui coûtent des points au Bac]
\begin{enumerate}
\item \textbf{Mauvaise préposition avec les verbes de participation.} *\eng{participate to}, *\eng{take part to}, *\eng{help to somebody} sont fautifs. Correct : \eng{to participate in}, \eng{to take part in}, \eng{to help somebody} (sans préposition) ou \eng{help somebody (to) do something}.
\item \textbf{Faux amis.} \eng{Actually} ne veut pas dire « actuellement » (\eng{currently / at present}) ; \eng{to assist} ne veut pas dire « assister à » (\eng{to attend}) ; \eng{common} (ordinaire, fréquent) n'est pas \eng{communal} (collectif) ; \eng{quarter} (trimestre, quartier historique) n'est pas \eng{neighbourhood / area} (quartier d'habitation).
\item \textbf{Calques du français.} « Faire un effort » → \eng{to make an effort} (jamais *\eng{do an effort}) ; « donner un coup de main » → \eng{to lend a hand} (jamais *\eng{give a hand}) ; « travail communautaire » → \eng{communal work} ou \eng{community work} (jamais *\eng{communitarian work}) ; « actuellement » → \eng{currently} (pas \eng{actually}).
\item \textbf{Orthographe des mots-clés.} \eng{community} (un m, un t), \eng{neighbourhood} (\emph{-gh-} muet), \eng{volunteer} (deux e à la fin), \eng{organisation} (avec s en anglais britannique). Une faute sur le mot central du thème est très mal perçue.
\item \textbf{Temps verbaux dans le récit d'une action passée.} Pour raconter une journée de travail communautaire, utilise le prétérit simple (\eng{past simple}) : \eng{We gathered, cleaned and planted.} Ne mélange pas présent et prétérit dans le même récit, et n'oublie pas les irréguliers : \eng{take → took}, \eng{make → made}, \eng{build → built}, \eng{give → gave}, \eng{lend → lent}.
\end{enumerate}
\end{attention}

\begin{exercice}[Entraînement personnel — À rendre au professeur]
\textbf{Exercice 1 — Vocabulary.} Match each verb (1–5) with its correct collocation (a–e).
\begin{enumerate}
\item to raise \quad a. a hand
\item to lend \quad b. funds
\item to carry \quad c. part in
\item to take \quad d. out a project
\item to make \quad e. a contribution
\end{enumerate}

\textbf{Exercice 2 — Word Formation.} Complete the sentences with the correct form of the word in CAPITALS.
\begin{enumerate}
\item The \_\_\_\_\_\_\_\_\_\_ of the new bridge took six months. (CONSTRUCT)
\item Every \_\_\_\_\_\_\_\_\_\_ received a certificate of appreciation. (PARTICIPATE)
\item It was a \_\_\_\_\_\_\_\_\_\_ decision to clean the market. (COLLECT)
\item We need more \_\_\_\_\_\_\_\_\_\_ for the health campaign. (VOLUNTEER)
\end{enumerate}

\textbf{Exercice 3 — Translation / Writing.} Translate into English.
\begin{enumerate}
\item Les villageois ont uni leurs forces pour réparer le toit de l'école.
\item Chaque famille doit apporter une contribution au projet d'eau.
\item Actuellement, le quartier est sale, mais nous allons organiser une journée de nettoyage.
\item Les bénévoles ont aidé les ouvriers à construire le pont.
\item Grâce à cet effort collectif, le risque de choléra a diminué.
\end{enumerate}

\textbf{Exercice 4 — Production écrite.} Your school organised a \eng{"Clean-Up Day"} last Saturday. Write a report (100–120 words) for the school magazine describing the event. Use at least 8 words/expressions from the lesson. Structure: Title → Introduction → Activities → Result/Conclusion.
\end{exercice}

\begin{corrige}[Exercice 1]
1–b (\eng{raise funds}), 2–a (\eng{lend a hand}), 3–d (\eng{carry out a project}), 4–c (\eng{take part in}), 5–e (\eng{make a contribution}).
\end{corrige}

\begin{corrige}[Exercice 2]
1. \textbf{construction} (nom action, suffixe \emph{-tion}).
2. \textbf{participant} (nom personne, suffixe \emph{-ant}).
3. \textbf{collective} (adjectif, suffixe \emph{-ive} ; \eng{collective decision}).
4. \textbf{volunteers} (nom personne au pluriel ; \eng{volunteering} serait le nom d'action, mais ici "more" attend un nom comptable pluriel).
\end{corrige}

\begin{corrige}[Exercice 3]
1. \emph{The villagers joined forces to repair the school roof.} (ou \emph{pulled together to repair...})
2. \emph{Every family must make a contribution to the water project.}
3. \emph{Currently, the neighbourhood is dirty, but we are going to organise a clean-up day.} (Attention : \eng{currently} pas \eng{actually} ; \eng{neighbourhood} pas \eng{quarter} ; \eng{clean-up day}).
4. \emph{The volunteers helped the workers (to) build the bridge.} (Sans \eng{to} après \eng{help} possible).
5. \emph{Thanks to this collective effort, the risk of cholera has dropped.} (ou \eng{decreased}).
\end{corrige}

\begin{corrige}[Exercice 4 — Modèle de rédaction]
\textbf{A Successful Clean-Up Day at GHS Mfandena}

\emph{Last Saturday, our school organised a major clean-up day on the campus. The event was organised by the Environmental Club with the support of the Principal. Over two hundred volunteers, including students and teachers, took part in the activity. We gathered at 7 a.m. and shared the tasks: some swept the classrooms, others cleaned the gutters and collected rubbish. The council provided gloves and wheelbarrows. There was a wonderful community spirit; everyone lent a hand. As a result, the school environment is now clean and safe. Many hands make light work, and we are proud of our collective effort.}
\end{corrige}

\newpage

\coursec{Exercices}

\courssub{Je vérifie mes connaissances}

\begin{exercice}[QCM : le bon mot]
Choisis la bonne réponse (A, B ou C) pour compléter chaque phrase.

\begin{enumerate}
\item Every Saturday, young people in our quarter do \ldots\ work to clean the market.\\
A. volunteer \quad B. voluntary \quad C. voluntarily
\item The chief \ldots\ the villagers to take part in the tree-planting campaign.\\
A. mobilized \quad B. attended \quad C. demanded
\item Many people \ldots\ the funeral of the old teacher last Sunday.\\
A. assisted \quad B. attended \quad C. helped
\item The women \ldots\ money to repair the village well.\\
A. raised \quad B. rose \quad C. lifted
\item When my neighbour was building his house, the whole family \ldots\ him a hand.\\
A. gave \quad B. lent \quad C. borrowed
\end{enumerate}
\end{exercice}

\begin{exercice}[Vrai ou faux : attention aux faux amis]
Réponds par \textbf{True} ou \textbf{False} et justifie chaque réponse en une phrase (corrige la phrase si elle est fausse).

\begin{enumerate}
\item \eng{To assist a meeting means to attend it.}
\item \eng{“Voluntary work” is work that people do freely, without being paid.}
\item \eng{To demand politely means the same thing as to ask politely.}
\item \eng{“To participate to the meeting” is a correct English sentence.}
\item \eng{A “community” is a group of people living in the same area or sharing the same interests.}
\end{enumerate}
\end{exercice}

\begin{exercice}[Vocabulaire : trouve le mot]
Trouve dans la leçon le mot ou l'expression anglaise qui correspond à chaque définition.

\begin{enumerate}
\item A person who works without being paid, of his or her own free will.
\item A long open channel at the side of a road, used to carry away rainwater.
\item A small cart with one wheel, used to carry sand, stones or rubbish.
\item A day when all the inhabitants of a village or quarter work together to clean their environment.
\item Money collected from many people to pay for a common project.
\item The fact of working together with other people for a common goal.
\end{enumerate}
\end{exercice}

\begin{exercice}[Complète avec le mot de la même famille]
Complète chaque phrase avec la forme correcte du mot entre parenthèses (nom, verbe, adjectif ou adverbe).

\begin{enumerate}
\item The \ldots\ of the villagers made the project a success. (to participate)
\item Our school needs more \ldots\ to clean the classrooms. (voluntary)
\item The bridge was built with the \ldots\ of the whole community. (to help)
\item The headmaster \ldots\ thanked all the parents for their gifts. (gratitude)
\item They worked very \ldots\ together to finish the road before the rainy season. (hard / choose the right adverb form)
\end{enumerate}
\end{exercice}

\courssub{Je m'entraîne}

\begin{exercice}[Traduction : du français vers l'anglais]
Traduis les phrases suivantes en anglais en utilisant le lexique du thème, sans calquer le français.

\begin{enumerate}
\item Les habitants du quartier ont participé aux travaux communautaires samedi dernier.
\item Nous avons collecté des fonds pour construire un pont sur la rivière.
\item Tout le monde a donné un coup de main pour nettoyer le marché de Mokolo.
\item Le chef a mobilisé les jeunes pour réparer l'école du village.
\item Beaucoup de volontaires ont assisté à la réunion organisée par la mairie.
\end{enumerate}
\end{exercice}

\begin{exercice}[Texte à trous : a keep-clean day]
Complète le texte suivant avec les dix mots de la banque. Attention : chaque mot ne peut être utilisé qu'une seule fois.

\begin{center}
\fbox{\eng{gutters — wheelbarrow — raised funds — lent a hand — mobilized — rubbish — voluntary — swept — community — proud}}
\end{center}

\begin{textebox}[A keep-clean day in Buea]
Last Saturday, our quarter organized a keep-clean day. The quarter head had \ldots\ (1) all the young people a week before. At six o'clock in the morning, everybody was in the streets with brooms, cutlasses and spades. Some people cleaned the \ldots\ (2) so that rainwater could flow easily. Others picked up the \ldots\ (3) that people had thrown on the ground. An old man pushed a \ldots\ (4) full of sand to cover the mud near the market. Even the children \ldots\ (5) a hand: they \ldots\ (6) the front of the shops. This \ldots\ (7) work was not paid, but everybody was happy. Last month, the same \ldots\ (8) had \ldots\ (9) to buy new dustbins for the quarter. At the end of the day, we were tired but very \ldots\ (10) of our clean streets.
\end{textebox}
\end{exercice}

\begin{exercice}[Appariement : reconstitue les collocations]
Relie chaque début de collocation (colonne A) à sa fin (colonne B) pour former huit expressions correctes.

\begin{center}
\begin{tabularx}{\linewidth}{|X|X|}
\hline
\rowcolor{popblueL} \textbf{Column A} & \textbf{Column B} \\
\hline
1. to carry out & a. funds \\
\hline
2. to raise & b. a hand \\
\hline
3. to lend & c. a project \\
\hline
4. to take part & d. a meeting \\
\hline
5. to attend & e. in communal work \\
\hline
6. to clean up & f. the environment \\
\hline
7. to work & g. together \\
\hline
8. to contribute & h. to the development of the village \\
\hline
\end{tabularx}
\end{center}

Puis utilise quatre de ces collocations dans des phrases complètes de ton invention.
\end{exercice}

\begin{exercice}[Transformation de phrases]
Réécris chaque phrase en remplaçant le mot ou l'expression souligné(e) par un mot ou une expression du thème, sans changer le sens.

\begin{enumerate}
\item \eng{The villagers \underline{helped} the carpenter to build the classroom.} (use a phrasal expression with “hand”)
\item \eng{They \underline{collected} 200,000 francs to repair the well.} (use a verb + “funds”)
\item \eng{Everybody \underline{was present at} the meeting about the new market.} (use a more formal verb)
\item \eng{The students cleaned the school \underline{without being paid}.} (use an adverb)
\item \eng{The mayor \underline{asked the people to come together} for the tree-planting campaign.} (use a verb meaning “to organize and encourage people to act”)
\end{enumerate}
\end{exercice}

\courssub{Je me prépare au Bac}

\begin{exercice}[Compréhension de texte et questions de langue]
Lis attentivement le texte suivant, puis réponds aux questions.

\begin{textebox}[Communal work: a living tradition in Cameroon]
In many villages and towns of Cameroon, communal work is an old and respected tradition. Long before modern town councils existed, communities already worked together to build houses, clear farms, repair roads and dig wells. This spirit of solidarity is still alive today.

In Bamenda, for example, the inhabitants of each quarter meet once a month for a “clean-up campaign”. Men, women and children sweep the streets, cut the grass and empty the gutters. Nobody is paid: the work is entirely voluntary. “When we work together, we finish quickly and we become closer,” explains Mama Gladys, a shopkeeper who never misses a campaign.

Communal work is not only about cleaning. In the village of Bafut, the population recently raised funds to renovate the community hall. Every family contributed what it could: some gave money, others offered bags of cement, and the young men lent a hand during the construction. The village chief mobilized everybody, including the students on holiday.

According to community leaders, these activities teach young people important values such as cooperation, generosity and respect for common property. “A village that works together develops faster,” says the chief of Bafut. “Development does not always come from outside. Sometimes it starts with our own hands.”
\end{textebox}

\textbf{A. Comprehension questions} — Answer in complete sentences, in your own words.

\begin{enumerate}
\item Give two examples of communal work mentioned in the text.
\item What do the inhabitants of Bamenda do during the clean-up campaign? (two activities)
\item How did the population of Bafut contribute to the renovation of the community hall? (three ways)
\item According to the chief of Bafut, where does development sometimes start?
\item Find in the text a word or expression that means:
\begin{enumerate}
\item work done freely, without payment (paragraph 2);
\item to encourage and organize people to act (paragraph 3);
\item things that belong to everybody (paragraph 4).
\end{enumerate}
\end{enumerate}

\textbf{B. Language work}

\begin{enumerate}
\item Complete: \eng{The villagers raised funds \ldots\ renovate the hall.} (one word)
\item Put into the plural: \eng{The child swept the street with a brush.}
\item Turn into the passive voice: \eng{The young men repaired the village well.}
\item Choose the correct word: \eng{Many people (assisted / attended) the funeral and (raised / rose) money for the family.}
\end{enumerate}
\end{exercice}

\begin{exercice}[Traduction]
Traduis en anglais le passage suivant, en utilisant le lexique du thème et en évitant les calques du français.

\begin{textebox}[Texte à traduire]
Dans mon village, les travaux communautaires ont lieu le dernier samedi de chaque mois. Tous les habitants participent à ces travaux : les hommes réparent les routes, les femmes nettoient le marché et les enfants ramassent les ordures. L'année dernière, nous avons collecté des fonds pour construire un puits. Grâce à la solidarité de tous, notre village devient de plus en plus propre et beau.
\end{textebox}
\end{exercice}

\begin{exercice}[Composition : sujet type Baccalauréat]
\textbf{Topic:} \eng{Describe a communal work activity in your village or quarter and explain why such activities are important for the community.}

Rédige une composition de \textbf{80 à 100 mots}. Tu organiseras ta production en trois paragraphes :

\begin{enumerate}
\item \textbf{Introduction :} présente l'activité (quoi, où, quand, qui).
\item \textbf{Développement :} décris ce que chacun a fait et comment la communauté s'est organisée (utilise au moins quatre mots ou expressions du thème : \eng{voluntary, to raise funds, to lend a hand, to mobilize, to take part in, to carry out}\ldots).
\item \textbf{Conclusion :} explique pourquoi les travaux communautaires sont importants (solidarité, développement, propreté, valeurs transmises aux jeunes).
\end{enumerate}

\textbf{Conseils :} utilise le prétérit pour la description, des connecteurs (\eng{first, then, finally, moreover, that is why}) et vérifie l'accord sujet-verbe et l'emploi des prépositions.
\end{exercice}

\newpage

\coursec{Corrigés des exercices}

\begin{corrige}[Exercice 1 — QCM : le bon mot]
\begin{enumerate}
\item \textbf{B. voluntary} — \eng{Every Saturday, young people in our quarter do voluntary work to clean the market.} \emph{Voluntary} est l'adjectif qui qualifie le nom \emph{work}. \emph{Volunteer} est un nom (un volontaire) ou un verbe ; \emph{voluntarily} est un adverbe. \cle{Adjectif vs Nom/Verbe}
\item \textbf{A. mobilized} — \eng{The chief mobilized the villagers to take part in the tree-planting campaign.} \emph{To mobilize} = rassembler et encourager les gens à agir. \emph{Attended} (assister à) et \emph{demanded} (exiger) ne conviennent pas au sens. \cle{Verbe d'action}
\item \textbf{B. attended} — \eng{Many people attended the funeral of the old teacher last Sunday.} Faux ami : \emph{to attend} = assister à (être présent) ; \emph{to assist} = aider (sauf dans la tournure formelle \emph{assist at}). \cle{Faux ami : attend / assist}
\item \textbf{A. raised} — \eng{The women raised money to repair the village well.} \emph{To raise} (raised, raised) est transitif : lever/collecter des fonds. \emph{Rose} est le prétérit de \emph{to rise} (s'élever, intransitif) ; \emph{lifted} = soulever physiquement. \cle{Transitivité : raise / rise}
\item \textbf{B. lent} — \eng{When my neighbour was building his house, the whole family lent him a hand.} Expression figée : \emph{to lend somebody a hand} = donner un coup de main. \emph{Borrowed} signifie emprunter (sens inverse). \cle{Expression idiomatique}
\end{enumerate}
\end{corrige}

\begin{corrige}[Exercice 2 — Vrai ou faux : attention aux faux amis]
\begin{enumerate}
\item \textbf{False.} \eng{To assist} means to help, not to attend. Correction : \eng{To attend a meeting means to be present at it.} (Faux ami avec le français « assister à ».) \cle{Faux ami}
\item \textbf{True.} \eng{Voluntary work is done freely, without payment.} C'est bien la définition du travail volontaire/bénévole. \cle{Définition}
\item \textbf{False.} \eng{To demand} is strong and means to claim firmly; \eng{to ask} is the polite, neutral verb. Correction : \eng{To ask politely is not the same as to demand.} (Faux ami avec « demander ».) \cle{Nuance : ask vs demand}
\item \textbf{False.} \eng{To participate} is followed by the preposition \emph{in}, never by \emph{to}. Correction : \eng{To participate in the meeting} (ou mieux : \eng{to take part in the meeting}). \cle{Préposition : participate in}
\item \textbf{True.} \eng{A community is a group of people living in the same area or sharing the same interests.} C'est la définition correcte de \emph{community}. \cle{Définition}
\end{enumerate}
\end{corrige}

\begin{corrige}[Exercice 3 — Vocabulaire : trouve le mot]
\begin{enumerate}
\item \eng{A volunteer} — un volontaire / bénévole. \cle{Nom d'agent}
\item \eng{A gutter} — un caniveau (prononciation : « gueu-teur »). \cle{Prononciation}
\item \eng{A wheelbarrow} — une brouette (prononciation : « ouil-ba-ro »). \cle{Prononciation}
\item \eng{A keep-clean day} (ou \eng{a clean-up campaign}) — une journée de nettoyage communautaire. \cle{Expression composée}
\item \eng{Funds} (money raised) — des fonds collectés ; on parle aussi de \eng{a fund-raising}. \cle{Indénombrable / Nom composé}
\item \eng{Cooperation} (ou \eng{working together}) — la coopération, le travail en commun. \cle{Nom abstrait}
\end{enumerate}
\end{corrige}

\begin{corrige}[Exercice 4 — Complète avec le mot de la même famille]
\begin{enumerate}
\item \eng{The \textbf{participation} of the villagers made the project a success.} — nom formé sur le verbe \emph{to participate} (participate $\rightarrow$ participation). \cle{Dérivation : -tion}
\item \eng{Our school needs more \textbf{volunteers} to clean the classrooms.} — nom de personne, au pluriel (accord avec \emph{more}). \cle{Pluriel / Dérivation}
\item \eng{The bridge was built with the \textbf{help} of the whole community.} — \emph{help} fonctionne ici comme nom (invariable), précédé de l'article \emph{the}. \cle{Nom indénombrable}
\item \eng{The headmaster \textbf{gratefully} thanked all the parents for their gifts.} — adverbe en \emph{-ly} formé sur l'adjectif \emph{grateful} (attention à l'orthographe : \emph{grateful}, pas « greatful »). \cle{Orthographe : grateful / Adverbe en -ly}
\item \eng{They worked very \textbf{hard} together to finish the road before the rainy season.} — piège classique : l'adverbe de \emph{hard} est \emph{hard} ; \emph{hardly} signifie « à peine ». \cle{Adverbe irrégulier : hard / hardly}
\end{enumerate}
\end{corrige}

\begin{corrige}[Exercice 5 — Traduction : du français vers l'anglais]
\begin{enumerate}
\item \eng{The inhabitants of the quarter took part in communal work last Saturday.} — \emph{to take part in} + nom, sans « to » ; \emph{last Saturday} se place en fin de phrase. \cle{Expression : take part in}
\item \eng{We raised funds to build a bridge over the river.} — \emph{to raise funds} = collecter des fonds ; \emph{over} (et non « on ») pour « sur la rivière ». \cle{Préposition : over / Collocation : raise funds}
\item \eng{Everybody lent a hand to clean the Mokolo market.} — expression idiomatique \emph{to lend a hand} ; ne pas traduire mot à mot « gave a hand kick ». \cle{Idiome : lend a hand}
\item \eng{The chief mobilized the young people to repair the village school.} — \emph{to mobilize somebody to do something}. \cle{Structure : mobilize + OBJ + to INF}
\item \eng{Many volunteers attended the meeting organized by the town council.} — faux ami : « assister à » = \emph{to attend}, pas \emph{to assist} ; « la mairie » = \emph{the town council} (ou \emph{the council}). \cle{Faux ami : attend / Vocabulaire : town council}
\end{enumerate}
\end{corrige}

\begin{corrige}[Exercice 6 — Texte à trous : a keep-clean day]
\begin{enumerate}
\item \textbf{mobilized} — \eng{The quarter head had mobilized all the young people a week before.} (plus-que-parfait : had + participe passé) \cle{Temps : Past Perfect}
\item \textbf{gutters} — \eng{Some people cleaned the gutters so that rainwater could flow easily.} \cle{Pluriel}
\item \textbf{rubbish} — \eng{Others picked up the rubbish that people had thrown on the ground.} (\emph{rubbish} est \cle{indénombrable}) \cle{Indénombrable}
\item \textbf{wheelbarrow} — \eng{An old man pushed a wheelbarrow full of sand.} \cle{Singulier après 'a'}
\item \textbf{lent} — \eng{Even the children lent a hand.} (prétérit de \emph{to lend} : lent, lent) \cle{Verbe irrégulier : lend / lent / lent}
\item \textbf{swept} — \eng{they swept the front of the shops.} (prétérit de \emph{to sweep} : swept, swept) \cle{Verbe irrégulier : sweep / swept / swept}
\item \textbf{voluntary} — \eng{This voluntary work was not paid.} \cle{Adjectif}
\item \textbf{community} — \eng{the same community had raised funds.} \cle{Nom}
\item \textbf{raised funds} — \eng{had raised funds to buy new dustbins.} \cle{Collocation + Past Perfect}
\item \textbf{proud} — \eng{we were tired but very proud of our clean streets.} (\emph{to be proud of} = être fier de) \cle{Adjectif + préposition : proud of}
\end{enumerate}
\end{corrige}

\begin{corrige}[Exercice 7 — Appariement : reconstitue les collocations]
Appariements corrects :
\begin{center}
\begin{tabularx}{\linewidth}{|X|X|}
\hline
\rowcolor{popblueL} \textbf{Collocation} & \textbf{Traduction} \\
\hline
1. to carry out \textbf{+ c. a project} & réaliser un projet \\
\hline
2. to raise \textbf{+ a. funds} & collecter des fonds \\
\hline
3. to lend \textbf{+ b. a hand} & donner un coup de main \\
\hline
4. to take part \textbf{+ e. in communal work} & participer aux travaux communautaires \\
\hline
5. to attend \textbf{+ d. a meeting} & assister à une réunion \\
\hline
6. to clean up \textbf{+ f. the environment} & nettoyer l'environnement \\
\hline
7. to work \textbf{+ g. together} & travailler ensemble \\
\hline
8. to contribute \textbf{+ h. to the development of the village} & contribuer au développement du village \\
\hline
\end{tabularx}
\end{center}

Exemples de phrases (corrections possibles) :
\begin{itemize}
\item \eng{The villagers carried out a project to build a new bridge.} « Les villageois ont réalisé un projet de construction d'un nouveau pont. »
\item \eng{We raised funds to buy books for the school library.} « Nous avons collecté des fonds pour acheter des livres pour la bibliothèque scolaire. »
\item \eng{Everybody took part in communal work last Saturday.} « Tout le monde a participé aux travaux communautaires samedi dernier. »
\item \eng{If we work together, we will finish before noon.} « Si nous travaillons ensemble, nous finirons avant midi. »
\end{itemize}
\end{corrige}

\begin{corrige}[Exercice 8 — Transformation de phrases]
\begin{enumerate}
\item \eng{The villagers \textbf{lent a hand to} the carpenter to build the classroom.} (ou \eng{lent the carpenter a hand}) — expression figée \emph{to lend somebody a hand}. \cle{Structure : lend sb a hand / lend a hand to sb}
\item \eng{They \textbf{raised funds of} 200,000 francs to repair the well.} (ou \eng{They raised 200,000 francs in funds to repair the well.}) — \emph{to raise funds}, prétérit \emph{raised}. \cle{Collocation + Prétérit}
\item \eng{Everybody \textbf{attended} the meeting about the new market.} — verbe formel, transitif direct, sans préposition (jamais « attended to » ni « attended at »). \cle{Transitivité directe : attend}
\item \eng{The students cleaned the school \textbf{voluntarily}.} — adverbe en \emph{-ily} formé sur l'adjectif \emph{voluntary} (changement y $\rightarrow$ i). \cle{Dérivation adverbiale : -ily}
\item \eng{The mayor \textbf{mobilized} the people for the tree-planting campaign.} — \emph{to mobilize} = rassembler et encourager à agir. \cle{Verbe + préposition : mobilize for}
\end{enumerate}
\end{corrige}

\begin{corrige}[Exercice 9 — Compréhension de texte et questions de langue]
\textbf{A. Comprehension questions} (réponses modèles)
\begin{enumerate}
\item \eng{The text mentions sweeping the streets and cutting the grass during clean-up campaigns; it also mentions building houses, repairing roads, digging wells and renovating the community hall.} (deux exemples suffisent)
\item \eng{During the clean-up campaign in Bamenda, the inhabitants sweep the streets, cut the grass and empty the gutters.} (deux activités suffisent)
\item \eng{The population of Bafut contributed in three ways: some families gave money, others offered bags of cement, and the young men lent a hand during the construction.}
\item \eng{According to the chief of Bafut, development sometimes starts with our own hands, that is to say with the efforts of the community itself.}
\item \begin{enumerate}
\item \eng{voluntary} (the work is entirely \emph{voluntary});
\item \eng{mobilized} (the village chief \emph{mobilized} everybody);
\item \eng{common property} (respect for \emph{common property}).
\end{enumerate}
\end{enumerate}

\textbf{B. Language work}
\begin{enumerate}
\item \eng{The villagers raised funds \textbf{to} renovate the hall.} — \emph{to} + base verbale exprime le \cle{but} (« pour »).
\item \eng{The children swept the streets with brushes.} — pluriel irrégulier \emph{child} $\rightarrow$ \emph{children} ; \emph{brush} $\rightarrow$ \emph{brushes} (noms en \emph{-sh} : + \emph{es}). \cle{Pluriels irréguliers / en -sh}
\item \eng{The village well was repaired by the young men.} — \cle{passif} au prétérit : sujet + \emph{was/were} + participe passé ; \emph{well} est singulier, donc \emph{was}.
\item \eng{Many people \textbf{attended} the funeral and \textbf{raised} money for the family.} — \emph{attend} = assister à ; \emph{raise} (transitif) = collecter, alors que \emph{rise} (rose, risen) est \cle{intransitif}.
\end{enumerate}
\textbf{Barème indicatif (type Bac) :} compréhension 10 pts (2 pts par question de 1 à 4, 2 pts pour le vocabulaire) ; langue 10 pts (2,5 pts par item).
\end{corrige}

\begin{corrige}[Exercice 10 — Traduction]
Proposition de traduction :

\eng{In my village, communal work takes place on the last Saturday of every month. All the inhabitants take part in it: the men repair the roads, the women clean the market and the children pick up the rubbish. Last year, we raised funds to build a well. Thanks to everybody's solidarity, our village is becoming cleaner and more beautiful.}

Points de vigilance :
\begin{itemize}
\item « avoir lieu » = \emph{to take place} (jamais « to have place ») ; \cle{Verbe : take place}
\item « participer à » = \emph{to take part in} ou \emph{to participate in} (préposition \emph{in}) ; \cle{Préposition : in}
\item « les ordures » = \emph{rubbish}, \cle{indénombrable} (pas de \emph{-s}) ;
\item « collecter des fonds » = \emph{to raise funds} ; \cle{Collocation}
\item « de plus en plus propre » = \emph{cleaner and cleaner} (\cle{comparatif de progression} répété) ; « plus belle » = \emph{more beautiful} (comparatif simple) ou \emph{more and more beautiful} (progression) ;
\item « grâce à » = \emph{thanks to}. \cle{Connecteur}
\end{itemize}
\end{corrige}

\begin{corrige}[Exercice 11 — Composition : sujet type Baccalauréat]
\textbf{Proposition de rédaction modèle (environ 95 mots) :}

\eng{Last month, the inhabitants of my quarter in Yaoundé organized a keep-clean day. The quarter head had mobilized everybody a week before, and many volunteers took part in the activity.}

\eng{Early in the morning, the men cleaned the gutters and repaired the road, while the women swept the market. The children picked up the rubbish, and even the old people lent a hand. Moreover, some traders raised funds to buy new dustbins. All this work was voluntary, but everybody was happy to contribute.}

\eng{Communal work is very important because it keeps our environment clean and teaches young people solidarity and cooperation. That is why our quarter develops faster when we work together.}

\textbf{Barème indicatif (10 pts) :} respect du sujet et de la longueur 2 pts ; richesse et exactitude du lexique du thème (au moins quatre expressions) 3 pts ; correction grammaticale (prétérit, accords, prépositions) 3 pts ; organisation en trois paragraphes et connecteurs 2 pts.

\textbf{Points de vigilance :}
\begin{itemize}
\item employer le \cle{prétérit} pour la description (\eng{organized, mobilized, cleaned, raised}) ;
\item \eng{to take part in}, \eng{to contribute to}, \eng{to be proud of} : vérifier les \cle{prépositions} ;
\item ne pas écrire « assist the meeting » pour « assister à la réunion » mais \eng{attend the meeting} (\cle{faux ami}) ;
\item \eng{rubbish, work, money} sont \cle{indénombrables} : pas de \emph{-s} ;
\item varier les \cle{connecteurs} : \eng{first, then, moreover, finally, that is why}.
\end{itemize}
\end{corrige}

\newpage

\coursec{Fiche bilan}

\begin{popfigure}
\begin{tikzpicture}[mindmap, grow cyclic, every node/.style={concept, text=white, font=\small\bfseries, align=center},
root concept/.append style={concept color=popblue, font=\bfseries},
level 1/.append style={level distance=3.4cm, sibling angle=51.5},
level 2/.append style={level distance=1.9cm, sibling angle=45, font=\scriptsize}]
\node[root concept]{Communal work}
child[concept color=poporange]{ node {Activities}
child{ node {clean-up} }
child{ node {plant trees} }
}
child[concept color=popgreen]{ node {Taking part}
child{ node {volunteer} }
child{ node {join in} }
}
child[concept color=popgold]{ node {Helping}
child{ node {lend a hand} }
child{ node {support} }
}
child[concept color=poppurple]{ node {Organising}
child{ node {plan} }
child{ node {assign tasks} }
}
child[concept color=poppink]{ node {Collocations}
child{ node {community spirit} }
child{ node {team up} }
}
child[concept color=popdark]{ node {Faux amis}
child{ node {assist $\neq$ assister} }
child{ node {actually $\neq$ actuellement} }
}
child[concept color=popblue!70]{ node {Word families}
child{ node {help, helpful, helpless} }
};
\end{tikzpicture}
\legende{Carte mentale : le vocabulaire du travail communautaire}
\end{popfigure}

\begin{bilan}
\textbf{1. Lexique clé à connaître par cœur}
\begin{center}
\begin{tabularx}{\linewidth}{|X|X|}\hline
\rowcolor{popblueL} \textbf{English} & \textbf{Français} \\ \hline
communal work / community work & le travail communautaire \\ \hline
a volunteer (n.) / to volunteer (v.) & un bénévole / se porter volontaire \\ \hline
to take part in / to participate in & participer à \\ \hline
to lend a hand / to give a hand & donner un coup de main \\ \hline
a clean-up campaign & une campagne de nettoyage \\ \hline
to plant trees & planter des arbres \\ \hline
to repair a bridge / a road & réparer un pont / une route \\ \hline
to raise funds & collecter des fonds \\ \hline
community spirit & l'esprit communautaire \\ \hline
to team up with & faire équipe avec \\ \hline
to assign tasks & attribuer les tâches \\ \hline
a neighbourhood & un quartier \\ \hline
\end{tabularx}
\end{center}

\textbf{2. Collocations et expressions idiomatiques}
\begin{itemize}
\item \eng{many hands make light work} — « À plusieurs, le travail est plus léger » (l'union fait la force).
\item \eng{to pull together} — « se serrer les coudes ».
\item \eng{to do one's bit / one's share} — « faire sa part ».
\item \eng{to roll up one's sleeves} — « retrousser ses manches », se mettre au travail.
\end{itemize}

\textbf{3. Règles à retenir}
\begin{center}
\begin{tabularx}{\linewidth}{|X|X|}\hline
\rowcolor{popblueL} \textbf{Règle} & \textbf{Exemple} \\ \hline
\eng{to take part in} + nom/\textsc{v-ing} (jamais sans \eng{in}) & \eng{She took part in cleaning the market.} \\ \hline
\eng{to look forward to} + \textsc{v-ing} & \eng{We look forward to helping.} \\ \hline
\eng{help} + (to) + base verbale & \eng{They helped (to) build the school.} \\ \hline
\end{tabularx}
\end{center}

\textbf{4. Faux amis}
\begin{itemize}
\item \eng{to assist} = aider ; « assister à » = \eng{to attend}.
\item \eng{actually} = en fait ; « actuellement » = \eng{currently}.
\item \eng{a community} = une communauté ; attention à l'orthographe (un seul \emph{m} redoublé : \emph{comm-unity}).
\end{itemize}

\textbf{5. Familles de mots} : \eng{help} $\rightarrow$ \eng{helpful / helpless / helper} ; \eng{volunteer} $\rightarrow$ \eng{voluntary / voluntarily} ; \eng{organise} $\rightarrow$ \eng{organisation / organiser / organised}.

\textbf{6. Méthode express} : pour décrire une action communautaire : (1) situer (où, quand, qui), (2) énumérer les tâches avec des verbes d'action, (3) conclure sur le bénéfice pour la communauté.
\end{bilan}

\begin{aretenir}[Checklist avant le Bac]
\begin{itemize}
\item[$\square$] Je connais au moins 15 mots du champ lexical du travail communautaire.
\item[$\square$] Je sais employer \eng{to take part in} avec \eng{in} suivi d'un nom ou d'un \textsc{v-ing}.
\item[$\square$] Je distingue \eng{volunteer} (nom et verbe) et je le prononce « vo-lon-tire ».
\item[$\square$] Je ne confonds pas \eng{assist} et « assister à » (\eng{attend}).
\item[$\square$] Je ne traduis pas « actuellement » par \eng{actually} mais par \eng{currently}.
\item[$\square$] Je sais utiliser \eng{to lend a hand} et \eng{to team up with} dans une phrase.
\item[$\square$] Je connais les familles de mots : \eng{help/helpful}, \eng{organise/organisation}.
\item[$\square$] Je sais que \eng{look forward to} est suivi de \textsc{v-ing}.
\item[$\square$] Je sais décrire une action communautaire en trois étapes (lieu, tâches, bénéfice).
\item[$\square$] Je peux citer deux expressions idiomatiques sur l'entraide (\eng{many hands make light work}, \eng{pull together}).
\end{itemize}
\end{aretenir}

\findecours

\end{document}
